% !TeX program = lualatex
\documentclass[11pt,a4paper]{ltjsarticle}
\usepackage[margin=25mm,headheight=17pt]{geometry}
\usepackage[haranoaji]{luatexja-preset}
\usepackage{amsmath,amssymb,amsthm,mathtools,mathrsfs,booktabs,array,enumitem,longtable}
\usepackage[unicode,hidelinks]{hyperref}
\usepackage{fancyhdr}
\pagestyle{fancy}\fancyhf{}\fancyfoot[C]{\thepage}
\fancyhead[L]{\small\nouppercase{\leftmark}}
\renewcommand{\sectionmark}[1]{\markboth{#1}{}}
\setlength{\emergencystretch}{3em}\allowdisplaybreaks[2]
\setlength{\parskip}{.15em}
\newtheoremstyle{jp}{.6em}{.6em}{\normalfont}{}{\bfseries}{.}{.5em}{}
\theoremstyle{jp}
\newtheorem{theorem}{定理}[section]
\newtheorem{proposition}[theorem]{命題}
\newtheorem{lemma}[theorem]{補題}
\newtheorem{corollary}[theorem]{系}
\newtheorem{definition}[theorem]{定義}
\newtheorem{remark}[theorem]{注意}
\newtheorem{example}[theorem]{例}
\newtheorem{question}[theorem]{問題}
\renewcommand{\proofname}{証明}
\newcommand{\kk}{\mathbb C}\newcommand{\PP}{\mathbb P}
\newcommand{\OO}{\mathcal O}\newcommand{\A}{\mathcal A}
\newcommand{\B}{B_{\mathrm{VK}}}\newcommand{\Blo}{B_{p,q}}
\newcommand{\HH}{\operatorname{HH}}\newcommand{\Hom}{\operatorname{Hom}}
\newcommand{\End}{\operatorname{End}}\newcommand{\Sym}{\operatorname{Sym}}
\newcommand{\Perf}{\operatorname{Perf}}\newcommand{\Der}{\operatorname{Der}}
\newcommand{\Inn}{\operatorname{Inn}}\newcommand{\gr}{\operatorname{gr}}
\newcommand{\rk}{\operatorname{rank}}\newcommand{\Spec}{\operatorname{Spec}}
\newcommand{\id}{\operatorname{id}}\newcommand{\Ext}{\operatorname{Ext}}
\newcommand{\Tor}{\operatorname{Tor}}\newcommand{\Ann}{\operatorname{Ann}}
\newcommand{\coker}{\operatorname{coker}}\newcommand{\Supp}{\operatorname{Supp}}
\newcommand{\pd}{\operatorname{pd}}\newcommand{\depth}{\operatorname{depth}}
\newcommand{\Frac}{\operatorname{Frac}}\newcommand{\Aut}{\operatorname{Aut}}
\newcommand{\Out}{\operatorname{Out}}\newcommand{\GL}{\operatorname{GL}}
\newcommand{\PGL}{\operatorname{PGL}}\newcommand{\Hilb}{\operatorname{Hilb}}
\newcommand{\dashsim}{\mathrel{\overset{\mathrm{bir}}{\sim}}}
\setcounter{tocdepth}{1}

\hypersetup{pdftitle={Veronese--Kronecker研究ノート v40}}
\title{Veronese--Kronecker 研究ノート\\\large 確定した構造、未決定の障害、次の判断基準}
\author{}\date{2026年10月30日\quad 第40版}
\begin{document}\maketitle
\begin{abstract}
再利用する証明機構、確定した構造、未決問題の正確な条件を記す。
全重みのCGG四項成分について、純粋三変数部分の内在的持上げと通常hullの分解を得る。
同じ重みの全共鳴は同型な通常hullを持つが、対称性保持部分は一致するとは限らない。
代数と導分を同時に変形する障害複体から、六項側の結合性障害を四項側の対称性障害へ比較する。
端点の一の冪根量子化では、$H=f(Q,u)+xG$ の全中心と、$f=0$ の全ファイバーを決定する。
同じ関数斜体を持ちながら全中心の特異次元が異なる二例を与える。
最深点の完成後にもBrauer周期・指数が最大値を保つことを、新しい付値とノルム条件で証明する。
一般境界では、単純加群の巡回被覆が可換解消の例外因子上のBrauer留数も支配する。
従来の対称性保持部分の二重構造、固定量子hull、Poisson葉空間、Prym配置、K3境界period、
七次被覆の成果は、対象と同値関係を区別して保持する。
\end{abstract}
\tableofcontents

\section{読者が最初に固定すべき区別}

本研究で使う四つの代数は、互いに異なる対象である。
\begin{center}\small
\begin{tabular}{p{24mm}p{46mm}p{74mm}}\toprule
記号&対象&変形・同値の意味\\\midrule
$B_{\rm VK}$&次数 $(0,1,3,4)$ の有限切り出し&四頂点を標識した結合代数の変形。標準部分代数の一意性を適用する。\\
$A_E$&第二関係成分 $Y_C$ 上の有限代数&三つの矢空間の基底変更を商する。埋め込まれた二十一点配置から復元する。\\
$T/B$&残留十一点の可換有限被覆&$B=\kk[[t_1,\ldots,t_{16}]]$。相対 potential はこの被覆の勾配表示の問題。\\
$\mathcal A$&明示残留点の完備非可換局所代数&固定した二生成元代数。全中心、PI 性、連続 CY 性、非可換 potential を判定する。\\\bottomrule
\end{tabular}
\end{center}

さらに公表原稿IVの $A_F$ は、$\mathcal A$ と同じ三次接錐を持つ\emph{比較族}である。
$A_F$ の全分類を、そのまま固定代数 $\mathcal A$ の分類へ移してはならない。
また、有限代数の同型、例外列を保持する圏の同値、例外列を忘れた導来同値は別の同値関係である。

\section{標準六項への補完を次に使うための要点}

今回の主対象は $\Lambda=R_{(0,1,2,3,4,5)}$ と $B=e\Lambda e$ の組である。
可換ファイバーの標準埋め込みを固定する。
次の三つを別々の問題として扱う。
\begin{enumerate}
\item 一次方向を周囲へ持ち上げること。これは $\rho^2$ の同型により常に可能。
\item 一次方向を固定し、補完可能な全形式族を何か一つ作ること。必要十分条件はPoisson性。
\item 高次の乗法まで指定済みの全形式 $B$ 変形を補完すること。一次Poisson性だけでは足りない。
\end{enumerate}
第二区分の十分性は発散零代表と量子化による。必要性は二次のHochschild障害による。
第三区分については、下の二変数族が全Artin基底での正確な試金石になる。

\subsection{二次的基底の認定方法}
六項のminiversal基底と四項の基底は、compatibleな選択のもとで
\[
 R_\Lambda\simeq k[[u_1,\ldots,u_{53}]]/(q_1,\ldots,q_{38}),\qquad
 R_B\simeq k[[u_1,\ldots,u_{53}]]
\]
となり、制限は基底上の閉埋め込みになる。
三十八式はJacobi式の独立な同次二次係数である。
基底の座標は標準的でなく、hullは自己同型を持つ変形関手のfine表現ではない。

この結果の認定補題は他の問題にも再利用できる。
同次二次生成イデアル $J\subset P=k[[u]]$ の零点全体上に接写像が同型の全形式族があれば、
miniversal基底 $P/I$ からの写像を取り、形式座標変更により $I\subset J$ とできる。
さらに二次障害で $I_2=J_2$ が分かれば
$J\cap\mathfrak m^3=\mathfrak mJ$ とNakayamaから $I=J$ となる。
高次の逐次消去を続ける代わりに、二次軌跡全体の全形式実現を探すことが有効である。

\subsection{全Artin基底での二変数モデル}
$[x_1,x_2]=2s x_1x_3$、$x_3$ 中心のOre部分と、
$x_iw_1=(1+t\delta_{i2})v_i$、$w_0x_i=(1-t\delta_{i2})v_i$ という平方零部分から
$B_{s,t}$ を構成する。補完存在条件は正確に $st=0$ である。
これは $s,t$ が独立な形式変数の場合だけでなく、任意の局所Artin環における値にも成立する。

必要性の証明で再利用すべき道具は二つある。
第一に、$HH^1$ 制限の全射性と $HH^2$ 制限の単射性から、corner同型を周囲の同型へ持ち上げられる。
第二に、任意の $st\neq0$ に対して
\[
 J=T(st),\qquad T'=T/\mathfrak m_TJ,
 \qquad I=J/\mathfrak m_TJ\neq0
\]
を取ると、$\mathfrak m_{T'}I=0$ の小拡大に還元できる。
そこで純粋な $s$ 項・$t$ 項は軸上の結合律で消え、三次以上の混合項も消えるため、
唯一残る障害は $(st)[\alpha_s,\alpha_t]$ である。
Ore表示の失敗だけでなく、全ての周囲の変形に対する不可能性を証明するにはこの還元が必要である。

$B_{\hbar,\hbar^r}$ ($r\ge2$) は非零Poisson一次方向を持ち、
$\hbar^{r+1}$ を法として補完可能、$\hbar^{r+2}$ を法として補完不能である。
従って、特定次数まで補完できたという計算結果だけから全次数の補完を推測しない。
ただしこれは有限本の方程式の不存在ではない。ここでは全条件が一式 $st=0$ である。

\subsection{今回の原稿構成の判断}
公表稿VIIIは、三次制限・二次的hull・非純次元性・節点モデルを一つの補完問題としてまとめる。
公表稿Iの標準部分代数の一意性とLO比較とは、主定理も必要な背景も異なるため独立稿とする。
既存資料にある四項側の無障害性は、VIIIで短い自己完結した証明として再構成した。
I--VIIの証明範囲は今回変更していない。
量子中間点の随伴次数同定を元のLO代数の実現へ移さないこと、
標準六項の補完不能性を全ての導来埋め込みの不存在へ移さないことを維持する。

\section{三次制限写像の確定事項と継続時の注意}
\label{sec:v31c-notes}

今回の追加計算で、標準四対象部分圏への制限は
\[
 \rho^1:\mathbf k^{19}\xrightarrow{\sim}\mathbf k^{19},\qquad
 \rho^2:\mathbf k^{53}\xrightarrow{\sim}\mathbf k^{53},\qquad
 \operatorname{rank}\bigl(\rho^3:\mathbf k^{41}\to\mathbf k^{39}\bigr)=1
\]
と確定した。恒等慣性成分だけでも既に像の直線全体を生成する。

\subsection{再利用する表現と重み}
$V$ を座標関数 $x_i$ の表現とする。
$w$ の拡大に関する $B$ の重み分解は
\[
\begin{array}{c|rrr}
w\text{-重み}&\dim\operatorname{HH}^1&\dim\operatorname{HH}^2&
 \dim\operatorname{HH}^3\\\hline
-1&10&35&0\\
0&9&18&28\\
1&0&0&11
\end{array}
\]
である。正の三次表現は
\[
 \operatorname{HH}^3(B)_{+1}
 =\det(V)^{-1}\oplus\operatorname{Sym}^3V\otimes\det(V)^{-2}.
\]
外側の正の三次表現は $\bigl(\det(V)^{-1}\bigr)^{\oplus3}$ であるため、
同変性だけで制限の階数は高々一となる。

全核は四十次元だが、primary obstruction の線形包は三十八次元である。
恒等慣性成分を $R_6$ と同一視すると後者は
\[
 (x_1,x_2,x_3)R\cap R_6
 =\operatorname{Sym}^6V\oplus w\operatorname{Sym}^3V
\]
であり、stacky point で消える切断である。
三十八個の二次式の線形独立は確認したが、正則列性や locus の余次元を
意味するものではない。

\subsection{保持する短い証明装置}
Euler 場との wedge の primitive は
\[
 f(a:i\to j,b:j\to k)=(h_i+h_j+h_k)A(a,b),\qquad
 d^2f=(E\wedge A)|_B
\]
である。この式は全ての重み零の多項式 biderivation に適用できる。
また、十八項の汎関数
\[
 \mathcal L(\varphi)=\sum_{\sigma\in S_3}\operatorname{sgn}(\sigma)
 \sum_{\ell=1}^3[x_\ell w]
 \varphi(x_{\sigma(1)},x_{\sigma(2)}x_\ell,x_{\sigma(3)})
\]
は $\mathcal L(d^2f)=0$、$\mathcal L(w\partial_{123})=24$ を満たす。
三次制限の非零性はこの式だけで証明できる。

\subsection{計算上の再現と範囲}
\texttt{verify\_ambient\_cohomology\_v31.py} と
\texttt{verify\_full\_rho3\_v31.py} を、その隣の JSON certificate とともに保存する。
Python 標準ライブラリだけで実行でき、通常実行は読み取り専用である。
前者は四対象の全複体と恒等慣性成分の制限を、後者は六対象の正の重み全体を
検証する。浮動小数点の結果や推定階数には依存しない。

後者の階数計算は次の線形代数を用いる。
$Q$ を $\operatorname{HH}^3(B)_{+1}$ の双対基底を restriction で引き戻した
十一行の行列とすると、
\[
 \operatorname{rank}(\rho^3_{+1})=
 \operatorname{rank}\begin{bmatrix}d^3_\Lambda\\Q\end{bmatrix}
 -\operatorname{rank}(d^3_\Lambda)=1.
\]
この計算法は、source の三次 cocycle 基底を大量に書き出す必要がない。
restriction が chain map であることと $Qd^2_\Lambda=0$ も検証する。

二つの非自明慣性成分の個々の HKR 基底が $B$ でどの scalar 倍の
$[w\partial_{123}]$ に写るかは同定していない。
全制限の階数・核次元・primary obstruction の計算には不要である。
HKR の splitting と restriction の両立をさらに調べるときに限り、
この同定を継続課題とする。

\subsection{既存資料の監査結果}
元の追加検証器の全有限検査は、そのままのファイルを読み取り専用で実行して
成功した。$\pi=x_1x_3\partial_{12}+x_2w\partial_{2w}$ に対する
Ore--平方零族の一次積と $\pi|_B$ の差は
$d(x_1x_3\partial_{x_1}\partial_{x_2})$ である。
Ore 部分の PBW 性と、$w$ を含む非恒等射三本の可合成列が存在しないことから、
全次数平坦性の構成にも欠落は見つからなかった。

v22 の量子多項式環の同定が直接与えるのは
$\operatorname{gr}_J B_{p,q}\simeq(Q_{p,q})_{(0,1,3,4)}$ である。
量子中間代数が非可換重み付き射影空間の半直交成分であることから、
元の LO 代数自体について同じ埋め込みが存在すると推論してはならない。

\section{全重みで増大する厚みと明示的な量子代数}
\label{sec:allm-notes-v33}

この節の結果は公表稿IX\cite{PaperIX}に証明をまとめる。
一般重みの厚みは仮説ではなく、全ての整数 $m\ge3$ に対する定理となった。
今後は、この機構と明示的な代数を使って別の変形問題を判断する。

\subsection{全方程式と、重み三だけに残る慣性成分}
\[
 X_m=\mathbb P_{\mathrm{stack}}(1,1,1,m),\quad
 S_d=\mathbb C[x,y,z]_d,\quad N_m=\binom{m+2}{2}
\]
と置く。$m\ge3$ では標準形の座標は
\[
 c\in S_3,\qquad D\in\mathfrak{sl}_3,\qquad
 [F]\in S_{m+1}^{\oplus3}/\mathbf xS_m
\]
であり、全Jacobi式は
\[
 D(c)=0,\qquad
 \frac{m+3}{3}F(c)-\frac m3(\mathbf x\times D)\cdot F-c\operatorname{div}F=0.
\]
変数数と独立な二次式の数は
\[
 d_m=m^2+6m+26,\qquad q_m=10+\binom{m+5}{2}.
\]
二次式の個数を余次元と取り違えない。
$m=1,2$ では二次場の係数に追加の $w$ 項があるので、上の標準形をそのまま使わない。

標準 $(m+3)$ 項の代数 $\Lambda_m$ は
$\dim\Lambda_m=\binom{m+6}{4}+15$ であり、そのminiversal基底は
このJacobiスキームの原点完備化である。
証明に再利用する機構は、Euler発散による一意な持ち上げ
\[
 L(U)=U-\frac1{m+3}E\wedge\delta U
\]
と、全Jacobi軌跡上の量子化、二次生成性に対するNakayamaの補題である。
全高次障害を個別に計算する必要はない。

$HH^2(\Lambda_m)$ は $d_m$ 次元である。
$m\ge4$ では $HH^3(\Lambda_m)$ は $q_m$ 次元で、一次方向間の括弧の像が全体を張る。
一方 $m=3$ では $\dim HH^3=41$、括弧の像は三十八次元である。
この差には、重み六の $w^2$ と、二つの非自明慣性成分が寄与する。
一般重みで慣性成分を数える際は、各点 $B\mu_m$ 上の法束行列式の文字
$\chi^3$ に対して\emph{全安定化群}の不変部分を取る。
個々の慣性元の位数だけを見ると、$m$ が三の倍数の場合に誤った寄与を加えてしまう。

\subsection{一般点のArtin環を決める三つの情報}
基準点は全重みを通じて
\[
 c_0=xy(x+y),\qquad D_0=(2x+5y)\partial_z,\qquad F=0
\]
に取れる。この点を含む適切な開集合で、全完備局所環は
\[
 \boxed{\mathbb C[[s_1,\ldots,s_{N_m+8},a,b,u]]/(ab,b^2,a^{m+2}u)}.
\]
その証明を再利用するには、次の三点を保つ。
第一に、基底 $D(c)=0$ の切片は $m$ によらず $(ab,b^2)$ であり、
平方零方向 $b$ を消去せずに混合式へ渡す。
第二に、混合核 $N_m+1$ 方向のうち、次数 $m$ の $w$ 移動の
$N_m$ 方向は非被約基底上でも実際の核の自由な直和因子である。
残る一方向は $F=(0,0,x^{m+1})$ で、$z$ 重量は $-1$。
余核の四方向の重量は $m+1,m+2,m+2,m+3$ である。
従って $a$ の障害は $m+2$ 次以上、$b$ の障害は零となる。

第三に、$b=0$ の三角形 $c=xy(x+y+az)$ の三交点での評価を使う。
Schur消去後の出力を
\[
 r_0x^2z^{m+1}+r_1xz^{m+2}+r_2yz^{m+2}+r_3z^{m+3}
\]
とすると、交点評価から $r_1=ar_0,r_2=r_3=0$ となる。
消去行列の重量から $r_0$ の $a$ 次数は $m+2$ 以下であり、
$a\ne0$ 側の階数が正確に一つ増えるので、$a^{m+2}$ 係数は必ず非零である。
ここが有限個の重みの検算を全重みの証明へ置き換えた箇所である。

引戻し成分の一般点では $u$ が単元なので
\[
 K[a,b]/(a^{m+2},ab,b^2),\qquad
 1,a,\ldots,a^{m+1},b
\]
が得られる。成分次元は $N_m+9$、長さは $m+3$、socle次元は二である。
被約化は滑らかであっても、この厚みは一般点に残る。

\begin{center}\small
\begin{tabular}{rrrrr}\toprule
$m$&$d_m$&$q_m$&引戻し成分次元&一般点の長さ\\\midrule
3&53&38&19&6\\
4&66&46&24&7\\
5&81&55&30&8\\
6&98&65&37&9\\\bottomrule
\end{tabular}
\end{center}

primary分解と交叉の厚みも全重みで一様である。
\[
 (ab,b^2,a^{m+2}u)=(a^{m+2},ab,b^2)\cap(b,u).
\]
二つの被約枝は引戻し型と $L(1,1,1,m)$ 型であり、被約交叉は滑らか。
primary成分どうしの交叉は長さ $m+2$、一般超平面切断は長さ $m+4$ である。
局所環はCohen--Macaulayだが完全交叉ではない。

Poissonスキームは斉次錐なので、原点での形式化後にも対応する成分と一般点の長さが残る。
証明では、次数付き整域の完備化も最小斉次項の積により整域となること、
原点の局所環からの忠実平坦性、組成列の基底変換を用いる。
従って $\Lambda_m$ のminiversal基底にも、一般点の長さが任意に大きい成分が現れる。
この引戻し成分の原点基底の計算だけでは、固定した非零量子代数の変形空間は決まらない。
接線共鳴の固定量子化については、公表稿Xで別途形式弧移送により決定した。

\subsection{全重みで使える明示的AS正則4-Calabi--Yau代数}
二変数の $c_3,L_1,f_{m+1}$ を任意に取り、
\[
 \begin{split}
 \delta_m(x)&=c_y/3+mxL/(3(m+3)),\\
 \delta_m(y)&=-c_x/3+myL/(3(m+3)),\\
 \delta_m(w)&=f-mLw/(m+3)
 \end{split}
\]
と置く。$\operatorname{div}\delta_m=0$ であり、
\[
 S_m=\mathbb C[x,y,w][z;\hbar\delta_m]
\]
は引戻し族の適合座標表示を量子化する。
PBW基底は $x^iy^jw^kz^\ell$、Hilbert級数は $((1-t)^3(1-t^m))^{-1}$。
noetherian整域で、AS正則次元四、通常の4-Calabi--Yau、Gorenstein parameter $m+3$ である。
星積は
\[
 a\star b=\sum_{r\ge0}\frac{\hbar^r}{r!}
                 (\partial_z^r a)(\delta_m^r b)
\]
であり、各多項式対について有限和になる。
抽象的な形式存在を超えて、任意の複素数 $\hbar$ で座標環を持つ点が有用である。

標準列 $(\mathcal O,\ldots,\mathcal O(m+2))$ はfull strongであり、
$B_m=(S_m)_{(0,1,m,m+1)}$ は階数 $2m^2+6m+20$ の四対象代数を与える。
$\operatorname{Perf}(B_m)$ は $D^b(\operatorname{qgr}S_m)$ の許容部分圏である。
$m=3$ は元のVK問題に対する、新しい具体的な標準六項実現である。
$m>3$ で $B_m$ 自身の変形が無障害であることや、外側との接空間同型は未確認である。
この二つは今回の量子化の存在からは従わない。

\subsection{接線共鳴は全重みの成分系列を作る}

$q=xz-y^2$、$N_0=x\partial_y+2y\partial_z$ とする。
$0\le j\le\lfloor(m-1)/2\rfloor$、$r=m-2j>0$ に対して
\[
 c=xq,\qquad D=a_jN_0,\quad a_j=-\frac{4m+3-6j}{2m},\qquad
 F_j=q^j(0,yz^r/2,z^{r+1})
\]
を基準点とする。$(c,D)$ の七次元軌道上の混合核束の閉包は、
次元 $N_m+8$ の既約成分 $\mathscr E_{m,j}$ である。
全 $m,j$ について一般点の局所環と $F=0$ の全形式近傍が
\[
 K[\gamma]/(\gamma^{j+1}),\qquad
 \mathbb C[[s_1,\ldots,s_{N_m+7},\kappa,\gamma,u]]/
                    (\kappa u,\gamma^{j+1}u)
\]
と決まった。固定した $m$ で相異なる $\lceil m/2\rceil$ 成分を与え、
長さは $1,\ldots,\lceil m/2\rceil$ を尽くす。
$j=0$ は滑らか、$j\ge1$ は一般点でも非被約である。
$u=0$ の枝は $N_m+9$ 次元の $L(1,2,m)$ 成分となる。

再利用すべき証明の機構は、次の具体的な四点である。
基底 $D(c)=0$ は七軌道方向を除くと
\[
 c=x(xz-y^2+\gamma z^2),\qquad
 D=a((x+2\gamma z)\partial_y+2y\partial_z)
\]
の二変数となる。混合核を $F_x=0$ のEuler代表で書くと
\[
 (F_y,F_z)=(xA+yB,2yA+2zB),\qquad
 N_0(A)=(T-2j)B,\quad T=2x\partial_x+y\partial_y
\]
である。実際の $w$ 移動は
\[
 A=-\tfrac12(T+1-2j)g,\qquad B=-\tfrac12N_0(g)
\]
となり、階数 $N_m$ の自由な直和因子をなす。
例外的な $T$ 重量 $2j-1\to2j$ の二重対角行列は単射で、
余核の一元が $q^jz^r$ である。

混合余核の $H_0=x\partial_x-z\partial_z$ 重量は
$-m-3,-m-2,-m-1,1-r$ である。
重量を保つ単元ブロックの逆行列は $\gamma$ の有限多項式なので、
残余列の形は正確に
\[
 f(a),\qquad \gamma^{j+1}g(a),\qquad0,\qquad\gamma^{j+2}h(a)
\]
へ制約される。$f$ は $a-a_j$ の単元倍である。
二つの特異点 $[0,\pm\sqrt\gamma,1]$ での評価から
$h(a_j)=-g(a_j)$、横断二次曲線の場合の混合核を次数降下で全て解くことから
$g(a_j)\ne0$ と分かる。この二つを合わせて全理想が
$(a-a_j,\gamma^{j+1})$ と確定する。
この方法は最初の障害次数だけでなく、高次の隠れた条件がないことまで保証する。

primary分解は
\[
 (\kappa u,\gamma^{j+1}u)=(u)\cap(\kappa,\gamma^{j+1})
\]
である。被約交叉は滑らかな $N_m+7$ 次元、primary成分の交叉には
長さ $j+1$ の厚みが残る。
$r=0$ は $F_x$ の消去が成立しないため、この成分定理の対象から外す。
量子化と中心の定理は、この端点も含めて別に成立する。

\subsection{重み三の第二共鳴と、二次で止まる方向}

$m=3,j=1$ の同じ成分は、$qz$ 型と等変同値な $y^2z$ 型で
\[
 D=-\tfrac32N_0,\qquad F=(0,y^3z/2,y^2z^2)
\]
とも書ける。この点の完備局所環は
$\mathbb C[[s_1,\ldots,s_{18},b]]/(b^2)$ である。
直接の有理係数消去でも、この全局所環を独立に確認した。
特に
\[
 \begin{split}
 c_b&=x(xz-y^2-bz^2/3),\\
 D_b&=(0,-3x/2+bz,-3y),\\
 F_b&=(0,y^3z/2-byz^3/12,y^2z^2+bz^4/6)
 \end{split}
\]
について $D_b(c_b)=0$、$M(c_b,D_b;F_b)=-7b^2xz^5/36$ である。
接触を解く一次方向が二次で止まることを、今後の障害計算の最小例として使える。
この成分を加え、重み三の認定済み大域成分は少なくとも十個となった。

非零の接線パラメータの四共鳴は、滑らかな十七変数を除いて
\[
 \begin{array}{c|c}
 a=2&(kv_1,\ldots,kv_{10})\\
 a=-1&(kv_1,\ldots,kv_4)\\
 a=-5/2&(kv,bv)\\
 a=-3/2&(kv,b^2v)
 \end{array}
\]
という局所方程式を持つ。最初の二つは $R(1,5),R(2,4)$ との交叉である。
非共鳴で混合核の次元が十であっても、常に全てを $w$ 移動と同一視してはいけない。
$a=1$ では $p=(q,0,0)$ となり、移動核には $x^3$ が入る。
欠ける方向 $x^3N_0$ は
$p_a\times\nabla(x^3)=(1-a)x^3N_0$ の有限極限として得られ、
全空間はこの点でも滑らかな十九次元である。

\subsection{非線形CGG族のNakayama自己同型は全て決まる}

$m=2j+r\ge3$、$j,r\ge0$ とし、$P=q^jz^r$ または $y^{2j}z^r$ とする。
\[
 N=N_0-P\partial_w,\qquad d=m+r+1,
\]
\[
 X_\alpha=\alpha x\partial_x+(\alpha+1)y\partial_y+(\alpha+2)z\partial_z
                   +(m\alpha+d)w\partial_w
\]
に対する有限CGG積を考える。
$T=1+sN$ は星積の自己同型で、変数を $x,y,z,w$ の順に加えると反復Ore表示を得る。
全 $\alpha,s$ でnoetherian整域かつAS正則四次元であり、
Hilbert級数は $((1-t)^3(1-t^m))^{-1}$、Gorenstein parameterは $m+3$ となる。
そのNakayama自己同型は全次数で
\[
 \boxed{\mu=T^\nu,\qquad \nu=(m+3)\alpha+m+r+3=\operatorname{div}X_\alpha-1}
\]
と決まる。従って $s\ne0$ で通常の4-Calabi--Yauとなるのは
$\alpha=-(m+r+3)/(m+3)$ の場合に限る。

再利用できる剛性は、次数一で生成される三変数部分代数を各元ごとに固定する
次数付き自己同型が恒等になることである。
$\psi(w)=uw+b$ を、$N$ 不変な奇数次数の $J$ に作用させる。
$x,q$ との可換性から $N(\psi J-uJ)=0$、$y$ との交換関係から
その $H_0$ 重量が零となる。
奇数次数・零重量の空間上の $N_0$ は二重対角の正則行列であり、従って差は零である。
$J\bmod x$ が非零であることから $u=1,b=0$ となる。
この議論は $\alpha=0,-1$ も含み、Ore定理だけでは未決の最後の平行移動項を消す。

Lecoutre--Sierraの三変数Nakayama公式を用いる部分と、この四変数剛性を区別する。
通常のCY性を発散零の一次テンソルだけから推論したわけではない。
$s\ne0$ の次数付き自己同型群は
$\mathbb C^\times\times(\mathbb C,+)$ であり、作用は
重み付きスカラー変換と $T^t$ である。
全ての族は標準 $(m+3)$ 項のfull strong列と四対象の許容成分を持つ。

\subsection{多項式中心の三分岐とDu Val曲面}

$q$ 型で
\[
 c_r=(-1)^r\frac{(2r)!!}{(2r+1)!!},\qquad
 J=x^{r+1}w+q^j\sum_{k=0}^r\binom rk\frac{q^{r-k}y^{2k+1}}{2k+1}
\]
とすると、
\[
 \Delta=(J^2+c_r^2q^{m+r+1})/x^{r+1}
\]
は多項式で、
\[
 \ker N=\mathbb C[x,q,J,\Delta]
 \simeq\mathbb C[X,Q,U,V]/(X^{r+1}V-U^2-c_r^2Q^{m+r+1})
\]
である。$x$ を反転した核だけでは生成性は決まらない。
$x=0$ でのcusp表示と $\Delta$ の非零の $w$ 一次係数から
$B/xB\to A/xA$ の単射性を示し、$x$ に関する飽和で核全体を得るのが重要である。

Calabi--Yauのパラメータでは、Poisson中心と量子中心が同じ部分環として
$\ker N\cap\ker X_\alpha$ に一致する。
$2(m+3)H=- (m+3)X_\alpha$ の重量は
\[
 x:m+r+3,\quad q:2r,\quad J:r(m+r+1),\quad
 \Delta:(r-1)m+(r-3)(r+1)
\]
なので、$r\ge2$ では中心は定数環、$r=0$ では $\mathbb C[q,J]$ である。
$r=1$ では
\[
 S=q^2\Delta,\quad P=2xJ\Delta^{(m+3)/2},\quad
 R=q(x^2\Delta+J^2)\Delta^{(m+3)/2}
\]
が中心全体を生成し、その唯一の関係は
\[
 \boxed{R^2=SP^2+\frac{16}{81}S^{m+3}}
\]
である。これは $D_{m+4}$ 型のDu Val曲面である。
生成性は、$\mathbb C[x,q,\Delta]\oplus J\mathbb C[x,q,\Delta]$ の
零重量単項式を指数の偶奇で完全に分けることで証明する。

$y$ 型との等変変換も多項式で与えられる。
\[
 \lambda_{j,r}=(-1)^j\frac{(2j-1)!!(2r+1)!!}{(2j+2r+1)!!},\qquad
 w'=\lambda_{j,r}^{-1}(w+b_{j,r}),
\]
\[
 b_{j,r}=(J_0^{(y)}-\lambda_{j,r}J_0^{(q)})/x^{r+1}\in S_m.
\]
多項式性は二つのTaylor表示の $r$ 次までの係数一致から従う。
特に重み三の第二共鳴は $w'=-5w-yz^2$ で $qz$ 型となり、その量子中心は $D_7$ 型である。

一般点の長さと中心の型が同時に読める系列は
\[
 m=2\ell-1,\quad j=\ell-1,\quad r=1:\qquad
 \text{成分の長さ }\ell,\quad \text{量子中心 }D_{2\ell+3}\quad(\ell\ge2)
\]
である。前者は可換スタック側のPoissonスキームの一般点長さである。
公表稿Xでは形式弧移送と軌道方向の除去により、固定した非零量子化の
標準有限代数のminiversal環も正確に $\mathbb C[[\gamma]]/(\gamma^\ell)$ と同定した。

\subsection{共通の関数斜体と、同じものにしない不変量}

全 $j,r,\alpha$ について、$s\ne0$ なら次数零の関数斜体は
$\mathbb C(C)$ 上の第一Weyl斜体である。
\[
 t=x^{-1}\star y,\quad u=q/x^2,\quad v=J/x^d,\quad
 [t,u]=2su,\quad[t,v]=dsv,\quad[u,v]=0
\]
とし、$d=2h+1$ に対し $C=v^2/u^d$、$\eta=v/u^h$ と取れば
$\partial=\eta^{-1}t/s$ は $[\partial,\eta]=1$ を満たす。
有理Poisson中心と量子関数斜体の中心はともに $\mathbb C(C)$ である。
したがって多項式中心が定数・平面・Du Val曲面のどれであるかは、
この射影的な関数斜体の同型型には残らない。
異なる $m$ の射影圏は $K_0$ の階数 $m+3$ で区別できる。
同じ $m$ の異なる $j$ の射影圏を座標環の中心だけで区別することはできない。
公表稿Xの固定変形環と反標準Hochschild--Serre環を用いると、本族内部の全ての $j$ を区別できる。

補助次数 $(0,1,2,d)$ による変数の拡大縮小は、非零の $s$ を互いに移す。
また $X_{\alpha+\beta}=X_\alpha+\beta E$ とVandermonde恒等式から
\[
 a\star_{\alpha+\beta}b=T^{\beta\deg b}(a)\star_\alpha b
\]
となる。固定した $j,r$ の全 $\alpha$ はZhang twistで結ばれ、射影圏は同値である。
通常のCY性は、この同じ射影圏の中の特定の次数付き代表を選ぶ条件でもある。

\subsection{$c=0$ と残る大域分類}

重み三で $c=0$ の被約部分は、ちょうど次元 $35,20,18$ の三成分を持ち、
全スキームではそれぞれ $\operatorname{PB}(\mathbb P^2)$、$L(1,1,4)$、
$L(1,1,1,3)$ に含まれる。
従って $c=0$ に支持される未知の大域成分はない。
ここでの十八次元部分は、接線型の二つの大域成分とは別である。
再利用する事実は、巡回型の $D$ に対して
$(\mathbf x\times D)\cdot F=0$ の全多項式解が $F=f\mathbf x+gD$ と書けることである。
商体上の解を高さ一の局所環で確認し、正規性で多項式に戻す。

接線型の非零 $D$ で残っていた第二共鳴は完了した。
大域分類の残りは非被約三次形式、その他の特異三次形式の特殊係数、
およびそれらの全混合核の大域成分への所属である。
十成分で分類が尽きることと最大次元は、なお未確定である。

% Reusable results and decision criteria; derivations are in publication paper IX.
\section{重み三のPoisson空間：境界と留数の再利用}
\label{sec:poisson-research-notes-v32}

\subsection{全方程式と、商を取らないという規約}
\[
 c\in S_3,\qquad D\in\mathfrak{sl}_3,\qquad
 [F]\in S_4^{\oplus3}/S_3x,
 \qquad S_d=\mathbb C[x,y,z]_d
\]
を座標とし、
\[
 D(c)=0,\qquad
 M_{c,D}(F)=2F(c)-(x\times D)\cdot F-c\operatorname{div}F=0
\]
を課す。全ての分類は、この五十三変数・三十八二次式のスキームから出発する。
三十八は最小二次生成元の個数であり、余次元ではない。
以下の成分次元はスカラー倍や重み付き自己同型で商する前の値である。
基底 $D(c)=0$ が非被約なら、その冪零方向を混合写像に残したまま計算する必要がある。

再利用できる幾何的辞書は
\[
 \alpha=c\,dw-w\,dc+w(x\times D)\cdot dx+(F\times x)\cdot dx
\]
である。これは重み六の水平一形式で、Poisson方程式は
$\alpha\wedge d\alpha=0$ とスキーム論的に同値である。
三次式による $w$ の移動は、混合核に実際の核方向を作る。
固定した一次方向の量子化と、指定済みVK族の高次係数の補完とは区別する。

\subsection{認定済み成分の表}
\begin{center}\small
\begin{tabular}{lrl}\toprule
成分&次元&再利用する特徴\\\midrule
$\mathbb P^2$ からの引戻し&35&一般点で $c=D=0$\\
$R(1,5)$&28&直線・二次曲線の共鳴比 $a=2$\\
$L(1,1,4)$&23&対数三角形の一つの共鳴壁\\
$R(2,4)$&22&直線・二次曲線の共鳴比 $a=-1$\\
$R(3,3)$&19&一般の $c$ が滑らか、$D=0$\\
$L(1,2,3)$&19&一般の直線・二次曲線型\\
$L(1,1,1,3)$&19&一般の非共点三直線型\\
$\mathbb P(1,1,3)$ からの引戻し&19&一般点の横断環の長さ六\\
接線・二次曲線型 $a=-5/2$&18&一般点で滑らか、有理第一積分\\
接線・二次曲線型 $a=-3/2$&18&一般点の長さ二、量子中心は $D_7$ 型\\\bottomrule
\end{tabular}
\end{center}
十成分のうち八成分は一般点で滑らかであり、残る二成分の長さは六と二である。
十個が全てであること、三十五が最大次元であることは未証明である。
$R,L$ は古典的な有理型・対数型の構成を表し、構成自体を新規性とはしない。
本研究で必要なのは、この重みのJacobiスキームでの成分性・厚み・交叉の認定である。
証明と引用の位置付けは公表稿IX\cite{PaperIX}にまとめる。

\subsection{共点三直線の重み三での数値}
節~\ref{sec:allm-notes-v33}の全重みの結果を $m=3$ に特殊化すると、
\[
 \mathbb C[[s_1,\ldots,s_{18},a,b,u]]/(ab,b^2,a^5u)
\]
となる。一般点の長さは六、primary成分の交叉の長さは五、一般切断の長さは七である。
元の有理数証明書も比較用に保持する。
$(A,B)=(2,5)$ で、自由な二変数四次係数を $h_{ij}=[x^iy^j]F_z$ とすると、
Schur消去の非零式は
\[
 a^5L,\ a^6L,\qquad
 L=(20h_{13}-100h_{22}+300h_{31}-700h_{40})/9
\]
であり、$F=(0,0,x^4)$ の係数は $-700/9$ である。
この式は全重み検証器が従来の規約と一致することを確認する基準になる。
全重みの証明ではこの特定の係数値を必要とせず、重量と三交点評価によって
先頭係数の非零性を保証する。

\subsection{非共点三直線：五つの局所枝}
$c=3xyz,D=F=0$ の完備局所環は
\[
 \mathbb C[[z_1,\ldots,z_{16}]]\widehat\otimes
 \frac{\mathbb C[[a_i,\delta_i,\rho_i\ (1\le i\le3)]]}
 {(\sum\delta_i,\sum\rho_i,a_i\delta_i,a_i^2\rho_i,
     \delta_1\rho_2-\delta_2\rho_1)}.
\]
ここでは $a_1\rho_1$ が非零の平方零元である。
被約化の局所枝は、$R(3,3)$、$L(1,1,1,3)$、および
$L(1,2,3)$ の三つの局所枝で、全て十九次元となる。
後者を別々の大域成分と数えない。
交叉次元は、三次曲線型と対数型が十六、三次曲線型と各直線・二次曲線枝が十七、
対数型と各直線・二次曲線枝が十八、異なる直線・二次曲線枝どうしが十六である。
残余二方向は既知の対数型の境界に収まり、この点だけから新たな大域成分は生じない。

\subsection{被約な共鳴交叉と明示量子化}
\[
 c=x(x^2+yz),\qquad D=a(y\partial_y-z\partial_z)
\]
では混合核次元が、一般の $a$ で十、$a=2$ で二十、$a=-1$ で十四となる。
局所環は滑らかな十八変数を除いて
\[
 \mathbb C[[t,v_1,\ldots,v_r]]/(tv_1,\ldots,tv_r),
 \qquad r=10\ (a=2),\quad r=4\ (a=-1)
\]
である。一般対数三角形の壁 $3\alpha+t=0$ では同じ形の $r=5$ が現れる。
これらは全て被約であり、共点三直線の厚みとは異なる局所模型である。

$r=5$ の一般点は、二変数四次式 $f$ と二つの非零係数 $a,b$ による
\[
 S=\frac{\mathbb C\langle y,z,x,w\rangle}
 {(yz-zy,\ xy-ayx,\ xz-bzx,\ wy-a^{-1}yw,\ wz-b^{-1}zw,\ wx-abxw+f)}
\]
で量子化できる。これは重み $(1,1,1,3)$ のPBW基底を持つnoetherian整域で、
AS正則次元四、通常の意味で4-Calabi--Yau、Gorenstein parameter六である。
補助次数 $\deg' x=1,\deg' w=-1,\deg' y=\deg' z=0$ は、
Ore定理で未決のNakayama平行移動項を消すために有効である。
Orlovの定理により標準六項列を持ち、$S_{(0,1,3,4)}$ の圏は実際の半直交成分となる。
任意の形式量子化のAS正則性や、全てのPoisson成分の明示量子化はまだ示していない。

\subsection{$D=0$ の全混合核}
非零三次式 $c=\prod f_i^{m_i}$ に対し
\[
 c_{\mathrm{red}}=\prod f_i,\quad d=\deg c_{\mathrm{red}},\quad
 h=c^2/c_{\mathrm{red}}
\]
と置く。Euler商の全解は
\[
 \frac{(F\times x)\cdot dx}{c^2}
 =d\left(\frac ph\right)+\sum_i\lambda_i\frac{df_i}{f_i},
 \quad p\in S_{6-d}/\mathbb Ch,\quad
 \sum_i\deg(f_i)\lambda_i=0
\]
で与えられる。核次元は $\dim S_{6-d}+r-2$、三次移動の像は常に九次元である。
既約三次、直線と二次曲線、相異なる三直線、二重直線と別の直線、三重直線の順に
核次元は $9,10,11,15,20$ となる。
ただしこのファイバーの次元増加だけでは、全Poisson空間の成分性は分からない。
$D$ を含む横断変形と混合方程式を同時に調べる必要がある。

\subsection{標識と自己同型を保持するための注意}

三次式による $w$ の変更で混合項を消す議論は、重み付き座標変換を許した標準形である。
その座標変換が特殊ファイバー上で恒等的とは限らないので、
標識付き変形の同値と取り違えてはいけない。
また、二次 Poisson 錐と外側の miniversal hull の同型は、
普遍量子化と hull の座標の選択に依存する。
指定された VK 族の補完条件をその座標で書く際には、その族を表す高次係数も追跡する必要がある。

\subsection{別の対象へ一般化する際の量子化経路}

今回の重み付き射影スタックでは、BV 恒等式
$\delta(E\wedge U)+E\wedge\delta U=6U$ によるアフィン持ち上げが最短の証明である。
Euler 場を持たない対象への一般化には、Pridham の量子化定理
\cite[Corollary~3.38]{Pridham2025} が別の入力になる。
同定理は完全な余接複体を持つ導来 Artin スタックの Poisson 構造を量子化し、
強準コンパクトの場合には perfect 圏の curved $A_\infty$ 変形を与える。

強例外列がある場合、各対象の曲率を逐次除く障害は自己 $\operatorname{Ext}^2$ に入る。
これが消滅すれば対象を持ち上げられ、全 morphism cohomology が零次に集中する場合には
最小 $A_\infty$ 構造に $m_2$ 以外の演算が残らない。
この経路は一般化に有用だが、本稿の weighted lift の証明の代わりに
曖昧な「スタック上の formality」を引用する必要はない。

\subsection{変形関手の比較をさらに強める場合}

発散零の重み零多重ベクトル場は Schouten 括弧で閉じ、各次数で
$H^0(X,\wedge^pT_X)$ と同一視できる。
アフィン Kontsevich $L_\infty$ 写像をこの部分代数に制限し、
得られる重みを保つ polydifferential cochain を六対象の可合成列に評価すれば、
外側の Hochschild 複体への $L_\infty$ 写像を構成できる。
これを用いて変形groupoidそのものの比較を述べるには、
評価の第一成分と untwisted HKR の自然性、零次 cochain から生じる内部自己同型、
および負次数を扱う規約を明記する必要がある。
現時点の hull 同定には、この追加のgroupoidレベルの主張は不要である。


\section{量子射影圏の離散分類と次に残る変形問題}
\label{sec:categorical-notes-v35}

\subsection{族内で完全に決まった同値関係}

$m=2j+r\ge3$、$j,r\ge0$、$s\ne0$ とし、CGG族の射影圏を
\[
 \mathcal Q_{m,j}(\alpha,s)
 =\operatorname{qgr}\mathcal S_{j,m-2j}(\alpha,s)
\]
と書く。本族内部では、$\mathbb C$ 線形アーベル圏としての同値と
有界導来圏の $\mathbb C$ 線形三角同値は、どちらも
\[
 \boxed{(m,j)=(m',j')}
\]
によって完全に特徴付けられる。固定した $m$ には
$\lfloor m/2\rfloor+1$ 個の相異なる射影圏がある。
$\alpha$ と非零の $s$ は圏の同値類を増やさない。
一致する場合の同値はZhang twistと非零量子化係数の次数付き同型で与えられる。

逆向きには $K_0$ と全共鳴のhullを使う。
標準full exceptional collectionから $K_0\simeq\mathbb Z^{m+3}$ となるため、まず $m$ が決まる。
$r>0$ では標準有限代数のminiversal環
\[
 R_{m,j}\simeq\mathbb C[[\gamma]]/(\gamma^{j+1})
\]
の長さが $j$ を決める。
偶数重みの端点 $r=0$ は、$m+2$ 次元の正則なhullを持つため、
Artin環を持つ残りの共鳴から分離される。
反標準環は、分類に加えて射影圏の曲線幾何を与える別の内在的不変量である。
Hochschild複体の高次構造とminiversal環の不変性には
Kellerのderived invarianceを用いる\cite{KellerDerived2003}。
三角同値が存在すれば、tilting対象の像とRickardの定理から
何らかのdg Morita同値が存在するので、元の同値のdg持上げを仮定する必要はない。

全族が同じ次数零のWeyl関数斜体を持つことは、この分類と両立する。
この例では双有理な関数体を固定しても、射影導来圏には離散的な相違が残る。

\subsection{射影圏に内在する反標準環}

$\ell=m+3$ とし、通常CY代表
$\alpha_0=-(m+r+3)/(m+3)$ のdg射影圏を $\mathcal C$ とする。
Serre関手は $\mathbb S=(-\ell)[3]$ である。
Lin--ZhangのHochschild--Serre環\cite{LinZhangSerre2025}の対角
\[
 \operatorname{AC}(\mathcal C)
 =\bigoplus_{n\ge0}
 \operatorname{Hom}_{D(\mathcal C^{\mathrm{op}}\otimes\mathcal C)}
 (I_{\mathcal C},\mathbb S^{-n}[3n])
\]
は、構造層や次数移動を標識しない導来不変量である。本族では
\[
 \operatorname{AC}(\mathcal C)\simeq Z(\mathcal S_{j,r}(\alpha_0,s))^{(\ell)},
\]
右辺の次数を $\ell$ で割る。
\[
\begin{array}{c|c|c}
 \text{条件}&\text{反標準環}&\text{生成次数}\\\hline
 r\ge2&\mathbb C&\text{正次数なし}\\[1mm]
 r=1&\mathbb C[S,P,R]/(R^2-SP^2-\frac{16}{81}S^{m+3})
 &(2,m+2,m+3)\\[1mm]
 r=0&\mathbb C[U,V,W]/(UV-W^{m+3})&(2,m+1,1)
\end{array}
\]
したがって特異点型はそれぞれ $D_{m+4}$、$A_{m+2}$ となる。
特に $r=0$ の座標環中心 $\mathbb C[q,J]$ は滑らかだが、
反標準環はそのVeroneseであり、$A$ 型特異点を持つ。
最初の二次数の次元は、$r\ge2,r=1,r=0$ の順に
$(0,0),(0,1),(1,2)$ である。

正次数部分を持つ場合、内在的なstacky Projも計算できる。
$r=1$ は粗空間 $\mathbb P^1$、安定化群位数 $(2,2,m+2)$ のorbifold曲線、
$r=0$ は $\mathbb P_{\mathrm{stack}}(2,m+1)$ となる。
この曲線の構成は、元の非可換圏から曲線への全域的fibrationの構成を含まない。

\subsection{別の族へ移すために保持する証明機構}

反標準環の計算では、標準tilting対象 $T$ に対し
$M_D=\operatorname{RHom}(T,T(D))$ が $D\ge0$ で次数零に集中することを先に使う。
双加群準同型は各対象の $c_i\in A_D$ と可換図式で表される。
正規元 $x$ の射は $c_i=x^ic_0x^{-i}$ を強制し、次数二の射 $q$ は
$[c_0,q/x^2]=0$ を強制する。
局所表示
\[
 A[x^{-1}]=\mathbb C[x^{\pm1},q,J][t;\delta],\qquad
 [t,q/x^2]=2s(q/x^2)
\]
で $t$ の最高次数を見ると、この中心化環は係数環そのものである。
従って全 $c_i$ が一致し、残る生成元の射からその共通値が中心元と分かる。
再利用すべき入力は、正の移動に対するExt消滅と、
有限例外列の中で中心化環を十分に小さくする射が存在することである。

固定量子化の変形環では、可換基点のhullを非零の複素数で評価してはならない。
CGG弧を $h\pi+O(h^2)$ としてJacobi錐へ移し、
$\mathbb C((h))$ 上で $h$ を割って $\pi+O(h)$ の全形式近傍を用いる。
非零点での環 $\mathbb C[[z,\gamma]]/(\gamma^{j+1})$ は、
被約な弧上で $\gamma=0$ となるため、その厚みを保って移送できる。
さらに $h^{-1}\log(1+hN)$ が外的導分を与え、最初のHochschild差分の上界と合わせて
$HH^1=1$ を得る。この次元によって軌道変数の数を確定し、
埋込次元とHilbert--Samuel重複度から残るhullの長さを決める。
形式弧の移送と、軌道方向を除く操作は別々に検証する必要がある。

\subsection{固定CGG圏の変形環は端点で型を変える}
$m=2j+r\ge3$、$s\ne0$ の標準全列有限代数を $A_{j,r}$ とする。
全ての $\alpha$ について、hullは次の二種類になる。
\[
\begin{array}{c|c|c|c}
 &R_{A_{j,r}}&\dim HH^1&\dim HH^2\\\hline
 r>0&\mathbf C[[\gamma]]/(\gamma^{j+1})&1&\mathbf1_{j\ge1}\\
 r=0\ (m=2j\ge4)&\mathbf C[[u_1,\ldots,u_{m+2}]]&1&m+2
\end{array}
\]
$r=0$ の $HH^3$ は $m+5$、高次は零である。
前者には被約な一パラメータ形式変形がなく、後者は $m+2$ 次元の滑らかな形式変形空間を持つ。
同じ明示的なCGG量子化式を使っても、端点の局所変形問題は連続的に外挿できない。

再利用しやすい核心は、非退化二次形式 $q$ に対する
\[
 L(F)=\frac{m+3}{2}F(q)-q\operatorname{div}F,
 \qquad W_m=S_{m+1}^{\oplus3}/S_m\mathbf x
\]
という写像である。偶数 $m=2j$ では
\[
 \ker L=\{[\nabla q\times\nabla H]:H\in S_{m+1}\},
 \qquad\operatorname{coker}L=\mathbf C\cdot[q^{j+1}].
\]

$q$ を平方和にすると、調和分解 $H=q^e h_\ell$ に対して
$L(\nabla H)=\ell(\ell+1)H$ となる。
余核の非零性はFischer双線形形式により全整数について証明できる。

端点切片では全混合式が
\[
 M_{\kappa,\gamma}(F)=xL_\gamma(F)
  +\kappa\left(q_\gamma F_x-\frac x2F(q_\gamma)\right)
\]
であり、$H_0=q^j y$ に対応する核ベクトルの $\kappa$ 微分は
$-xq^{j+1}$ となる。
これは余核で非零であるため、非零混合点での全Schur残差イデアルは $(\kappa)$ となる。
この証明は接空間の一致にとどまらず、全形式局所環の正則性を与える。
Poisson側の次元 $N_{m+1}+8$ から座標変換軌道の次元 $N_m+8$ を除いた
$m+2$ が量子圏のhullの次元となる。


この滑らかなhullは、後述する明示Ore族によって全て実現できる。
Poisson側の有理型表示 $2q\,d(xw+H)-(m+1)(xw+H)\,dq$ では、
$w$ 移動により $H$ を $xS_m$ で割り、非零の倍率で一点の係数を固定できる。
これと直線と二次曲線の接触を解く一方向が、量子側でも独立な $m+2$ 個のパラメータになる。


\subsection{別の族へ進むための条件}

重み三の二つの四項圏はGerstenhaber作用により区別できる。
一般の部分圏の反標準環を調べる場合には、許容埋込み $i$ に対して
Serre関手が $i^!\mathbb S i$ となることを保持する。
外側の次数移動をそのまま制限した関手ではない。
特定の外側への補完軌跡の差を、標識を忘れた部分圏単体の不変量として使わない。

\paragraph{他のPoisson成分における固定量子化の変形。}
引戻し型・対数型などへ進む場合は、明示的な量子化弧、非零係数間の同型、
非零Poisson点の全局所環、量子側 $HH^1$ の四つを先に揃える。
一般点の長さだけでは固定量子化のhullは決まらず、
正次元の横断方向がある場合は全てのsmooth変数がゲージになるとは限らない。
この条件が揃う成分を優先すれば、既存のPoisson成分表を
実際の非可換射影圏のmoduliと高次障害へ移すことができる。

以上の分類は明示したCGG族の内部についての完全分類であり、
全非可換重み付き射影空間の分類ではない。
反標準環は $r\ge2$ 同士を区別せず、その部分を分離するのは固定変形環の長さである。

\section{固定引戻し量子圏のhull：再利用可能な論証}
\label{notes:pullback-fixed-hulls-v37}

\subsection{確定した対象と結論}

$m\ge3$、$N_m=\binom{m+2}{2}$、$c=xy(x+y)$、$L=Ax+By$ とする。
$\deg(x,y,z,w)=(1,1,1,m)$ に対し
\[
 \delta(x)=\frac{c_y}{3}+\frac{m}{3(m+3)}xL,\quad
 \delta(y)=-\frac{c_x}{3}+\frac{m}{3(m+3)}yL,\quad
 \delta(w)=u x^{m+1}-\frac{m}{m+3}Lw
\]
と置き、$S=\mathbf C[x,y,w][z;s\delta]$ の標準 $m+3$ 対象列の有限代数を
$\mathcal A$ とする。$u,s\ne0$ とし、下記の明示開条件を課すと、
\[
 R_{\mathcal A}\simeq
 \mathbf C[[t_1,t_2,a,b]]/(a^{m+2},ab,b^2).
\]
二つの滑らかな変数は実際のOre係数 $A,B$ の変化によって実現される。
$HH^\bullet$ の次元は $m=3$ で $(1,2,4,9)$、$m\ge4$ で $(1,2,4,6)$、
次数三を超える群は零である。
これは全結合変形のhullであり、次数付き座標環の小さなパラメータ族だけの計算ではない。
ただし対象は標準有限代数であって、無限次元環 $S$ の全非次数付き変形とは区別する。

十分な開集合は次の通りである。
\[
 \kappa=\left(1+\frac m3(1-B),\,1+\frac m3(1+A),\,
                  1+\frac m3(1-A+B)\right),
\]
\[
 (B-2A)(A-2B)(A+B)\kappa_1\kappa_2\kappa_3
 [m(B+2)+6]\prod_{k=0}^{m}[m(B+2)-3k]\ne0,
\]
かつ $\kappa$ が和 $m+3$ の正整数三つからなる有限集合を避けることを要求する。
$(A,B)=(2,5)$ は全ての $m\ge3$ でこれを満たす。
この条件集合の外における例外パラメータは分類していない。

\subsection{非零混合点への局所環の移動}

共点三直線での既存の全Poisson局所環は
\[
 \mathbf C[[s_1,\ldots,s_{N_m+8},a,b,u]]/(ab,b^2,a^{m+2}u)
\]
である。$u=0$ での形式級数に $u_0\ne0$ を直接代入することはできない。
再利用すべき性質は、証明全体が混合変数 $F$ に関して線形である点である。
$(c,D)$ の切片は $F$ に依存せず、Schur消去、移動像の直和因子の分離、
残余座標 $u$ の選択は $F$ の係数について線形である。
従って基底 $(c,D)$ だけを完備化し、$F$ は多項式変数のままとして消去できる。
その後 $F=u_0x^{m+1}\partial_z$ で完備化し、$u=u_0+v$ と置くと
\[
 \widehat{\mathcal O}_{\mathscr P_m,\pi}
 =T_m[[z_1,\ldots,z_{N_m+9}]],\qquad
 T_m=\mathbf C[a,b]/(a^{m+2},ab,b^2)
\]
を得る。これが非零混合点での定理の入力となる。

開条件の由来も明示できる。二変数移動の双対角行列は
\[
 V(x^{m-k}y^k)=
 \left(\frac{m(A+1)}3-k\right)x^{m+1-k}y^k+
 \left(\frac{m(B+2)}3-k\right)x^{m-k}y^{k+1}
\]
であり、第二係数の全てが非零なら単射かつ余核は $x^{m+1}$ で生成される。
最高重量 $m+1$ で、源を $z^{m+1}\partial_x,z^{m+1}\partial_y$、
消す標的を $xyz^{m+1},y^2z^{m+1}$ にすると、小行列式は
\[
 -[m(1-B)+3][m(B+2)+6]/9
\]
となる。低い重量の標的全射性は $\kappa_i\ne0$ から分かる。
従って不可視の追加小行列条件を置く必要はない。

\subsection{二つの安定化方向と二つの変形方向}

重み付き座標変換のLie代数は
$\mathfrak{sl}_3\oplus\mathbf Cw\partial_w\oplus S_m\partial_w$
で、次元 $N_m+9$ である。
$c$ の軌道が六次元、移動像が $N_m$ 次元、非零混合係数の倍率が一方向なので、
全Poisson軌道次元は $N_m+7$ である。
安定化Lie代数は $x\partial_z,y\partial_z$ の二次元である。

量子側でも $z\mapsto z+tx$、$z\mapsto z+ty$ は正確な自己同型である。
この二つの導分の生成元上の値は $z\mapsto x,y$ であり、他は零である。
正規順序のPBW多項式上で通常の $x\partial_z,y\partial_z$ と同一視してはならない。
高い $z$ 冪にはOre補正が生じる。
標準有限代数の冪等元を固定する内的導分は対角スカラーから来るので、
この二方向は独立な外的導分となる。
Hochschild複体の最初の差分 $[\pi,-]$ による上界と合わせて $HH^1=2$ を得る。

可換点の基底変換軌道次元を $b_0$ とすると、量子点での値は
$b_\eta=b_0+N_m+7$ である。
全結合積スキームのgeneric局所環には $b_0+N_m+9$ 個の滑らかな変数があるため、
二つがhullに残る。Artin環だけを最終hullとする結論は誤りである。

さらに $A,B$ の二つのKodaira--Spencer類は独立である。
Poisson側では、$c$ を固定する $\mathfrak{sl}_3$ 導分は二つの $z$ 移動だけで、
それらも $w$ 変換も $D$ を変えない。一方 $\partial_A D,\partial_B D$ は独立である。
量子側で可縮部分を除いた複体を取ると、次数一の微分は
$d_h=h(D_1+O(h))$ であり、$HH^1$ が先頭差分とgeneric fiberの双方で二次元なので、
$d_h/h$ と $D_1$ の階数が一致する。
従って $\operatorname{im}(d_h/h)$ は $\mathbf C[[h]]$ 上の直和因子である。
$A,B$ のKSコサイクルを $h$ で割ってから特殊化すると $\partial_A\pi,\partial_B\pi$ となる。
直和因子性によってgeneric fiberでの従属は先頭項の従属を引き起こすので、独立性を得る。
この議論は、最初のコホモロジー差分で階数が確定した明示量子族へ一般に再利用できる。

\subsection{形式移送で必要な三つの一般原理}

第一に、斉次錐の点 $\pi$ の局所環が $T[[z_1,\ldots,z_d]]$、$T$ が有限局所環なら、
$p(h)=h\pi+O(h^2)$ のgeneric局所環も $T_{\mathbf C((h))}[[z_1,\ldots,z_d]]$ である。
$p/h$ を $\pi$ で再中心化し、形式級数を $h$ 進に評価する。
被約環 $\mathbf C[[h]]$ への弧は $T$ の冪零元を殺すため、動くのは滑らかな座標だけである。

第二に、$T[[z_1,\ldots,z_d]]$ 上の滑らかな群作用で閉点軌道次元が $q$ なら、
横断切片を $T[[z_{q+1},\ldots,z_d]]$ に取れる。
被約な軌道上では $T$ の冪零元が消えるので、軌道接空間は $z$ 方向だけに入る。
$q$ 個の群座標を選び、最初の $q$ 個の $z$ を零にする方程式へ形式逆関数定理を使う。
これはArtin厚みの同型型を保存したまま軌道を除く論証であり、
滑らかな変数の数だけを数えて形式的な消去を推測する方法とは区別される。

第三に、複素係数の有限埋込次元完備Noether局所環二つが、拡大体上で形式同型なら、
元の複素係数でも非標準的に同型である。
定義関係、形式同型、逆同型、および関係所属の証人に現れる係数は可算個である。
定義関係の係数を含む可算部分体 $k_0\subset\mathbf C$ を固定し、
残りの係数を添加した可算体を $k_0$ 上で $\mathbf C$ に埋め込む。
$\mathbf C$ の超越次数と代数的閉性がその埋込みを保証する。
全係数恒等式が保たれるため、互いに逆な複素係数の形式写像を得る。
これは $h=1$ を形式級数へ代入する操作ではない。

\subsection{再利用する環の不変量と未決の範囲}

$T_m$ の基底は $1,a,\ldots,a^{m+1},b$、socleは
$\langle a^{m+1},b\rangle$ である。
hullは二次元Cohen--Macaulay、型二、重複度 $m+3$、埋込次元四であり、
非Gorensteinかつ非完全交叉である。
障害の三つの最小関係は $a^{m+2},ab,b^2$ である。
従って純粋な $a$ 方向は次数 $m+2$ で初めて止まり、$b$ は二次で止まる。

残る有用な問題は、$A,B$ 間の大域的な同値分類、明示開集合の外側のhull、
および二次元安定化群の非被約hull上の作用である。
本結果はmarked miniversal環の同型型を決定するが、
自己同型を含む変形groupoidや、大域的なcoarse moduliを分類しない。

検証コードは標準有限代数の矢印像を未知数として、全最小Ore関係のLeibniz式を有理数上で解く。
$m=3,4,5,6$ の $HH^1=2$、二つの明示外的導分、$A,B$ のKS類の独立性、
および別構成のPoisson Lie微分行列の軌道階数を照合する。
これらの有限計算は全重みの構造証明の監査であり、外挿の根拠には使わない。

\section{端点の明示的量子化から残す構造}
\label{sec:notes-endpoint-quantization-v37}

\subsection{再利用できる構成原理}

連結次数付きnoetherian AS正則整域 $B$、非零斉次正規元 $x$、Nakayama自己同型 $\sigma$ が
\[
 ux=x\sigma(u),\qquad B/(x)\ \text{可換}
\]
を満たせば、任意の斉次混合項 $H$ を
\[
 \delta_H(u)=x^{-1}[u,H],\qquad
 S_H=B[w;\sigma,\delta_H]
\]
によって組み込める。ここで $\deg H>\deg x$ とし、$\deg w=\deg H-\deg x$ とする。
左割算が可能なのは $B/(x)$ の可換性による。
選んだ一つの $H$ だけを扱う場合には、その像が $B/(x)$ の中心に入るという仮定で十分である。
中心元 $J=xw+H$ が存在し、Ore拡大のNakayama制限
$\mu_{S_H}|_B=\sigma^{-1}\mu_B=\operatorname{id}$ と
中心の各点固定性を組み合わせると、$\mu_{S_H}(w)=w$ まで決まる。
この方法は、発散零の半古典極限だけを根拠としてCY性を推測するものではない。
混合項が高次数でも、証明にはその次数による計算量の増加がない。

本例では $D=1-\gamma s^2$、
$a=s^2/D,b=s/D,c=(1+\gamma s^2)/D$ として
\[
 yx=cxy-bx^2-2\gamma bxZ,\quad
 Zx=cxZ-2bxy+ax^2,\quad
 Zy=yZ+axy-bxZ.
\]
唯一の重なり $Zyx$ の解消は
$c^2-4\gamma b^2=1$、$c-1=2\gamma a$、
$a(c+1)=2b^2$ に帰着する。
この小さい普遍計算を一度確認すれば、任意の混合次数に対するPBW性は
Ore拡大の構造から従う。

\subsection{生成元の変更と係数の正規化}

$m=2j$、$n=2j+1$、$\epsilon=(-1)^j$ とする。
CGGのCY代表との対応は $\gamma=0$、
\[
 Z=z+sy,\qquad Q=q=xz-y^2,\qquad H_*=Q^j y
\]
である。$H_* \equiv\epsilon y^n\pmod{xB}$ なので、
$G=x^{-1}(H_*-\epsilon y^n)$ とおけば、二変数形への変更は
$w_{\mathrm{new}}\mapsto w_{\mathrm{CGG}}+G$ となる。
左割算を使わない有限表示もある：
\[
 G=\sum_{\ell=0}^{j-1}(-1)^\ell
 Q^{j-1-\ell}(y-sx)^{2\ell}(Z-sy)y.
\]
これは
\[
 Q^j-(-y^2)^j
 =\sum_{\ell=0}^{j-1}(-1)^\ell Q^{j-1-\ell}y^{2\ell}(Q+y^2),
 \qquad Q+y^2=x(Z-sy),
\]
および $y^{2\ell}x=x(y-sx)^{2\ell}$ から従う。
この式は、CGG代表との同型を具体的な有限多項式として実装する際に使える。

全categorical hullの係数には
\[
 H=\epsilon y^n+\sum_{k=1}^{n}h_k y^{n-k}Z^k
\]
を使える。量子化係数の規格化は
\[
 \gamma'=s^2\gamma,\quad h_k'=s^kh_k,\qquad
 (x,y,Z,w)\mapsto(x,sy,s^2Z,s^n w)
\]
であり、$D\mapsto1-\gamma'$ となる。
従って $s=1$ での族は
$\mathbb C[\gamma,h_1,\ldots,h_n,(1-\gamma)^{-1}]$ 上に代数的に定義される。
形式変形を複素数の非零量子化係数へ評価する操作は不要である。

\subsection{Kodaira--Spencer移送で必要な飽和性}

半古典一次類の独立性だけでは、非零量子化における独立性は一般に従わない。
本例では、二段のGaussian消去と既知の次元等号がその不足を埋める。
以下の形で再利用できる。

\begin{lemma}[第一差分の上界が等号の場合の一次類移送]
$O=k[[s]]$、$K=k((s))$ とし、$C^\bullet$ を各次数有限自由な
$O$ 複体とする。特殊化のコホモロジーを $E_1^\bullet$、
最初の非零差分を $d_1$ とし、
$E_2^\bullet=H^\bullet(E_1,d_1)$ と置く。
必要な次数の近傍で二段のGaussian消去ができるとする。
$r=\dim_kE_2^q=\dim_KH^q(C^\bullet\otimes K)$ であり、
$q$ コサイクル $v_1,\ldots,v_r$ の $E_2$ への先頭像が基底を成すなら、
これらの $K$ 上のコホモロジー類は基底を成す。
\end{lemma}

\begin{proof}
特殊化の可縮対を消去すると、複体は
$E_1[[s]]$ 上の微分 $s d_1+s^2d_2+\cdots$ になる。
$K$ 上で微分を $s$ で割り、その後 $d_1$ の可縮対を $O$ 上で消去する。
残る台は $E_2[[s]]$、微分は $s$ で割り切れる。
次数 $q$ のコホモロジー次元は
\[
 \dim E_2^q-\operatorname{rank}_K(d^{q-1}_{\mathrm{res}})
              -\operatorname{rank}_K(d^q_{\mathrm{res}}).
\]
仮定の等号から二つの階数はともに零である。
$O$ は整域なので、対応する二行列は $O$ 上でも零である。
したがって次数 $q$ に境界による追加の商はなく、
先頭係数行列が可逆な $v_i$ は $O$ 上でも基底を成す。
\end{proof}

端点では $q=2$、$r=m+2$ である。
パラメータ微分によるKodaira--Spencerコサイクル自体は $s$ で割り切れ、
それを割ったコサイクルに補題を適用する。
$E_2^2$ の基底は $\gamma$ と
$S_{m+1}/(xS_m+\mathbb C q^jy)$ の二変数単項式である。
従って像の飽和性を別途仮定していない。
第一差分の上界を等号にする端点の $HH^2=m+2$ が、決定的な入力である。

\subsection{中心と斜体：混合項に依存しない部分}

以下は新族の内部での構造計算であり、パラメータ間の導来非同値を主張するものではない。
$x$ は全代数 $S$ では一般に正規元ではないが、
$B$ の正規局所化を使うと
\begin{equation}
\label{eq:notes-endpoint-local-embedding-v37}
 S_H\hookrightarrow B[x^{-1}][J],
 \qquad w\longmapsto x^{-1}(J-H)
\end{equation}
を得る。右辺の $J$ は独立な中心変数である。
Ore関係はこの代入で成立する。
さらに、$\sum_{i=0}^{e}b_iw^i$ の像の最高 $J$ 次係数は
$b_e x^{-e}$ であり、$b_e\ne0$ なら零にならない。
従ってこの写像は単射である。
これは $S$ における $x$ の正規性を仮定した局所化ではない。

\begin{proposition}[非共鳴パラメータでの中心]
$s\ne0$、$D\ne0$ とする。
$\gamma=0$、または
$\gamma\ne0$ で
\[
 t=\frac{1-s\sqrt\gamma}{1+s\sqrt\gamma}
\]
が一の冪根でないと仮定する。
任意の斉次混合項 $H$ について
\[
 Z(S_{\gamma,s,H})=\mathbb C[Q,J]
\]
である。
\end{proposition}

\begin{proof}
まず $\gamma\ne0$ とする。
$u=y+\sqrt\gamma Z+x/(2\sqrt\gamma)$、
$v=y-\sqrt\gamma Z-x/(2\sqrt\gamma)$ とおけば
\[
 ux=txu,\quad vx=t^{-1}xv,\quad
 uv=-Q-\frac{t}{4\gamma}x^2,\quad
 vu=-Q-\frac{t^{-1}}{4\gamma}x^2.
\]
従って $B[x^{-1}]$ は
$R=\mathbb C[Q,x,x^{-1}]$ 上の一般化Weyl代数であり、
$u$ の次数を一、$v$ を負一、$R$ を零次とする直和分解
\[
 R\ \oplus\ \bigoplus_{a\ge1}Ru^a\
   \oplus\ \bigoplus_{a\ge1}Rv^a
\]
を持つ。$x$ と可換な元は、$t$ が一の冪根でないことから零次成分 $R$ に入る。
さらに $u$ と可換であれば、Laurent多項式は
$f(tx,Q)=f(x,Q)$ を満たすため $x$ に依存しない。
従って $Z(B[x^{-1}])=\mathbb C[Q]$ である。

$\gamma=0$ では、$Q=xZ-sxy-y^2$ によって $Z$ を消去できる。
従って $B[x^{-1}]$ は
$\mathbb C[Q,x,x^{-1}][y;\delta]$、
$\delta(x)=-sx^2,\delta(Q)=0$ と書ける。
$y$ 次数正の元の $x$ との交換子は、最高次項から零にならない。
残るLaurent多項式が $y$ と可換である条件は
$-sx^2\partial_x f=0$ なので、やはり中心は $\mathbb C[Q]$ である。

式\eqref{eq:notes-endpoint-local-embedding-v37}の中で、
$S$ の中心元は $x^{-1}$ とも可換であり、
$B[x^{-1}][J]$ の中心に入る。
その中心は $\mathbb C[Q,J]$ である。
逆包含は既知の中心元 $Q,J\in S$ により成立する。
\end{proof}

一の冪根のパラメータはこの命題から除外した。
そこでは $x$ と可換な非零次数成分が残る。後述の中心局所化とMcKay模型が、その拡大を記述する。
除外条件はZariski開条件ではないので、
この命題を「一つのZariski開集合上の中心計算」と読み替えてはならない。

\begin{proposition}[次数零関数斜体の明示表示]
$m=2j$、$s\ne0$、$D\ne0$ とする。
$\gamma\ne0$ では、$S_{\gamma,s,H}$ の次数零関数斜体は
\[
 \mathbb C(C)(X)(U;\tau),\qquad \tau(X)=tX,\qquad
 C=\frac{J^2}{Q^{m+1}}
\]
である。すなわち $\mathbb C(C)$ 上の量子トーラス斜体
$UX=tXU$ で表示される。
$\gamma=0$ では、同じ中心変数 $C$ を係数とする第一Weyl斜体である。
いずれの表示も混合項 $H$ に依存しない。
\end{proposition}

\begin{proof}
式\eqref{eq:notes-endpoint-local-embedding-v37}と $J=xw+H$ から、
全商斜体は $\operatorname{Frac}(B)(J)$ となる。
$\gamma\ne0$ では
$uv=-Q-tx^2/(4\gamma)$ により $v$ を有理的に消去できるので、
\[
 \operatorname{Frac}(B)=\mathbb C(Q)(x)(u;\tau),\qquad ux=txu.
\]
これはPBW表示、または上の一般化Weyl代数表示から忠実な表示となる。

中心の次数一斉次有理元
\[
 \nu=\frac{J}{Q^j}
\]
を使い、$X=x/\nu$、$U=u/\nu$ とする。
これらは零次であり、$UX=tXU$ を満たす。
また
\[
 Q=\frac{\nu^2}{C},\qquad J=\frac{\nu^{m+1}}{C^j}.
\]
従って全ての斉次生成元は、零次斜体と中心変数 $\nu$ から回復できる。
例えば $v$ は $uv$ の式から、$w$ は $xw=J-H$ から回復できる。
よって零次斜体は表示した量子トーラス斜体に等しい。

$\gamma=0$ では、同じ $\nu$ によって
$X=x/\nu$、$Y=y/\nu$ とすれば
$[Y,X]=-sX^2$ である。
\[
 \eta=X^{-1},\qquad \partial=s^{-1}Y
\]
と置くと $[\partial,\eta]=1$ となり、逆に $X,Y$ を回復できる。
これが第一Weyl斜体の表示である。
\end{proof}


関数斜体の表示から、後述のスカラー交換子不変量と整数変換によって全パラメータの同型判定が得られる。
$H$ の方向を検出する問題は、関数斜体の先に残る。

\subsection{検証と射程の固定}

付属コードは $\mathbb Q(\gamma,s)$ 上で唯一の重なり、$Q$ の中心性、
$\sigma$ の逆、対角Ore表示、および係数スケーリングを検証する。
混合項の具体例では左割算、$\sigma$ 導分恒等式、中心 $J$、
CGG生成元との同一視を有理数上で照合する。
全 $m$ のAS正則性とCY性は一般Ore補題によるものであり、
有限個の次数からの外挿ではない。
miniversalityもHochschild複体の構造と次元等号で証明している。

固定 $\alpha$ の無限次元次数付き環と、その標準有限代数を区別する。
一般 $\alpha$ に対して今回得られたのは、特殊ファイバーの有限Zhang同型を
使って同一視したcategorical hullの実現である。
一般混合項に対するCGG自己同型の延長や、
固定 $\alpha$ の全次数付き変形の分類は未解決である。

\section{端点族で再利用できる双有理不変量とPIパラメータ}
\label{sec:notes-endpoint-birational-v38}

\subsection{中心変数を持つ量子トーラス斜体の分類}

$m=2j\ge4$ の端点族では、
\[
 \nu=J/Q^j,\quad \chi=J^2/Q^{m+1},\quad
 Q=\nu^2/\chi,\quad J=\nu^{m+1}/\chi^j
\]
により中心の次数一変数と零次変数を分けられる。
$w=x^{-1}(J-H)$ によって混合項を消すと、零次斜体は
\[
 \mathscr D_0=D_{\mathrm W}(\mathbb C(\chi)),\qquad
 \mathscr D_\gamma=D_t(\mathbb C(\chi)),\quad
 t=\frac{1-s\sqrt\gamma}{1+s\sqrt\gamma}\quad(\gamma\ne0)
\]
となる。$s\ne0$ を固定した場合の $\mathbb C$ 斜体の同型条件は、
零同士、非冪根で同じ $\gamma$、冪根で同じ位数、の三つだけである。
$H$ には依存しない。

必要条件には $G_{\mathbb C}(D)=[D^\times,D^\times]\cap\mathbb C^\times$ を使う。
量子Laurent級数の先頭項は交換子を $t$ の整数冪に送り、
$UXU^{-1}X^{-1}=t$ なので $G_{\mathbb C}(D_t)=\langle t\rangle$ である。
Weyl斜体では、微分作用素の先頭記号が可換体に値を取る乗法写像となるため
$G_{\mathbb C}(D_{\mathrm W})=\{1\}$ である。
中心の全体を固定する同型でなくても、複素数体が固定されればこの不変量を使える。

有限位数での十分条件には、次の小さい整数変換が有効である。
$t$ の位数が $\ell$、$\gcd(k,\ell)=1$ のとき、$kb-\ell a=1$ を選び、
\[
 X'=X^k\chi^a,\quad U'=U,\quad \chi'=X^\ell\chi^b;
 \qquad
 X=(X')^b(\chi')^{-a},\quad \chi=(X')^{-\ell}(\chi')^k
\]
とする。$\chi'$ は中心、$U'X'=t^kX'U'$ で、指数行列の行列式は一である。
従って中心変数一つを許せば、同じ位数の原始冪根の違いは消える。
この同型は $\chi$ を固定しない。
例えば位数五の二つの $\gamma=-5\pm2\sqrt5$ は、
$X'=X^2\chi,\chi'=X^5\chi^3$ により同じ関数斜体を与える。

\subsection{PI次数の集積と検証用の代数的方程式}

$t$ の位数が $\ell$ なら
\[
 Z(\mathscr D_\gamma)=\mathbb C(\chi,X^\ell,U^\ell),\qquad
 [\mathscr D_\gamma:Z(\mathscr D_\gamma)]=\ell^2.
\]
全商斜体に加わる $\nu$ は中心の超越元なので、座標環のPI次数も $\ell$ である。
$\gamma=0$ と非冪根の場合には、中心上次元は無限であり非PIである。
PIであることから、全座標環がその中心上有限であるとは、この議論だけでは結論しない。

$s=1$ に正規化したPIパラメータは
\[
 \gamma=-\tan^2(\pi k/\ell),\qquad
 \ell\ge3,\quad\gcd(k,\ell)=1.
\]
位数二はCayley変換の分母を零にするため現れない。
$k=1$ として $\ell\to\infty$ とすれば、
非PIな端点 $\gamma=0$ にPI次数が非有界なファイバーが集積する。
有限型のminiversal族でPI集合とその補集合はともにZariski稠密であり、
PI集合は構成可能ではない。
この例では、特定のZariski開集合を選ぶだけで全ての冪根を除くことはできない。

厳密計算には、円分多項式 $\Phi_\ell(t)$ と
$z(1+t)^2-(1-t)^2$ の終結式を用いればよい。
平方因子を一度だけ残した多項式の次数は $\varphi(\ell)/2$ である。
$\ell=3,4,5,6$ では、それぞれ
$z+3,z+1,z^2+10z+5,3z+1$ を得る。
付属コードは任意指数の交換子恒等式と逆変換を記号的に確認し、
有限個の位数について円分方程式を有理数上で照合する。
任意位数の主張は整数変換の証明に基づく。

残る幾何学的課題は、同じ関数斜体を持つ $H$ の方向を、
射影座標環または導来圏の不変量で区別することである。
関数斜体の分類を、そのまま導来圏の分類に読み替えてはならない。

\section{一の冪根で残る中心と表現の構造}
\label{sec:notes-root-unity-v38}

\subsection{中心局所化を使う理由}

端点の全変形族 $S_H=B[w;\sigma,x^{-1}[\,-,H]]$ では、
一般の $H$ に対して $x$ は正規元でない。
しかし対角パラメータ $t$ が位数 $\ell$ の一の冪根なら、
$\Xi=x^\ell$ は $Z(B)$ に入り、$\sigma$ が固定し、混合Ore導分も零にする。
従って $\Xi$ は全代数の中心元となる。
これにより
\[
 S_H[\Xi^{-1}]=B[x^{-1}][J],\qquad J=xw+H
\]
を正当な中心局所化として使える。
この手順は、正規元がOre拡大で正規性を失う場合にも、有限位数の冪が
中心性を回復すれば適用できる。

\subsection{周期が中心の特異点と表現次元を同時に決める}

$\beta=(4\gamma)^{-1}$、$U=u^\ell$、$V=v^\ell$ とすれば、
中心の唯一の関係は
\[
 UV=\prod_{i=0}^{\ell-1}(-Q-t^{2i+1}\beta x^2).
\]
$t^2$ の位数が $\ell/\gcd(2,\ell)$ であることから、
\[
 \begin{array}{c|c|c}
 \ell&\text{中心の関係}&\text{特異中心指標の単純加群}\\ \hline
 \text{奇数}&UV=-Q^\ell-\beta^\ell\Xi^2&\text{特異指標なし}\\
 2d&UV=(Q^d+(-1)^d\beta^d\Xi)^2&d\text{ 次元が二つ}
 \end{array}
\]
となる。ここでは $\Xi$ を可逆としている。
奇数位数では全中心上でAzumaya、偶数位数では中心の滑らかな部分が正確にAzumaya部分である。
この結論は一般的なAzumaya部分と滑らかな部分の一致定理に依存しない。
$x$ の固有値を分離する有限エタール被覆上で、周期quiverの可逆辺を縮約して証明できる。

偶数位数の特異ファイバーは $M_d(\mathcal R)$ であり、$\mathcal R$ は
二頂点・各向き二本の矢を持つ平方零代数である。
その全次元 $6d^2$、半単純商の次元 $2d^2$、根基の次元 $4d^2$ は区別して残す。
PI次数は $2d$ のままであり、特異点の単純加群の次元 $d$ と取り違えない。

\subsection{別の中心開部分への移行}

混合項 $H$ の全ての方向は、関数斜体だけでなく上の中心局所化からも消える。
v39では $U=u^\ell$ を反転する別の中心開部分を用い、$\Xi=0$ への延長が
$H\bmod x$ の巡回ノルムを記憶することを決定した。
全中心 $Z(S_H)=Z(B[x^{-1}])[J]\cap S_H$ の決定には、なお
$\Xi=U=V=0$ の深い境界を扱う必要がある。
非次数付き局所化の同型から、全次数付き環や射影圏の同値を推論しない。
根基の平方零模型は、境界上の単純加群の退化や拡張を比較する際の基準模型になる。

\section{偶数冪根の二頂点角とMcKay構造}
\label{sec:notes-root-McKay-v38}

偶数位数 $\ell=2d$ のchartでは、中心の有限エタール拡大
$\lambda^{2d}=\Xi$ を取ると、二頂点角が全体をMorita的に代表する。
$x$ の固有冪等元を $e_i$ とし、
\[
 e=e_0+e_d,\quad g=e_0-e_d,\quad P=eu^de,\quad R=ev^de
\]
と置けば
\[
 g^2=e,\quad gP=-Pg,\quad gR=-Rg,\quad
 PR=RP=f=(-1)^dQ^d+\beta^d\lambda^{2d}.
\]
一般化Weyl代数の次数が $d\mathbb Z$ のときだけこの角に残るので、
これらの関係は完全である。従って
\[
 eA'e\simeq
 \left(\frac{\mathbb C[Q,J,\lambda^{\pm1},P,R]}{(PR-f)}\right)\#\mu_2.
\]
括弧内の可換環は $\partial_\lambda f$ が単元なので滑らかである。
$e_0,e_d$ は、特異ファイバーの二本の $d$ 頂点鎖を一つずつ代表し、
滑らかな中心ファイバーでもfullである。
有限中心上の全閉点でfullであることから、$e$ は大域的にもfullとなる。
従って有限エタール基底変換の後で、chart全体が滑らかな二元群商スタックの
加群圏で記述できる。元のchartの大域的な行列環表示は、この議論には含まない。

特異中心の完成では $Q\ne0$ なので $PR=f$ から $Q$ を消去できる。
二本の鎖の可逆辺を縮約すると
\[
 \widehat A_p\simeq
 M_d\bigl(\mathbb C[[a,b,z_1,z_2]]\#\mu_2\bigr),
 \qquad (a,b)\longmapsto(-a,-b)
\]
を得る。中心は
$\mathbb C[[a^2,ab,b^2,z_1,z_2]]$、横断特異点は $A_1$ である。
作用は余次元一で自由なので、skew group algebraは
正規不変環上の自己準同型環
$\operatorname{End}_{L^{\mu_2}}(L)$、$L=\mathbb C[[a,b,z_1,z_2]]$ と一致する。
最大Cohen--Macaulay性と大域次元四により、完成代数は非可換クレパント解消となる。
その特殊ファイバーが
$M_d((\mathbb C[a,b]/(a,b)^2)\#\mu_2)$ であり、
二単純加群と平方零の四矢を同時に説明する。

この計算の再利用点は、巡回一般化Weyl代数の係数 $a_i$ が周期 $d$ を持ち、
零辺が二つであるとき、反対の二頂点を選ぶだけで、
長い巡回quiverを二元群のskew group algebraに置き換えられることである。
位数を変えても形式Morita型は変わらず、行列サイズだけが変わる。

\section{再利用する計算：特異部分への降下とBrauer周期}
\label{sec:descent-notes-v39}

偶数位数 $\ell=2d$ の中心chartについて、エタール被覆を降下する問題は
特異部分上で完全に解ける。$\Sigma$ の座標環は
$R=\mathbb C[Q^{\pm1},J]$、二単純加群の被覆は
\[
 L=R[\kappa]/(\kappa^2+tQ/\beta)
 \simeq\mathbb C[\kappa^{\pm1},J]
\]
である。$\iota(\kappa)=-\kappa$ とすると
\[
 A_\Sigma\simeq M_d\bigl(L\ltimes(L_\iota\oplus L_\iota)\bigr).
\]
平方零拡大の中心は $R$、根基を除いた商は $M_d(L)$ である。
これは閉点ごとの二頂点quiverに加え、二頂点が大域的には交換されることを記録する。

\subsection*{再利用できる座標と行列}

冪等元の降下には $x$ の固有値より $x^2$ の固有値を用いる。
\[
 z=-\frac{\beta x^2}{tQ},\qquad
 f_i=\frac1d\sum_{k=0}^{d-1}t^{-2ik}z^k,
 \qquad a_i=Q(t^{2i}-1).
\]
$a_i$ は $1\le i<d$ で単元であり、$f_i$ を結ぶ鎖から行列単位が得られる。
$f_0$ の角では $\kappa=xf_0$、$p=f_0u^df_0$、$r=f_0v^df_0$ で、
$p,r$ は $\kappa$ と反可換、全ての二次積が零である。
逆向きの表示に必要な係数は
\[
 c_d=\prod_{i=1}^{d-1}a_i=(-1)^{d-1}dQ^{d-1}.
\]
$x$ を $t^i\kappa$ の対角行列、$u$ を上向きの単位鎖と折返し $p$、
$v$ を逆向きの $a_i$ 倍の鎖と折返し $c_d^{-1}r$ とすれば、
元の代数の関係を満たし、$6d^2$ 個の自由基底を回収できる。

全ての単純加群を与える大域的表示族は $p=r=0$ で得られる。
被覆座標を内部的に読むには
\[
 \kappa=\frac{1-t}{2}\operatorname{Tr}(x)
\]
が有効である。$\operatorname{Tr}(x^d)$ は $d$ の偶奇で退化するので、
一様な座標としては使わない。
$Q$ の原点の周囲の一周は $\kappa\mapsto-\kappa$ であり、二単純加群を交換する。
$\lambda^{2d}=\Xi$ の被覆を $\Sigma$ に制限すると、
$Q=-t^{2i-1}\beta\lambda^2$ で定まる $d$ 個の連結成分に分かれ、各成分はこの二重被覆である。
特に $d$ が偶数のとき、単純加群の被覆は $\Xi$ の平方根ではない。

\subsection*{一般点の障害は全位数に有効}

一般点の中心体と代数は
\[
 K=\mathbb C(Q,J,\Xi,U),\qquad
 D=(\Xi,U)_t,\qquad
 x^\ell=\Xi,\quad u^\ell=U,\quad ux=txu
\]
である。整域性と $\dim_KD=\ell^2$ から指数は $\ell$。
$E=K(\Xi^{1/\ell})=\mathbb C(Q,J,\Xi^{1/\ell})(U)$ とし、
$U$ に関する付値を取れば
\[
 v_U(N_{E/K}(a))=\ell v_U(a).
\]
巡回代数のノルム判定により $[D]^r=0$ なら $\ell\mid r$、
また $U^\ell=N_{E/K}(U)$ なので周期も $\ell$ となる。
この付値の中心を現在のアフィンchart内の因子と同一視する必要はない。

同じBrauer類を持つ代数は、中心上の通常の自己準同型環
$\operatorname{End}_Z(M)$ とMorita同値にならない。
従って完成局所NCCRから大域NCCRへ移る際の障害が明示された。
有限エタール被覆 $\lambda^\ell=\Xi$ はAzumaya部分上でこの類を分裂させ、
位数 $\ell$ は一般点における最小有限分裂次数でもある。
分裂後に得られるのはベクトル束の自己準同型環であり、
その束が大域的に自由であるとは結論しない。

このBrauer計算は、$H$ に依存する別の中心chartを調べる際にも再利用できる。
一般点の記号代数が同じなら周期障害は同じであり、
$H$ の情報を区別するには中心の境界またはそのアフィンモデルを調べる必要がある。

\section{混合項を検出する境界の一般化Weyl表示}
\label{sec:boundary-notes-v39}

\subsection{正規元の選択を変える原理}
$S_H=B[w;\sigma,x^{-1}[-,H]]$、$J=xw+H$ において、
$x^\ell$ を反転した中心開部分は $H$ に依存しない。
一方 $u^\ell$ を反転すれば
\[
 B[u^{-\ell}]=\mathbb C[Q,u^{\pm1}][x;\tau],\qquad
 \tau(u)=t^{-1}u
\]
となり、$H=h(Q,u)+xG$ と一意分解できる。
$W=w+G$ とするだけで
\[
 xW=J-h(Q,u),\qquad Wx=J-h(Q,tu)
\]
という一般化Weyl表示になる。
この変換は、$x$ が全 $S_H$ の正規元であることを必要としない。
$u^\ell$ は本当に全代数の中心なので、中心の境界を扱う正当な局所化である。

\subsection{境界の記憶は有限ノルムに集約される}
$U=u^\ell$、$\Xi=x^\ell$、$\Omega=W^\ell$ と置くと、
\[
 Z(S_H)[U^{-1}]
 =\mathbb C[Q,J,U^{\pm1},\Xi,\Omega]/(\Xi\Omega-P_h),\qquad
 P_h=\prod_{i=0}^{\ell-1}(J-h(Q,t^iu)).
\]
同じノルムは、自由加群 $\mathbb C[Q,u^{\pm1}]$ の
$\mathbb C[Q,U^{\pm1}]$ 上の乗法行列から $\det(JI-M_h)$ と計算できる。
二方法を独立に比較すると、巡回向き、負冪、ノルムの符号を検出できる。

中心は $D=\mathbb C[Q,J,U^{\pm1},\Xi]$ の部分環 $D[P_h/\Xi]$ である。
従って、混合項はアフィン改変の中心 $(\Xi,P_h)$ を動かす。
$Q,U,J$ を固定した $P_h$ の一致は、体内で積の一つの因子が零になることから、
$h'=h(Q,t^ku)$ と同値である。
この復元性は標識付きであり、$J$ の移動を許す分類とは区別する。
特に中心元を $H$ に加えてもOre導分は変わらず、$J$ が移動するだけである。

\subsection{切断間隔が単純加群の次元を与える}
中心点で $\Xi=\Omega=0$ とすると、軌道上の条件
$J=h(Q,t^i\lambda)$、$\lambda^\ell=U$ の成立箇所が切断辺になる。
その個数を $r$、切断後の鎖の頂点数を $d_1,\ldots,d_r$ とする。
基本代数は両向きの巡回箙で、折返しと一周を零にした $\mathcal R_r$ である。
元のファイバーは $\bigoplus(e_i\mathcal R_r)^{d_i}$ の自己準同型代数であり、
\[
 \text{simple dimensions}=(d_i),\qquad
 \dim=2\ell^2-\sum_i d_i^2,\qquad
 \operatorname{nilindex}(\operatorname{rad})=r.
\]
非零辺を縮約する上界と、行列単位・切断矢を用いた明示表現の下界を両方示す。
単純加群の数だけからファイバー全体の次元を推測しない。
$r=1$ の近傍で行列環表示を作ることで、Azumaya性に必要な局所自由性も得られる。
この中心開部分ではAzumaya部分と滑らかな部分が一致する。

\subsection{一般の巡回モノドロミー}
$h$ の巡回安定化群の位数を $g$、$n=\ell/g$ とする。
被約スペクトル因子の正規化は $\mathbb C[Q,u^{\pm g}]$ であり、
$P_h=P_{\mathrm{min}}^g$ となる。
$z=u^g$、$h_i=h(Q,t^iu)$、
$\Delta=\prod_{i=1}^{n-1}(h_0-h_i)$ を用い、追加衝突を除く部分
$A=\mathbb C[Q,z^{\pm1},\Delta^{-1}]$ を取る。
$L=A[\kappa]/(\kappa^g-z)$ 上の角は
\[
 \mathcal T_g=
 L\langle P,R\rangle/
 (Pa-\sigma(a)P,\ Ra-\sigma^{-1}(a)R,\ PR,\ RP,\ P^g,\ R^g),
 \quad \sigma(\kappa)=t^{-n}\kappa.
\]
正規化上の全代数は $M_n(\mathcal T_g)$ である。
その冪零部分を除く商 $M_n(L)$ が、連結な巡回 $g$ 重被覆上の
実際の $n$ 次元単純表現族を与える。
これは任意の $g$ に対する構成であり、追加衝突点では切断間隔の一般則へ戻る。
$v38$ の偶数位数の二単純加群は、この仕組みの二次の場合に対応する。

\subsection{比較に使える最小例と残る問題}
重み $m=4$、冪根位数 $\ell=5$ で、
$H=u^5$ と $H=u^4v$ はそれぞれ
\[
 P_h=(J-U)^5,\qquad P_h=J^5+Q^5U^3
\]
を与える。$U\ne0$ の中心の特異部分の次元は二と一で異なる。
前者の一般境界ファイバーは五つの一次元単純加群を持ち、
後者は $Q\ne0$ で五次元単純加群を一つ持つ。
両者の関数斜体と $\Xi\ne0$ の中心開部分は同じである。
この例は、関数斜体による分類が混合項を検出しないことの具体的な補完となる。

混合項が中心元と $xG$ の和となる場合の全中心と深い境界は、後続節で解決する。
一般の混合項について残る問題は、$\Xi=U=V=0$ の深い境界における全中心と表現、
三つの中心開部分を含む大域的な延長、標識を忘れた座標環・射影圏の分類である。
今回のノルム復元だけから、全 $H$ に対する導来非同値は結論しない。

検証コードは \texttt{verify\_endpoint\_boundary\_v39.py}。
位数 $3,4,5,6,8$ の非可換ゲージ変換と独立な行列式、
位数三から八までの全 $498$ 個の非空切断配置、
および $60$ 個の単項式ノルムを有理数・円分体上で確認する。
一般巡回降下では $32$ 個の $(\ell,g)$ の組について、
$411$ 個の行列恒等式も厳密に検証する。
全位数の主張は原稿中の一般証明による。

\section{全中心と最深点の完備局所障害}
\label{sec:notes-global-center-v40}

対象は $t$ 原始 $\ell$ 乗根、$\ell\ge3$、$\beta\ne0$ の $S_0=B[w;\sigma]$。
より一般に $H=H_c+xG$、$H_c\in Z(B)$ は $W=w+G$ によって同じ代数となる。
任意の $H$ をこの範囲へ還元できるとは主張しない。
公表稿XIの対応節に全証明を記した。

\subsection{全中心と層別}

$R=\mathbb C[Q,x,w]$、$\tau(x)=tx$、$\tau(w)=t^{-1}w$、
$a=-Q-t^{-1}\beta x^2$ とすると、$S_0=R(\tau,a)$ である。
全中心は
\[
 Z_0=\frac{\mathbb C[Q,J,\Xi,\Omega,U,V]}
 {(\Xi\Omega-J^\ell,\ UV-P_\ell(Q,\Xi))},
 \qquad J=xw,\quad(\Xi,\Omega,U,V)=(x^\ell,w^\ell,u^\ell,v^\ell),
\]
\[
 P_\ell=
 \begin{cases}
 -Q^\ell-\beta^\ell\Xi^2& (\ell\text{ 奇数}),\\
 (Q^d+(-1)^d\beta^d\Xi)^2& (\ell=2d).
 \end{cases}
\]
$Z_0$ は四次元の正規Gorenstein完全交叉で、$S_0$ は中心上有限、PI次数 $\ell$。
$T=V(\Xi,\Omega,J)$、$D=V(Q,\Xi,J,U,V)$ とし、偶数時は
$E=V(U,V,Q^d+(-1)^d\beta^d\Xi)$ とする。
被約特異部分は奇数時 $T\cup D$、偶数時 $T\cup E$、Azumaya部分はその補集合である。

全閉ファイバーは、滑らかな点で $M_\ell$、
$T\setminus\{o\}$ と $D\setminus\{o\}$ で $\mathcal R_\ell$、
偶数時の $E\setminus D$ で $M_d(\mathcal R_2)$、原点で次の $\mathcal F_\ell$ となる。
$\mathcal R_r$ は $r$ 頂点の両向き巡回箙で、折返しと一周を零にした代数である。
$T\setminus\{o\}$ の一次元単純加群は
$(u,v)\mapsto(-uv,u^\ell,v^\ell)$ という連結エタール $\ell$ 重被覆、
$D\setminus\{o\}$ のものは $w\mapsto w^\ell$ という同様の被覆を成す。
両者は原点で一つの自明表現へ合流する。

偶数時の $E$ は $D$ に沿って正規でない。
$\alpha=(-1)^d\beta^d$ とすれば、正規化は
\[
 Qa=J^2,\qquad \Xi=-Q^d/\alpha,\qquad \Omega=-\alpha a^d,
\]
という $A_1$ 曲面であり、$D\setminus\{o\}$ の上に $d$ 点を持つ。

より一般に $H=f(Q,u)+xG$ の全中心も決定できる。
$W=w+G$、$J=xW+f$、$\Omega=W^\ell$ と置き、
\[
 N_f(Q,U,J)=\prod_{i=0}^{\ell-1}(J-f(Q,t^iu))
\]
とすれば、上の中心の第二の表示は
\[
 UV=P_\ell(Q,\Xi),\qquad \Xi\Omega=N_f(Q,U,J)
\]
となる。$W^\ell$ の中心性は $U\ne0$ の中心計算と整域性から全体へ降りる。
右辺の候補環は $J$ に関してモニックなので、$R_B[\Omega]$ 上自由となり、
$\Xi$ を可逆にして整域性が分かる。
$D(\Xi)$ と $D(U)$ は正規、残る $V(\Xi,U)$ は次元二だから、
四次元完全交叉の候補環は正規である。
全中心はその環上有限かつ同じ分数体を持つので、整閉性から一致する。
この方法は局所中心の単なる貼合せより強く、最深閉部分を含む全表示を与える。
$f=u^k$ では $N_f=(J^{\ell/g}-U^{k/g})^g$、$g=\gcd(\ell,k)$。
以下の全ファイバー分類は $H=0$ の場合を扱う。
最深点のファイバーと指数は、後述の通り $f\in Q\mathbb C[Q,u]$ にも拡張できる。

\subsection{最深ファイバーの再利用可能な基底}

$e_0=1$、$e_k=u^k$ ($k>0$)、$e_k=v^{-k}$ ($k<0$)、
$b_0=\ell$、$b_k=\min\{\ell,2(\ell-|k|)\}$ と置く。
\[
 \mathcal F_\ell=S_0/(Q,J,\Xi,\Omega,U,V)
\]
の基底は
\[
 \{x^ae_k:|k|<\ell,\ 0\le a<b_k\}
 \ \cup\ \{w^be_k:|k|<\ell,\ 1\le b<\ell\}.
\]
追加の $x$ 指数の上限は、$u^\ell v^{\ell-k}=0$ とその逆向きから来る。
次数を全生成元で $1$ とすれば、根基は正次数部分、最大次数は $2\ell-2$ である。
$w^{\ell-1}u^{\ell-1}\ne0$ だから冪零指数は正確に $2\ell-1$。
唯一の単純加群は一次元で、全次元は
\[
 \dim\mathcal F_\ell=
 \begin{cases}(7\ell^2-6\ell+1)/2&\ell\text{ 奇数},\\
 (7\ell^2-6\ell+2)/2&\ell\text{ 偶数}.
 \end{cases}
\]
付属の検証コードはこの基底上に左生成元行列を直接作り、
元の関係式と全中心関係、循環ベクトル、根基の各層を厳密に検証する。
$3\le\ell\le24$、$\beta=1,2$ の有限検算は一般証明の代わりではない。

\subsection{最深点の完備化でも指数が落ちない理由}

$\widehat Z_0$ に正の重み
\[
 \operatorname{wt}(Q,J,\Xi,U,V)=1,\qquad
 \operatorname{wt}(\Omega)=\ell-1
\]
を置く。$J^\ell$ と $UV$ を消す標準形は互いに素な先頭単項式を持つので、
付随次数環は
\[
 \frac{\mathbb C[Q,J,\Xi,\Omega,U,V]}
 {(\Xi\Omega-J^\ell,\ UV-c\Xi^2)},\qquad c=(-1)^\ell\beta^\ell.
\]
これは整域であり、最低重みは分数体の離散付値を定める。
剰余体は $\mathbb C(q,j,r)$、
$q=\overline{Q/\Xi}$、$j=\overline{J/\Xi}$、$r=\overline{U/\Xi}$。
一般点の代数 $(\Xi,U)_t$ について、付値完備化で
$x^\ell=\Xi$ は次数 $\ell$ の完全分岐拡大となる。
$U^a$ がノルムなら単元部分の剰余 $r^a$ が $\ell$ 乗となる必要があり、
$r$ の付値から $\ell\mid a$ を得る。
従って完備局所分数体上でもBrauer周期と指数はともに $\ell$ である。

ゆえに最深点の完備代数は、中心上の非零有限反射的加群の自己準同型環にも、
その中心線形Morita型にもならない。
これは $\Xi\ne0$ の閉点完備化における分裂した局所模型と両立する。
後者では $\Xi$ の定数項が非零であり、完備環内で $\ell$ 乗根を取れるためである。
中心自身についての別のNCCRの不存在や、任意の二変数混合項 $H$ の全中心は結論していない。

\subsection{$Q$ 可除混合項への拡張と全中心の区別}

$f=Qg$ なら $\delta_f=Q\delta_g$ なので、最深中心ファイバーは同じ $\mathcal F_\ell$ である。
さらに $f$ の中心項 $cQ$ を $J$ の変更で除くと、残る各単項式 $Q^au^k$ は
$a+k/\ell>1$ を満たす。$\operatorname{wt}(u)=1/\ell$ によって巡回ノルムを展開すれば、
$N_f=J^\ell+\text{高重み項}$ となる。
二つの中心関係の初期形は上記と同じ正則列なので、完備局所付値とノルム障害も同じになる。
従って最深点のBrauer周期・指数は依然として $\ell$。
これは全ファイバーが不変であるという主張ではない。

実際、$m=4,\ell=5,\beta=1$ で $f_0=0$ と $f_1=Qu^3$ を比べると、全中心はそれぞれ
\[
 \Xi\Omega=J^5,\quad UV=-Q^5-\Xi^2,
 \qquad
 \Xi\Omega=J^5-Q^5U^3,\quad UV=-Q^5-\Xi^2
\]
となる。前者の特異部分は次元二、後者は
\[
 Q=J=\Xi=0,\qquad UV=U\Omega=V\Omega=0
\]
という三本の軸の合併で次元一である。
従って、選んだ開部分や座標の標識を忘れても全中心は非同型である。
両代数の $D(\Xi)$ 局所化は $B[x^{-1}][J]$ で一致するが、
中心を保存するアフィンMorita同値・dg導来Morita同値は存在しない。
この区別は射影的な $\operatorname{qgr}$ の導来非同値を意味しない。
最深ファイバーも完備局所Brauer指数も同じままなので、
全中心の特異部分の次元が独立した識別手段になっている。

付属コードは両中心のJacobianイデアルを厳密に計算する。
三軸の被約イデアルの各生成元の冪がJacobianイデアルに属することと、逆の包含を確認し、
被約特異部分の等式を有限点の標本に頼らず検証している。

\section{境界のモノドロミーをBrauer留数として読む}
\label{sec:notes-boundary-ramification-v40}

\subsection{零ファイバーの巡回降下を全形式近傍へ延ばす}
$h$ の安定化群位数 $g$、$n=\ell/g$、$z=u^g$ を用いる。
正規化境界の一般点の係数体は $F=\mathbb C(Q,z)$、$U=z^n$ である。
$s=J-h_0$、$b_i=s+h_0-h_i$、$E=\prod_{i=1}^{n-1}b_i$ と置くと、
$E$ は単元であり、中心は $\xi\upsilon=s^g$、$\upsilon=\Omega/E^g$ と直せる。

零ファイバーで $PR=0$ だった角を、全形式近傍では $PR=RP=s$ とする。
具体的には $L=F[\kappa]/(\kappa^g-z)$、$\sigma(\kappa)=t^{-n}\kappa$ 上で
\[
 Pa=\sigma(a)P,\qquad Ra=\sigma^{-1}(a)R,\qquad PR=RP
\]
を課し、$(P^g,R^g,PR)$ で完備化した $\mathcal T_g$ を取る。
全代数は $M_n(\mathcal T_g)$ である。
行列表現は $u=\operatorname{diag}(t^i\kappa)$、
$x$ の内部下辺は $1$、巻戻り辺は $P$、
$W$ の内部上辺は $b_i$、巻戻り辺は $R$ である。
この規約で $x^n=P I$、$W^n=E R I$ となる。
零ファイバー専用の行列の係数をそのまま流用せず、$PR=s$ を保つ必要がある。

\subsection{閉点完成と一般点完成の区別}
複素閉点で $z\ne0$ なら、その根を完成係数環の中に取れる。
従って完成代数は
$M_n(\mathbb C[[a,b,q_1,q_2]]\#\mu_g)$ となり、
$\mu_g$ は $a,b$ に互いに逆の指標で作用する。
横断中心は $A_{g-1}$ 型で、この代数は完成中心のNCCRである。

一方、一般点の $F$ には $z$ の $g$ 乗根がない。
角の一般点は巡回代数 $(L/F,\sigma,\xi)$ となる。
最小解消の例外付値は
\[
 v_i(\xi)=i,\quad v_i(\upsilon)=g-i,\quad v_i(s)=1.
\]
従って $\chi(\sigma^{-1})=1$ の規約で留数は $i\chi$、位数は $g/\gcd(g,i)$ である。
$Pa=\sigma(a)P$ を採用したため、cup積に対応する指標では逆元を指定する。
特に最初の例外因子の留数は位数 $g$ で、可換解消上のAzumaya延長を妨げる。
その留数を消す最小の有限係数拡大は $L/F$ であり、
前版の単純加群を分離する被覆と一致する。
この議論は、閉点の形式NCCRから境界上の大域的なEnd表示を推論できない理由を示す。

\subsection{検算に残す式}
トーリックchartは
\[
 \xi=a_k^{k+1}b_k^k,\quad
 \upsilon=a_k^{g-k-1}b_k^{g-k},\quad s=a_kb_k
 \quad(0\le k<g).
\]
隣り合う格子ベクトルの行列式は一である。
$g=6$ の留数位数列は $(6,3,2,3,6)$。
検証コード \texttt{verify\_boundary\_ramification\_v40.py} は
非零の $PR$ を保った全形式行列、独立な可逆辺係数、全中心冪を厳密に確認する。
一般証明は公表稿XIの「境界の形式McKay模型と例外因子上のBrauer留数」に記した。
追加衝突点の形式代数は、本節の同一値が周期的に現れる場合の議論だけでは決まらない。

\section{中心特殊化の再利用できる区別}
\label{sec:central-fiber-notes-v35}

本節では中心・射影chart・導来不変量の違いと、再利用する局所表示を保持する。
以下では $r=1$、$m=2j+1\ge3$、通常のCY値、$s\ne0$ とする。
$\mathcal S$ の中心を $Z=D_{m+4}$、その原点を $\mathfrak m_0$ と書く。

\subsection{中心は有限な代数族の底ではない}

非原点での $\mathcal S/\mathfrak m_\chi\mathcal S\simeq\mathcal D(\mathbb G_m)$ により、
$\mathcal S$ は $Z$ 上有限加群ではない。
実際、有限加群なら各閉点の特殊化は有限次元でなければならない。
従って、この中心の滑らかな部分をAzumaya locusとして扱うことはできない。
今回得た幾何は、二次元の微分作用素環を各滑らかな中心点上に持つ族として解釈すべきである。
有限次元単純表現は全て中心原点上に集まり、$x=y=0$、$(z,w)\in\mathbb A^2$ で与えられる。
これは滑らかな中心上の単純性と矛盾しない。

\subsection{原点上のもう一つの局所化}

$x$ を反転した原点fiberの中心は $\mathbb C[c]$ であり、
$c\ne0$ の各特殊化は $\mathcal D(\mathbb G_m)$、$c=0$ は
冪零根基を除くと同じ環になる。
blowupの残る一点に対応する計算は、$x=0$ としてから $y$ を反転すれば得られる。

\begin{lemma}
原点fiberの $x=0,y\ne0$ 層は
\[
 \left(\mathcal S/(\mathfrak m_0\mathcal S+(x))\right)[y^{-1}]
 \simeq\mathcal D(\mathbb G_m)
\]
である。
\end{lemma}

\begin{proof}
$x=0$ の商では $q=-y^2$、
\[
 H=w+\frac{(-1)^j}{4}y^{m-2}z^2,
 \qquad \tau=(2y)^{-1}\star z
\]
を用いて
\[
 (\mathcal S/(x))[y^{-1}]
 =\mathbb C[y^{\pm1},H][\tau;\delta],\qquad
 \delta(y)=-\frac{s}{m+3}y,\quad
 \delta(H)=\frac{s(m+6)}{m+3}H
\]
と書ける。
中心生成元はこの商で
\[
 \mathsf S=-\frac43(-1)^j y^{m+6}H,\qquad
 \mathsf P=0,\qquad
 \mathsf R=-\frac49\mathsf S^{(m+3)/2}
\]
となる。$y$ は可逆なので、原点への特殊化は正確に $H=0$ を課す。
従って $\eta=y$、$E=-(m+3)\tau/s$ により所望の微分作用素環になる。
\end{proof}

ここでの $y$ の反転は $x=0$ の商で行う。
元の $\mathcal S$ の中で $y$ の冪がOre集合であるという主張は不要である。

\subsection{局所化と導来不変量の区別}

全次数付き局所化の中心
$Z(\mathcal S[x^{-1}])=A_{n-1}$ は、正規元 $x$ を選んだ後の不変量である。
一方、射影chartの次数零環の中心は $\mathbb C$ である。
従って前者の $A$ 型の添字だけを使って、標識を忘れた射影導来圏を分離することはできない。
射影圏に内在する分離には、Serre双加群から得る反標準環と固定変形環を使う。
特に $r\ge2$ 同士は反標準環が全て $\mathbb C$ となる。
その区別は、公表稿Xで求めた固定変形環の長さ $j+1$ が担う。

\subsection{文献上の位置付け}

微分Ore拡大の単純性判定は既存理論であり、
\"Oinert--Richter--Silvestrovの一般結果と整合する。
JordanのOre拡大とPoisson primeの比較は、中心指標を動かす全スペクトルの解析にも関係する。
今回独自に計算した部分は、本族の局所表示、
$D_{m+4}$ から $A_{m+3}$ への中心blowup chart、
二chartからなるcrepantな部分解消、
および元代数の全ての非原点central reductionを
拡大体を導入せず $\mathcal D(\mathbb G_m)$ と同定したことである。
一般のOre単純性やDu Val部分解消の存在自体に新規性を主張しない。

有限体上または有理数上での計算証明書は、本文の一般証明の代用ではない。
第35版の中心側verifierは全て有理数で、
535個の明示恒等式と12778個の半群単項式を照合する。
反標準側verifierの有限体系は、標数零の明示的下界と組み合わせて有限個の次数を照合する。
証明書と再現コードを保存する。これらの過去の検算は、新しい数学的リスクがある場合に再実行する。

\section{Primitive分類から次に使える構造}
\label{sec:primitive-notes-v36}

本節では $r=1$、$m=2j+1\ge3$、通常CY値の次数を忘れた量子代数 $\mathcal S$ を扱う。
非零の量子化係数は同型により移せるが、$\alpha$ を変えるZhang twistは
この次数を忘れたprimitive spectrumを保存するとは限らない。
公表稿XIの分類は、この対象の範囲で使う。

\subsection{単純な正規局所化は収縮のprimitive性に十分}

可算次元複素整域 $A$ の非零正規元 $u$ について、$A[u^{-1}]$ が単純ならば $A$ はprimitiveである。
再利用できる理由は、局所化後の単純加群の制限ではなく次の環論である。
非零イデアルは $u$ の冪を含み、零イデアル以外の全primitive idealは $u$ を含む。
一方、可算次元性と基礎体の非可算性からJacobson radicalの元は代数的、従って冪零であり、
整域ではradicalは零である。従って零イデアル自体がprimitiveとなる。
この議論は同時に零イデアルの局所閉性も与える。

\subsection{境界の有限整拡大を使う方法}

$\Delta=0$、$qJ=\lambda x$、$\lambda\ne0$ の $x$-saturationに対し、
\[
 q=-u^2,\quad J=2u^{m+2}/3,\quad
 x=-2u^{m+4}/(3\lambda),\quad y=uv
\]
という有限整拡大が有効である。
$\varepsilon=(-1)^j$、$\kappa=-2/(3\lambda)$ とすると
\[
 v^2=1+\kappa u^{m+2}z,
 \qquad
 w=-\frac{\varepsilon}{3}u^{m-2}z^2
                 \frac{v-2\varepsilon}{(v-\varepsilon)^2}.
\]
有限な境界点では $u=0,v=-\varepsilon$ が強制されるので、$w=0$ になる。
この方法はsaturationの巨大なGroebner基底を全重みで求めることを避け、
全ての $m$ に一つの境界証明を与える。
重要なのは $m-2>0$ であり、$m=1$ へそのまま外挿してはいけない。

\subsection{原点の根基と多項式イデアルへの帰着}

$\mathfrak m_0$ を $D_{m+4}$ 中心の原点とすると
\[
 \operatorname{Jac}(\mathcal S/\mathfrak m_0\mathcal S)
 =\operatorname{Nil}_*(\mathcal S/\mathfrak m_0\mathcal S)
 = (\Delta)/\mathfrak m_0\mathcal S,
 \qquad \mathcal S/(\Delta)\text{ は整域}
\]
である。$\Delta$ の正確な冪零指数は $(m+7)/2$ であり、次の所属式と評価で決定できる。
$x$ 局所化では指数が $L+1=(m+5)/2$ なので、これは全環の指数の下界となる。
中心生成元は $N$ 不変かつ $X$ 重量零なので、元の問題は通常の可換多項式環での
\[
 \min\{n:\Delta^n\in
       (q^2\Delta,\,xJ\Delta^L,\,q(x^2\Delta+J^2)\Delta^L)\}
\]
の計算に帰着する。局所化の指数を無条件に全環へ戻してはいけない。

\subsection{原点の位相から全域の葉空間へ}

中心原点上の $\mathbb P^1$ 族の一点閉包には同じ直線 $w=0$ が加わり、
その任意の無限部分集合は原点ファイバー全体に稠密である。
全正共鳴で量子・Poissonスペクトルの対応も得られ、
非原点を動く中心指標を含めた全閉包は、後述の葉の飽和とZariski閉包の公式で記述できる。
部分解消上の代数の層と単純加群自体の分類は、これらの結果とは別の問題である。

\subsection{局所化で一段消える原点根基の厚み}

$r=1$ の通常CY族では、原点根基 $\mathcal R=(\Delta)/\mathfrak m_0\mathcal S$ の
正確な冪零指数は $j+4=(m+7)/2$ である。
$x$ 局所化後の $j+3=(m+5)/2$ は全環の指数ではない。
再利用する機構は、星積イデアルをPBW多項式イデアルと同定してから、
上界を明示的な所属式、下界を可換多項式の評価で分離することである。

$f=xyz-2y^3/3$、$g=z^2(4xz-3y^2)/9$ とすると
\[
 \Delta^2-x(xw^3+q^jyz w^2)J=q^2T_j
\]
であり、$T_j$ は公表稿XIの明示式で与えられる多項式である。
これに $\Delta^{j+2}$ を掛けて上界を得る。
$y=z=0$ では原点イデアルが $(x^{2j+7}w^{2j+5})$ へ写り、
$\Delta^{j+3}$ の像がそのイデアルに入らないので下界を得る。
評価を非可換環の準同型として用いない点を維持する。

この指数の問題は解決済みである。
次に調べるなら、根基の全ての層 $\mathcal R^a/\mathcal R^{a+1}$ の
境界上の台と加群構造、または動く中心指標との貼合せを対象とする。

\section{正規スライスから回収する素イデアルと葉空間}

\subsection{CGG量子化のイデアルを元の多項式空間で扱う機構}

$N$ が局所冪零、$[N,X]=N$ のCGG積を考える。
$u\in\ker N$、$Xu=au$、$a\ne0$ で、$C[u^{-1}]=K[t]$、
$Nt=1,Xt=t$ がある場合、局所化の内部自己同型は
\[
 u^{-1}\star f\star u=(1+sN)^a f.
\]
両側イデアルに属する $f$ の有限な $N$ 軌道上で、
$((1+sN)^{an}f)_{n=0}^{d}$ から $(N^kf)_{k=0}^{d}$ を補間で回収する。
従って全両側イデアルは $N$ 安定であり、$t$ の最高係数を順に除けば
$I=Q[t]$、$Q\subset K$、$XQ\subset Q$ となる。
Poisson側でも $\{u,f\}=-auNf$ から同じ回収ができる。
さらに
\[
 (Q_1[t])\star(Q_2[t])=(Q_1Q_2)[t]
\]
なので、素性は微分素性へ移り、標数零の通常の素性と一致する。
この補題は、全局所冪零導分が線形であることを必要としない。
固有重量 $a=0$ のchartにはそのまま適用できない。

\subsection{正共鳴の全層と閉包}

$m=2j+r\ge3,r>0$、通常CY値、$s\ne0$ について、
$x$、商での $y$、必要な場合の商での $z$ を順に反転する。
スライスと不変環は
\[
\begin{array}{c|c|c}
 \text{部分}&t&K\\\hline
 x\ne0&y/x&\mathbb C[x^{\pm1},q,J]\\
 x=0,y\ne0&z/(2y)&\mathbb C[y^{\pm1},H_r]\\
 j=0,x=y=0,z\ne0&-w/z^r&\mathbb C[z^{\pm1}]
\end{array}
\]
であり、$H_r=w+(-1)^j y^{2j-1}z^{r+1}/(2r+2)$ である。
残る閉部分では $N=0$ となる。
この順序で、$\operatorname{Spec}S=\operatorname{PSpec}(C,N\wedge X)$ が
同じ多項式の部分空間として成立し、全素イデアルは完全素となる。
集合の対応だけでなく、全包含関係と位相も一致する。

整数重量 $(m+3)X$ と $N$ は二次元可解代数群 $G$ を与える。
階数零の集合は $F=V(x,y)$ ($j\ge1$)、$F=V(x,y,z)$ ($j=0$) であり、
他のシンプレクティック葉は二次元 $G$ 軌道である。
$F$ は $G$ の固定点集合ではない。$F$ の各点は一つの葉なので、葉の飽和は
\[
 \operatorname{Sat}_\pi(T)=G\cdot(T\setminus F)\cup(T\cap F)
\]
としなければならない。
葉閉包のイデアルを取る写像 $\kappa$ は商写像であり、任意の点集合 $T$ に対し
\[
 \overline{\kappa(T)}=\kappa(\overline{\operatorname{Sat}_\pi(T)})
\]
が成立する。原始スペクトルの全域的な位相はこの葉空間で決定した。
中心パラメータだけを使う閉集合の列挙と、部分解消上の代数層は別問題である。

原始スペクトルの各層の内部は、$n=(m+r+3)/\gcd(m+r+3,r)$ として
$A_{n-1}$ 曲面、$\mathbb A^1$、$j=0$ の場合だけ現れる一点、および $F$ である。
非文字表現の原始商は、該当する正規元を反転すると全て $\mathcal D(\mathbb G_m)$ となる。
特に $r\ge2$ で全中心が $\mathbb C$ であっても、原始スペクトルには二次元の開層がある。

\subsection{Dixmier--Moeglin同値に必要な最小の分数体情報}

局所素商の係数体 $K_0$ は整数重量の固有元で生成され、非零重量の可逆元 $u$ を持つ。
$u$ の重量を $a\ne0$、固有元 $f$ の重量を $b$ とすると、
$f^a/u^b$ は不変である。従って
\[
 \operatorname{trdeg}K_0^X=\operatorname{trdeg}K_0-1.
\]
この不変体が分数斜体の中心に入ることだけで、rational性から $\dim K/Q=1$ を得る。
分数斜体の全中心を先に計算する必要はない。
可逆固有元を持つ既約曲線では、全 $\mathbb G_m$ 軌道が開なので一軌道に限られ、
局所量子商は $\mathcal D(\mathbb G_m)$ となる。
これと正規局所化補題を合わせて、量子・Poisson双方のDixmier--Moeglin同値を証明できる。

\subsection{次に使える方向}

全域的な葉空間が得られたため、$r=1$ では $D$ 型中心の部分解消と
原始スペクトルの関係を、集合の対応より強い対象で記述することが次の候補となる。
中心の原点ファイバーの冪零厚みはZariski位相からは失われるので、
代数層を構成するなら冪零指数と両側加群の貼合せを同時に保持する必要がある。
$r=0$ や非CY値へ広げる場合は、ゼロ固有重量のchartを先に特定し、
今回のイデアル補間が成立する範囲を個別に判定する。

\section{固定量子VKの補完軌跡と再利用できる制限複体}

重み三のCGG共鳴 $(j,r)=(0,3),(1,1)$、$s\ne0$ に対し、
六項代数を $\Lambda$、$e=e_0+e_1+e_3+e_4$、$B=e\Lambda e$ とする。
固定した量子代数での結果は
\[
 HH^\bullet(B)=(1,1,35,39),\qquad
 R_B\simeq\mathbb C[[z_1,\ldots,z_{35}]],\qquad
 R_\Lambda\simeq\mathbb C[[\gamma]]/(\gamma^{j+1})
\]
である。適合的なhullでは制限を
\[
 \operatorname{Spf}R_\Lambda\hookrightarrow\operatorname{Spf}R_B
\]
と表示できる。$j=0$ では補完可能な一次方向は零であり、
$j=1$ では一次元だがその非零方向は二次で止まる。
可換基点の五十三次元の接空間同型とは、基点が異なる。

再利用すべき機構は、相対正規化複体の全射
\[
 C_E^\bullet(\Lambda_h)\longrightarrow C_{eE}^\bullet(B_h)
\]
とその核 $K_h^\bullet$ である。
可換基点での $HH^0,HH^1,HH^2$ の同型から $H^1(K_0)=H^2(K_0)=0$ が従う。
有限自由複体の階数特殊化によりgeneric fiberでも同じ消滅を得るため、
固定量子化での $HH^1$ 同型と $HH^2$ 単射が同時に得られる。
非零量子化係数の同型があれば、この結論を任意の $s\ne0$ へ移せる。

可換cornerの結合積スキームは次元 $53+598=651$ の滑らかな点を持つ。
量子化後もこの滑らかさが残り、軌道次元は
$620-1-3=616$ となるため、商のminiversal方向は三十五個である。
この計算は、障害群 $HH^3(B)$ が三十九次元残っていてもhullが滑らかになり得る
具体例として使える。

全 $\alpha$ に対して指定cornerを保つことは、単なる圏同値に頼らず確認できる。
\[
 a\star_{\alpha+\beta}b=T^{\beta\deg b}(a)\star_\alpha b,
 \qquad T=1+sN
\]
から、右加群の通常のEnd合成には $G_{ik}(a)=T^{-\beta i}(a)$ を使う。
経路順の積には $F_{ik}(a)=T^{\beta k}(a)$ を使う。
どちらも全冪等元を固定する。
Hochschild検算や既存quiver表示の規約を切り替える際、この向きの差を明示する。

補完像の差は、\emph{指定した外側と指定した埋め込み}を含む不変量である。
二つの $B$ 自体の導来非同値は、次節のGerstenhaber作用によって別に証明される。
今後、全ての最小補完を統一的に記述する関手を作る場合も、
この標準補完像と圏単体の不変量を区別する必要がある。

\subsection{四項VK圏の内在的区別と次の対象}

$m=3$ の二つの非零CGG量子化について、四項圏そのものの導来非同値が確定した。
外側の補完環の差を使う必要はない。
$HH^1$ は両方とも一次元なので、非零生成元 $\xi$ の作用
$[\xi,-]$ の階数と冪零Jordan分割が、生成元の選択によらない導来不変量になる。
$\alpha=0,s=1$ の代表で $\xi=[\log(1+N)]$ を使うと
\[
\begin{array}{c|c|c}
(j,r)&HH^2\text{ のJordan分割}&HH^3\text{ のJordan分割}\\\hline
(0,3)&(13,10,7,5)&(14,11,8,5,1)\\
(1,1)&(11,9,6,5,3,1)&(13,10,7,5,3,1)
\end{array}
\]
を得る。従って一次変形の固定方向は四次元と六次元である。
可換点も含めれば、この明示CGG族の四項VK圏にはちょうど三つの導来同値類がある。
どちらの非零点でも通常のhullは三十五次元の滑らかな環なので、
hullとHH次元が一致することは圏の同値を意味しない。

再利用できる計算法は、全cocycle空間の基底行列 $K$ を使う式
\[
 \operatorname{rank}(L:HH^2\to HH^2)
 =\operatorname{rank}[d^1\mid LK]-\operatorname{rank}d^1
\]
である。商座標による計算から独立に、第一階数だけで今回の非同値性を確認できる。
導分には $N$ でなく $\log(1+N)$ を使う。
補助次数 $(0,1,2,m+r+1)$ に対する下降作用が、有限複体上の冪零性を保証する。
反対代数への入力反転はこの作用を共役するので、End合成の規約によらない。

次に調べる価値があるのは、全ての $HH^2$ 方向の自己同型作用と、
標準補完方向がその中で占める位置の関係である。
一次の固定部分空間だけでなく、外的導分を持ち上げる条件の全完備局所環も決定した。
全次数の対称性保持族、純粋部分への還元、標準補完接線の二次消失を後述する。
一つの外部導分が与える作用は一般のVK変形へ移ると消滅し得るため、
共鳴点以外の層を分類するには、自己同型群とHochschild括弧の同時変化を追う必要がある。


全相対Hochschild dg Lie代数 $C_E^\bullet(B)[1]$ は、
どちらの非零成分でもhomotopy abelianではない。
$HH^1\times HH^2\to HH^2$ の非零括弧がその証拠であり、
異なる作用階数は二つのdg Lie代数がquasi-isomorphicでないことも示す。
これは三十五次元の滑らかなhullと両立する。
ここで全複体の $C_E^0$ は次数 $-1$ に残す。
$C_E^{\ge1}[1]$ への切捨てでは $H^0$ が $\operatorname{Der}_E(B)$ となり、$HH^1$ との混同が生じる。

\paragraph{元の射空間に見えるJordan鎖.}
$V_2=\langle x,y,z\rangle$ に $N_0=x\partial_y+2y\partial_z$ を作用させると、
通常の次数 $m$ の多項式空間はJordan鎖の直和
\[
 \mathbb C[x,y,z]_m
   =\bigoplus_{k=0}^{\lfloor m/2\rfloor}q^k V_{2(m-2k)}
\]
に分解する。ここで $V_{2r}=\langle N_0^a z^r:0\le a\le2r\rangle$ とする。
その鎖の長さは $2m-4k+1$ である。
$Nw=-q^jz^{m-2j}$ を加えると、$k=j$ の鎖だけが一つ長くなり、
長さ $2(m-2j)+1$ が $2(m-2j)+2$ へ変わる。
$\delta=N(1-N/2+\cdots)$ は $N$ と同じJordan分割を持つ。
従って $m=3$ の次数三射空間での分割は $(8,3)$ と $(7,4)$ となる。

これは普通の代数非同型の短い確認にも使える。
四頂点の到達関係は全順序なので、代数同型を内的同型で補正すれば各冪等元を固定できる。
そのとき $\operatorname{Der}_E(B)=\mathbb C\delta+\operatorname{ad}(E)$ であり、
$\delta$ の像は $c\delta+\operatorname{ad}(h)$（$c\ne0,h\in E$）である。
各射空間上では後項は $h_i-h_j$ というスカラーなので、
像が冪零であることから全差が零となり、像は $c\delta$ に限られる。
次数三射空間のJordan分割の差から代数非同型が従う。
この議論は三角同値には直接適用できないため、導来非同値には本文のHochschild作用を使う。
一般重みの四項cornerで同じ議論を内在化するには、まず $HH^1$ の一次元性が必要である。

\subsection{外的導分の跳躍を有限個の方程式へ縮約する}

量子VKの三十五変数hull $R$ 上では、相対Hochschild複体を単元消去すると
次数一から二への一列 $a:R\to R^{35}$ だけが非零の微分として残る。
次数零は $R$、次数三は $R^{39}$ である。
結合積スキームが六百五十一次元で滑らかなため、$d^2$ の階数は九百三十三で一定となる。
この点が、特殊ファイバーの階数だけから残余微分の消滅を推論しないために必要である。

$I=(a_i)=\operatorname{Fitt}_{34}H^2(C_R^\bullet)$ は、外的導分の持上げが存在する
閉形式部分 $Z$ を定める。列の一次項は $-L_\delta$、$\delta=\log(1+N)$ である。
従って $T_0Z=\ker L_\delta$、次元は $j=0$ で四、$j=1$ で六。
$Z$ の外では $HH=(1,0,34,39)$ となり、通常のhullは三十四変数で滑らかである。
$R$ 上の相対 $HH^1$ 自体は零なので、特殊化との非可換性を見落とさない。
跳躍は複体の導来基底変換で保持される。

独立な一次方程式を消すと、$Z$ の全環は
\[
 \mathbf C[[v_1,\ldots,v_k]]/(f_1,\ldots,f_k),\qquad f_i\in(v)^2,
 \qquad k=4\text{ または }6
\]
に縮約される。二次項と、後述の全次数の族を組み合わせて高次構造を限定できる。
$j=0$ では全て零、$j=1$ では一つの平方 $v_6^2$ が生成する。
この平方が消える超平面は、内在的には $\ker L_\delta\cap\operatorname{im}L_\delta$ である。
従って $j=1$ の外的導分跳躍部分は原点で特異で、次元は正確に五であり、滑らかな五次元の閉形式部分を持つ。
二次項が平方であることだけから全環の非被約性を推論しない。

再利用する二次障害の公式は、$u\in\ker L_\delta$ に対して
\[
 d^1\eta=L_\delta u,\qquad d^2v=u\circ u,
 \qquad
 \mathfrak q(u)=[L_\delta v+L_\eta u]\bmod\operatorname{im}L_\delta
\]
である。代数変形の無障害性と、導分を同時に持ち上げる条件は異なる。
独立な検算は $2556\times2204$ のblock行列
$\bigl(\begin{smallmatrix}d^2&0\\L_\delta&-d^1\end{smallmatrix}\bigr)$ で
二つの持上げ条件を同時に解く方法であり、第二次の積補正の選択に依存しない。

ここで保持するのは非零外的導分の存在であり、指定した $\mathbf G_a$ 作用や
局所冪零性を含む同変変形ではない。
全対称性保持部分は $j=0$ で滑らかな四次元、$j=1$ では滑らかな五次元部分の二重構造となる。
標準補完線と五次元接平面の位置関係も決定できる。
純粋三変数部分への還元と全二重構造の証明は後述する。

\subsection{原始射に支持を持つ変形による全次数の対称性保持}

四対象 $(0,1,3,4)$ で、原始射 $w_0:0\to3$ と次数一射の積だけを動かす
部分空間 $W$ は五十四次元、その出力を純粋な四次多項式に制限した $W_0$ は
四十五次元である。
三重合成 $0\to1\to3\to4$ の最初の積には $w_0$ が現れないため、
全 $F,G\in W$ について $d^2F=0$、$F\circ G=0$ である。
この支持の判定は、有限次数の検証を全次数の証明に変えるために再利用できる。

不変cocycle $U_i$ と導分補正 $E_i$ について $d^1E_i=[\delta,U_i]$ とする。
$E_i$ が $w$ を含む射を純粋多項式へ送るだけなら、
全 $V\in W_0$ に対して $[E_i,V]=0$ である。
$B_{0,3}$ の四本の不変方向は全て $W_0$ に入り、
$\mu+\sum u_iU_i$ と $\delta+\sum u_iE_i$ がそのまま全次数の族となる。
接空間次元と一致するため、対称性保持部分全体は滑らかな四変数形式空間である。

$B_{1,1}$ の $\ker L\cap\operatorname{im}L$ の五本では、
二次の導分誤差は $W_0$ に入り、階数三十八の
$L_\delta:W_0\to W_0$ の像に入る。
有理数上で $L_\delta V_{ij}=-([E_i,U_j]+[E_j,U_i])$ を解く
（対角項は一項だけ）と、積の二次補正だけで全次数の族を得る。
三次以上の導分条件と全結合則は支持から零である。
この族は滑らかな五次元閉形式部分を与え、対称性保持部分の次元が正確に五と決まる。
両族の導分は $w$ フィルトレーションを保ち、その補正は次数を下げるため、
一様に十二乗が零であり、実際の $\mathbf G_a$ 作用を持つ。

五次元の滑らかな部分を $t=0$ と直線化すれば、全跳躍環は
\[
 A[[t]]/(t(t-b),tJ)
   =A\oplus(A/J)t,\qquad
 A=\mathbf C[[u_1,\ldots,u_5]],\quad
 b\in\mathfrak m_A^2,\quad J\subset\mathfrak m_A^2
\]
の形になる。証明は、定義イデアルが $(t)$ に含まれることと、
二次項 $t^2$ を持つ関係を $t$ で割った元への形式的陰関数定理である。
この有限表示は、横断族を構成する前の候補を限定するために有用である。
v39では横断方向を含む全次数の有理式族により $b=J=0$ と決まった。
従って実際の環は $A[\varepsilon]/(\varepsilon^2)$ であり、$A$ 上自由階数二である。

\section{任意の重みで使える標準部分の還元}
\label{sec:notes-all-weight-filtered-v40}

\paragraph{固定する対象。}
$m=2j+r\ge3$, $j,r\ge0$、$\alpha=0,s=1$ のCGG座標代数の
$(0,1,m,m+1)$ 切出しを $B_{j,r}$、純粋三変数部分を $N_m$ とする。
$r=0$ を含む。$n_d=\binom{d+2}{2}$ と書く。
任意の標識付きArtin変形で $N_m$ は自由直和因子の部分代数として一意に持ち上がり、
二段・三段合成の像の階数は $n_m,n_m,n_{m+1}$ に保たれる。
同型と、冪等元を零にする正規化導分もこの部分を保つ。

\paragraph{行列計算から独立した消滅機構。}
可換積の正の $w$ 重みの二cocycleに対し、中間の対称積の $m-3$ 因子を固定する。
残る四変数への二重反対称化は重み三と同じで、係数 $n(n-2)$ により一次・二次が消滅する。
CGG積では補助重み $(0,1,2,m+r+1)$ を使う。$N$ は一つ下げ、$X_0$ は保つので、
有限フィルトレーションのassociated gradedは可換の複体である。
従って $Q=C/C_F$ に対し
\[
 H^1(Q)=H^2(Q)=0,\qquad \dim H^3(Q)=4m^2-6m-7.
\]
この二段階の証明は、非零量子点に個別の行基本変形を繰り返す必要をなくす。

\paragraph{相対安定化群を殺す明示的な汎関数。}
純粋部分を固定する正規化導分は、二つの原始射で共通に $w\mapsto aw+p$ と書ける。
$q,y$ との混合関係により
$Np=aq^jz^r$、$X_0p=(m+r+1)p$ である。
許される $p$ の基底は
$x^b y^{2j-1-2b}z^{b+r+1}$、$0\le b<j$。
$N$ の二重対角行列を左から零化する汎関数は
\[
 \ell_0=1,\qquad
 \ell_{b+1}=-\frac{2(b+r+1)}{2j-1-2b}\ell_b.
\]
$q^jz^r$ とのpairingは
$(-1)^j\sum_{b=0}^j|\ell_b|\binom jb\ne0$ だから $a=0$、続いて $p=0$ となる。
これは全共鳴に共通する有限次元の障害を、具体的な有理汎関数で表示している。

\paragraph{相対変形の次元と通常hull。}
純粋部分の変形とその同定を固定した相対変形は滑らかで
\[
 k_m=3n_{m+1}-2n_m+7=\frac{m^2+9m+28}{2}
\]
次元である。原始射を含む三重合成がないため相対結合性条件はなく、
上記の汎関数が相対自己同型を消す。
最初の値は $k_3=32,k_4=40,k_5=49,k_6=59,k_7=70$。
通常の変形接空間については
\[
 0\to HH^1(B_{j,r})\to HH^1(N_m)
 \xrightarrow{\partial}\mathbf C^{k_m}
 \to HH^2(B_{j,r})\to HH^2(N_m)\to0
\]
を使う。今回、$\partial=0$ も全ての重みで証明できた。
従って適合するminiversal環について
\[
 R_{B_{j,r}}\simeq R_{N_m}[[u_1,\ldots,u_{k_m}]]
\]
となる。$N_m$ と $k_m$ は $j,r$ によらないので、固定した重みの全共鳴は
端点も含め同じ通常変形環を持つ。対称性保持部分まで同じになるという主張ではない。

\paragraph{純粋部分の外的導分を一様に計算する方法。}
可換点の導分で三段積 $vtz=ztv$ を微分すると、左右の次数一作用の差はスカラーになる。
自然な $\mathfrak{gl}_3$ 作用と内的導分を引いた後、中間作用は
$xG(yt)=yG(xt)$ を満たすのでスカラーである。よって可換点の $HH^1$ は $\mathfrak{pgl}_3$。
三本の $x,x^{m-1},x$ を可逆にする普遍局所化は $M_4(\mathbf C[Y,Z])$ である。
$m\ge3$ が、$yzx^{m-3}$ から $YZ=ZY$ を導くために必要となる。
CGG族を $\mathbf C[h]/h^n$ 上で局所化すると、既知の平坦な座標chartへの全射は法 $h$ で同型。
核 $K$ は、平坦な商へのTor完全列から $K/hK=0$、従って $K=0$ になる。
元の局所化の平坦性を先に仮定しないことが重要である。
一次で持ち上がる導分は $2(Z-Y^2)\partial_Y\wedge\partial_Z$ を保つ。
射影的ベクトル場の中でその安定化群は $\partial_Y+2Y\partial_Z$ の一次元である。
頂点内的導分の像は三本の隣接射の自己係数により自由直和因子と分かる。
その商での飽和された導分核への特殊化と、実際の $h^{-1}\log(1+hN)$ により
$HH^1(N_m)=\mathbf C$ が全ての重みで従う。
その生成元は四変数側にも存在するため、上の接続写像 $\partial$ は零である。

\paragraph{次の比較問題。}
純粋部分の二次Hochschild群と全変形環を非可換射影平面のものと比較することが有望である。
独立な有理数計算では $m=3,4,5$ の $N_m$ について
\[
\begin{array}{c|c|c|c}
 m & (HH^1,HH^2)\text{、可換点}&(HH^1,HH^2)\text{、非零CGG点}&\dim HH^3\\\hline
 3&(8,10)&(1,3)&28\\
 4&(8,10)&(1,3)&108\\
 5&(8,10)&(1,3)&270
\end{array}
\]
となる。一次群の全重みの主張は上で証明済みだが、表の二次群と滑らかさを一般の $m$ へ外挿しない。
特に $(\mathcal O,\mathcal O(1),\mathcal O(m),\mathcal O(m+1))$ は
$\mathbb P^2$ 上では逆向きの高次Extを持つので、通常のtilting同値をそのまま使えない。
一般比較を証明できれば、未決の $R_{N_m}$ 自体と全cornerの通常hullを決定できる。
対称性保持曲線が原始射の接合によって何次に切り詰められるかは、その後にも残る別の問題である。

\paragraph{検証の範囲。}
通常ゲートは $3\le m\le7$ の全十六共鳴と、五つの可換正重み複体、
七十七の極端重量の汎関数、および射影的Poisson安定化群の階数七を厳密に確認する。
全整数の証明は上記の構造による。純粋部分の表は別の任意実行プローブで再現でき、
一般定理の証拠へ読み替えない。

\section{全形式的な純粋部・原始射分解}
\label{sec:vk-filtered-reduction-notes-v39}

今回の対称性保持環の決定と独立に、全変形を小さい問題へ還元する構造が得られた。
今後の分類や別の量子VKに適用するため、必要な情報を残す。
ここで扱う形式近傍の基点は $s=1$ の非零量子点 $B_{0,3},B_{1,1}$ である。
可換VKの $53$ 次元形式近傍やLO型の別の基点については、
本節の商複体の消滅を別に調べる必要がある。

\paragraph{フィルトレーションは追加の仮定ではない。}
$w$ の個数による二段のフィルトレーションを保存するHochschild複体を $C_F$ とする。
両 $B_{0,3},B_{1,1}$ に対して商 $C/C_F$ は
\[
 (\dim Q^1,\dim Q^2,\dim Q^3)=(65,216,162),\qquad
 (\operatorname{rank}d^1,\operatorname{rank}d^2)=(65,151)
\]
を持つ。$H^1(Q)=H^2(Q)=0$ により、全ての標識付き形式変形を次数別の基底変換で
保存的にでき、保存的な変形同士の同型も保存的である。
頂点冪等元を零にする正規化導分にも同じ次数別議論が適用される。
任意の外的導分類は内的導分を差し引いてそのように正規化できる。
従って、原始射に関する支持やフィルトレーションを用いる探索は、
全変形問題の一部だけを見る制限ではなく、正当な形式正規化として用いられる。
さらに純粋部分は、二段合成の像 $N_{03},N_{14}$ と三段合成の像 $N_{04}$ から
復元される。基点の階数 $10,10,15$ は保存的表示とNakayamaの補題により
全ての形式変形で不変である。従ってこのフィルトレーションは合成に内在する。

\paragraph{三方向と三十二方向。}
純粋な $x,y,z$ の四対象代数 $N$ は
$HH^*(N)=(1,1,3,28)$ を持ち、その変形hullは滑らかな三変数環である。
相対複体の次元は $(76,108,0)$、最初の微分は単射であり、相対変形は滑らかな
$32$ 次元である。左原始合成を固定すれば、右原始合成の $54$ 係数から
二つの原始射のスカラー変更と純粋三次式の追加、合計 $22$ 係数を除けばよい。
純粋出力に限定した $45$ 次元部分のHochschild像は $24$ 次元である。
従って原始合成だけを用いて全 $35$ 方向を求めることはできない。

\paragraph{残る三方向の対称性。}
$L_\delta$ は $HH^2(N)$ で単一の長さ三のJordan鎖を持つ。
純粋部の対称性保持部分全体は滑らかな $\gamma$ 曲線であり、
公表稿の関係式と
\[
 D_\gamma(x)=0,\qquad D_\gamma(y)=x+2\gamma Z,\qquad
 D_\gamma(Z)=2y
\]
がその全次数の模型になる。
次数一の作用の特性多項式は $z(z^2-4\gamma)$ である。
これは純粋部の外的導分保持の曲線であり、一般の $\gamma$ で冪零導分を要求する
曲線とは区別する。

\paragraph{有限行列問題への具体的な書換え。}
以下は $B_{1,1}$ の場合である。
純粋部とその導分を $(C_\gamma,D_\gamma)$ に固定し、左原始合成を基底として
\[
 w_0v=M(v)w_1+H(v),\qquad
 M\in\operatorname{End}(V),\quad H\in\operatorname{Hom}(V,C_{\gamma,4})
\]
と書く。導分の原始射上の値を
$\delta(w_i)=\lambda_iw_i+P_i$、$P_i\in C_{\gamma,3}$ とすると、
導分条件の $w$ 成分は
\[
 [D_1,M]+(\lambda_1-\lambda_0)M=0.
\]
$M$ は基点で可逆である。$M^{-1}$ を掛けてトレースを取れば
$\lambda_0=\lambda_1=\lambda$ となる。
$D_1=D_\gamma|_V$ は $Z$ を巡回ベクトルとするので、その中心化環は
基底環上でも自由な $\langle I,D_1,D_1^2\rangle$ である。
従って
\[
 M=aI+bD_1+cD_1^2,
\]
純粋成分の導分条件は
\[
 D_4H(v)-H(D_1v)-\lambda H(v)+M(v)P_1-P_0v=0
\]
だけになる。基点では $M=I-2D_1+2D_1^2$ である。
これは今後、正規化された対称性保持族を小さい線形代数で探索するための式である。
式だけから全基底上の可解性を推論せず、今回の全次数解は別の厳密有理式検証により保証する。

\paragraph{一次と二次の障害の所在。}
同変完全列
$0\to V_{j,r}\to HH^2(B_{j,r})\to HH^2(N)\to0$ において、
不変部分から純粋部への写像の階数は $j$ である。
$j=0$ では純粋 $\gamma$ 方向に対する相対接合障害が一次で現れる。
$j=1$ では一次で持ち上がるが、既出の二次障害は純粋部へ写すと零になり、
原始合成の相対余核の非零クラスとして残る。
この違いは二つの対称性保持部分の形を説明する際に利用できる。

\paragraph{再現用資料。}
\texttt{verify\_vk\_filtered\_reduction\_v39.py} は通常実行で読み取り専用であり、
商複体、相対複体、完全列と $L_\delta$ の同変性、純粋PBW重なり、
全三関係の導分条件を再計算する。
\texttt{vk\_filtered\_reduction\_v39\_certificate.json} は
有理数上の制限写像、原始合成の像、純粋接方向、および障害の制限を保存する。

\subsection{対称性跳躍の全環と全次数の検証法}

$B_{1,1}$ の対称性保持環は完全に決まり、
\[
 R_Z=\mathbf C[[u_1,\ldots,u_5,\varepsilon]]/(\varepsilon^2)
\]
である。従って前版の $A[[t]]/(t(t-b),tJ)$ という表示の未知量は
$b=J=0$ となる。被約部分は前版の五変数 $\mathbf G_a$ 族である。
全体は次元五の局所完全交叉、Gorenstein、重複度二で、埋込成分を持たない。

再利用すべき判定は次のものである。
埋込次元 $k+1$ の完備局所環 $R$ に非零な二次関係があり、
$\mathbf C[[u_1,\ldots,u_k,\varepsilon]]/(\varepsilon^2)$ 上の族へ
接空間同型を与える写像があれば、$R$ 自体がこの平方零環に同型となる。
最小表示からの合成を形式的逆関数定理で標準化すると、
定義イデアルは $(\varepsilon^2)$ に含まれる。
非零な二次関係は $\varepsilon^2$ の単元倍なので逆包含も成立する。
有限jetだけの一致からこの結論を推論しない。

今回の族は、既存の五変数族を $\mu_u,D_u$ とし、
\[
 \widetilde\mu=\mu_u+\varepsilon\Phi/\Delta,
 \mathcal D=D_u+\varepsilon\Psi/\Delta,
 \quad\varepsilon^2=0,\quad\Delta(0)=1
\]
とする有理式で構成した。
原点で $\Phi=U_6$、$\Psi=E_6$ となる。
二cochainの追加補正は原始射の右合成に支持を持つ二十二座標、
一次cochainの追加補正は百座標だけである。
分母は全次数十二、分子は十四以下であり、
元のCGG積から直接再構成した九百七十二個の結合則と
千五百八十四個の導分方程式が、分母除去後に多項式恒等式として全て零となる。
この直接検証は、構成時に選んだ百二十二行の線形方程式から独立である。
全有理係数は \texttt{vk\_double\_symmetry\_v39\_witness.json}、
検証コードは \texttt{verify\_vk\_double\_symmetry\_v39.py} に保持する。

$D_u^{12}=0$ と $\varepsilon^2=0$ から、補正 $\Psi$ の形によらず
$\mathcal D^{24}=0$ が従う。
一般に $D^n=0$ かつ $\varepsilon^2=0$ なら
$(D+\varepsilon E)^{2n}=0$ である。
一個の $\varepsilon$ を含む項は $D^aED^b$、$a+b=2n-1$ だからである。
したがって全非被約基底上で有限和の指数写像による $\mathbf G_a$ 作用を得る。

通常の三十五変数hullの形式座標と残余列の基底も選び直せる。
\[
 R_B=\mathbf C[[u_1,\ldots,u_5,\varepsilon,z_1,\ldots,z_{29}]],
 \quad I_B=(z_1,\ldots,z_{29},\varepsilon^2),
 \quad a=(z_1,\ldots,z_{29},\varepsilon^2,0^5)^t.
\]
従って任意の形式弧に沿う対称性消失次数は正確に
\[
 \nu=\min\{\operatorname{ord}z_1,\ldots,
                  \operatorname{ord}z_{29},2\operatorname{ord}\varepsilon\}
\]
であり、有限なら相対 $HH^2$ の捩れは $\mathcal O/(s^\nu)$ である。
無限大の弧は被約部分 $Y$ を通る。
また $\operatorname{Ext}_{R_B}^1(HH^2_{R_B},R_B)=R_B/I_B$ によって、
二重構造を相対Hochschild加群の双対から回復できる。

純粋三変数部分の対称性保持曲線を $\mathbf C[[\gamma]]$ とすると、
上の同型はこの環上で
$R_Z=\mathbf C[[u_1,\ldots,u_5,\gamma]]/(\gamma^2)$ と取れる。
理由は、被約五変数族が純粋部分を固定するので $\gamma=\varepsilon h(u)$ となり、
純粋部分への接写像が非零であるため $h$ が単元になることである。
$B_{0,3}$ の四変数対称性保持部分は純粋曲線の原点へ全次数で写る。
純粋部分では滑らかな曲線が存在するが、原始射の接合がその方向を
原点または平方零近傍へ切り詰める、という形で二つの共鳴の違いを説明できる。

標準六項補完の像も同時に直線化できる。
純粋部分への制限行列に補完接線を代入すると
$(229/120,117/20,-1571/60)^t$ となり、純粋 $\gamma$ 曲線の接線と一致する。
六項hullは $\mathbb C[[\gamma]]/(\gamma^2)$ なので、
この環への像は $u_i=a_i\gamma$ となる。
$u_i$ を $u_i-a_i\gamma$ へ変えれば、標準補完像のイデアルは
$(u_1,\ldots,u_5,z_1,\ldots,z_{29},\gamma^2)$ になる。
被約五次元部分は $\gamma=z_1=\cdots=z_{29}=0$ のままである。
従って標準補完像は、全対称性保持部分の横断的な長さ二の切片である。
この同時正規形は座標の選択に依存し、別の例外列を許す補完の分類ではない。

次に調べる価値があるのは、この平方零の厚みを一般の重み・共鳴へ拡張することと、
純粋三変数部分の対称性保持曲線を原始射の接合がどの次数で切り詰めるかという問題である。
今回の厚み二は全次数で確定したものであり、既存の外側の
$\mathbf C[[\gamma]]/(\gamma^{j+1})$ と同じ次数が一般に続くかは未決である。

\section{標準補完接線と対称性障害の比較資料}
\label{sec:completion-notes-v38}

\subsection{継続に必要な確定事項}

$m=3$, $j=r=1$ の量子VK corner $B$ と標準六項代数 $\Lambda$ について、
補完接線 $\ell=\operatorname{im}(HH^2(\Lambda)\to HH^2(B))$ は
\[
 \ell\subset K=\ker[\delta,-],\qquad
 \ell\cap\operatorname{im}[\delta,-]=0,
 \qquad \delta=\log(1+N)
\]
を満たす。従って、v37で残っていた「唯一の補完接線は、五次元の
$K\cap\operatorname{im}[\delta,-]$ に入るか」という問題の答えは否定である。
一次方向を失わずに、二次で外的導分を失う。

非零の補完初速度を持つ任意の形式弧 $\mathcal O=\mathbf C[[t]]$ に沿って、
対称性保持のFittingイデアルは $(t^2)$ へ引き戻される。
従って最小相対複体は
\[
 \mathcal O[0]\oplus
 (\mathcal O\xrightarrow{t^2}\mathcal O)_{1,2}
 \oplus\mathcal O^{34}[-2]\oplus\mathcal O^{39}[-3]
\]
となる。相対 $HH^2$ の捩れは正確に $\mathcal O/(t^2)$ であり、
一般ファイバーの $HH^1$ は零となる。
この結論は、二次以上の積補正や、滑らかな三十五次元hull内の弧の選択によらない。

初速度だけで弧を分類すると、$B_{1,1}$ では対称性消失次数は
$u\notin K$ で一、$u\in K\setminus\operatorname{im}L$ で二に固定される。
$u\in K\cap\operatorname{im}L$ では、任意の整数 $n\ge2$ と無限大が実現する。
後者は、滑らかな五次元部分 $Y$ 上の弧 $c(t)$ に
$t^n v$（$v\notin K$）を加えることで有限次数 $n$、$c(t)$ 自身で無限大を得る。
従って「二次対称性障害が零」という条件は、二次補正を適切に選べることを意味し、
任意の二次補正が対称性を保つことを意味しない。
$B_{0,3}$ の四次元の不変接空間についても、同じ次数の全範囲が実現する。

\subsection{再使用できるOre構成}

本計算では $\alpha=-1$, $s=1$ の表示で一次の外側変形を作り、
有限Zhang同型で v37 の $\alpha=0$ 表示へ戻した。
三変数代数は、既存の $C_{\gamma,s}$ の $s=1$, $\gamma^2=0$ への特殊化である。
$ux=x\sigma(u)$ を満たす正規元 $x$ を用いて
\[
 \mathcal S=C[w;\sigma^2,D_H],\qquad
 D_H(u)=x^{-2}(uH-Hu)
\]
とする。必要なのは $[C,H]\subset x^2C$ だけであり、
このとき $x^2w+H$ が中心元となる。
この補題は任意の正整数 $k$ に対し、
\[
 [C,H]\subset x^kC,\qquad
 D_H=x^{-k}[-,H],\qquad
 C[w;\sigma^k,D_H],\qquad J=x^kw+H
\]
へそのまま一般化できる。
$r>0$ の他の共鳴で $k=r+1$ を使う場合に再利用できる。
端点 $r=0$ の構成は $k=1$ に当たり、商 $C/(x)$ の可換性だけで
任意の $H$ を許せた理由もこの形で統一される。

今回の $H_0$ はCGG中心元
\[
 J=x^2w+q\left(xyz-\frac23y^3\right)
\]
のPBW表示から得た。一次補正 $H_1$ は、次数五のPBW元を未知数として
$[x,H],[y,H],[Z,H]$ の $x$ 次数零・一の係数を消す線形方程式を解いて得た。
未知数二十一、制約行三十九、係数行列と拡大行列の階数は共に十である。
自由変数を零にした解は公表稿に明示した五項の $H_1$ である。
次数三の $G$ に $H=QG$ という因子分解を課すと、
対応する行列の階数は六、拡大行列は七となり一次ですでに解がなくなる。
従って、$H_\varepsilon$ が二次の中心元 $Q_\varepsilon$ の倍数であるという
余分な制約を課すと、存在する補完接線を誤って排除する。
この計算は構成法の選択に関わるため保持する。

\subsection{計算規約と保存データ}

経路順の積を用い、cornerの対象は $(0,1,3,4)$ とする。
$\alpha=0$ から $\alpha=-1$ への有限同型は、$i\to k$ の射上で
$F_{ik}=T^{-k}$ である。
さらに、PBW元 $Z$ を $z+y$ へ送る固定した基底変換を用いる。
生成した一次cochain $U$ は $1584$ 成分を持つ。
v37の核基底における第六座標は $-1571/60$ であり、
余核の二次障害の第二座標は
\[
 \frac{432000}{2468041}\left(-\frac{1571}{60}\right)^2=120.
\]
この一つの非零係数だけでも、補完接線が二次対称性障害を持つことが分かる。

\texttt{verify\_vk\_completion\_symmetry\_v38.py} は通常実行時に書込みを行わず、
PBWの重なり、正規性、自己同型、全六対象における特殊ファイバーの積と一次結合律、
cornerのcocycle条件、核・余核座標、二種類の二次障害計算を行う。
証明書 \texttt{vk\_completion\_symmetry\_v38\_certificate.json} は
$U$、$d\eta=L_\delta U$ の解、$dV=U\circ U$ の解、
$L_\delta V+L_\eta U$ のコホモロジー座標を保持する。
比較先のv37証明書のSHA--256も検査する。

\subsection{全局所環の判定と残る障害の比較}

補完接線上の二次障害が非零であることだけから、全対称性保持部分の
完備局所環や非被約性を推論してはならない。
今回の全環の決定には、五変数の厳密な $\mathbf G_a$ 族に
平方零横断方向を加え、全ての結合則と導分条件を多項式恒等式として確認する議論が必要であった。
その結果、$A=\mathbf C[[u_1,\ldots,u_5]]$ に対し全環は
$A[[\gamma]]/(\gamma^2)$ と決まり、以前の有限表示における $b,J$ は共に零となる。
純粋部分への制限と適合する座標では、標準六項補完像は $u_1=\cdots=u_5=0$、
被約部分 $Y$ は $\gamma=0$ である。
この埋込みを含む同時正規形までが確定した範囲である。

一方、外側の結合性障害は $HH^3(\Lambda)$、内側の対称性障害は
$\operatorname{coker}(L|HH^2(B))$ に値を持つ。
両者の自然な比較は、後続の導分付き障害複体によって解決した。
外側の結合性障害に導分の一次補正を加えた類は、全充満部分圏への制限で自然に写る。
内側の結合性誤差を消すと、対称性作用の余核に一意の類が残る。
標準六項では一次元の一次障害線上に自然な非零線形写像を得る。
七次元の $HH^3$ 全体への比較写像や、一般重みの高次比較はこの主張に含めない。

\section{導分付き障害による補完障害の比較}
\label{sec:notes-obstruction-transgression-v40}

再利用する符号規約は $d=[\mu,-]$、$L=[\delta,-]$ である。
通常のHochschild微分とは $C^n$ 上で $(-1)^{n-1}$ だけ異なる。
導分付きの障害複体は
\[
 \mathcal P^n=C^n\oplus C^{n-1},\qquad
 \partial(f,g)=(df,Lf-dg)
\]
であり、自然な完全列
\[
 0\to\operatorname{coker}(L|HH^2)\to H^3(\mathcal P)
   \to\ker(L|HH^3)\to0
\]
を持つ。一般のAssDer複体と変形障害については
Das--Mandal\cite{DasMandal2026}を参照する。

一次対 $da=0$, $d\eta=La$ の障害は
\[
 (O,Q)=\bigl(\tfrac12[a,a],[\eta,a]\bigr),\qquad dQ=LO.
\]
$HH^1=\mathbf C\delta$ ならば、そのcone類は一次類だけから定まる。
$\eta$ を $c\delta+db$ だけ変更したときの第二成分の差は
$d(c\eta+[b,a])$ である。
$a'=a+df$, $\eta'=\eta+Lf$ としたときは
\[
 (O',Q')-(O,Q)=\partial(X,Y),\quad
 X=-[a,f]+\tfrac12[f,df],\quad
 Y=-[\eta,f]+\tfrac12[f,Lf].
\]
これらの式は基底選択からの独立性を直接確認するときに使える。

全充満部分圏への相対Hochschild制限 $\rho$ は微分・括弧・$L$ と可換である。
$\rho O$ を $dv=-\rho O$ により消せるとき、
\[
 q=L_Bv+[\rho\eta,\rho a]
\]
の $\operatorname{coker}(L_B|HH^2(B))$ における類は一意であり、
増強された外側の障害の制限そのものである。
固定した積補正の導分障害は $HH^2$ に値を取り、積補正を選び直す自由を
許した障害がこの余核に値を取る。
同じ議論は任意の小拡大に通用するが、出発点として導分付き変形が必要である。

$(\Lambda,B)=(A_{1,1},B_{1,1})$ では $HH^2(\Lambda)$ が一次元で、
$L_\Lambda$ は冪零なので零である。
従って一次導分補正が存在し、外側の非零平方障害線
\[
 O_{\rm std}=\operatorname{im}(\operatorname{Sym}^2HH^2(\Lambda)\to HH^3(\Lambda))
\]
から内側の余核への標準的な非零線形写像を得る。
既存の接線規格化での像は
\[
 (0,120,1640,-5917/12,-2269/2,42347/6)^t
\]
である。
この比較は七次元の $HH^3(\Lambda)$ 全体への写像ではない。
外側の一般重み $\gamma^{j+1}$ と内側の高次対称性障害を比較する際には、
各段階の増強障害を追跡する。
普通の結合性障害の消滅だけから、次の段階の導分持上げを結論しない。


\section{公表用原稿の構成と確定範囲}

十一本の原稿の境界は、主要な定理と読者が必要とする背景に合わせた。
比較と Hochschild 消滅を読むために、数百頁の局所 jet 計算を読む必要はない。
一方、点配置と相対微分は同じ点被覆から生じるため、定義の整合性を保持する。
\begin{enumerate}[label=\textbf{\Roman*.}]
\item \textbf{標準部分代数と変形障害}：正の重みの $HH^1,HH^2$ の消滅、任意 Artin 基底上の
標準部分代数の一意持上げ、非被約基底での合成階数、Loginov--Orlov 代数との非同値性・特殊化障害、
量子中間点と共鳴する節点族を一つの論証にする。
\item \textbf{二十一点配置と復元}：左関係の Hilbert--Burch 表示と右核の枠を用いて、
平面内の点配置から関係代数を復元する。$S_{21}$-monodromy を全外部自己同型群の消滅へ使い、
左右の標準対応が射影変換でないことから反対代数の対合の非自明性を証明する。
反対合を固定する三つの滑らかな族では、同型を除いた次元を七、十一、八と求める。
三族全てで被約二十一点配置が実在し、固定点数と行列式曲線の次数を署名による一つの式で説明する。
最大の自己反対族では、重み付き五直線の標準形から全monodromyと内在的復元を得る。
残る二族のPrym幾何の詳細は原稿V・VIIへ分け、共通の点配置の基盤を本稿に置く。
\item \textbf{相対微分と corank 二}：普遍点対象のSerre固定性から $\psi=0$ を全次数で証明する。
四枚被覆を四変数の最小versal族と九つの自明な方向へ分け、
$S_4$-monodromy、分岐自由因子とその正規化を記述する。
四枚被覆全体の非critical性と三変数sliceのcritical性を分類し、
critical族のcorank二からも同じ非critical被覆が生じることを明示例で証明する。
元の十一点族の相対potentialは高さ四以上の循環torsionイデアルの比較として保持する。
\item \textbf{完備 CY 代数と共鳴}：generalized down-up/GWA の明示形式族で、全 potential、
全中心、PI 性の必要十分条件を同時に記述する。任意に高い同次 jet を持つ非 PI の連続 $3$-CY 例を得る。
\item \textbf{四次Prymデータに付随する線織面と種数二十K3の境界周期}：四次の支配性とHermite曲線、
固定Gramの三成分閉包とK3平滑化を扱う。行列式線束による二つの線織面と二乗商、
端成分の半豊富性、偏極付き極限格子の根系・貼合せ・monodromyを決定する。
さらに全 $T^1$ 切断の実現、次数九のperiod isogeny、境界支配性とHalphen層を扱う。
\item \textbf{種数三Prymモジュライの相対有理性と二次被覆}：一般netの位数十八の安定化群を含めて
$\mathcal R_3\dashrightarrow Y_0(2)$ の一般ファイバーの有理性を証明する。
一般関係ファイバーのBrauer障害、周期・指数三と二の積による指数六、
Pl\"ucker二次曲線が与える内在的な二次分裂体とノルムトーラスを扱う。
逆Gramの四十五次元支配性から第二の楕円moduliを得る。
三次pencilでKummer写像を明示し、第三の二次被覆のAlbaneseと二次同種写像を決定する。
\item \textbf{六次Prym曲線の標準七次被覆と分岐曲線}：抽象Prym曲線からの復元、
十九変数・九余因子方程式、標準十八点の自由分解をまとめる。
五次netの有限平坦七次被覆と $S_7$ を同じ選択成分内で証明し、
分岐源の滑らかさ、種数三十七、分岐曲線の次数二十四と双対次数四十八を決定する。
\item \textbf{半直交成分の変形と重み付き射影空間への補完}：Hochschild制限を三次まで決定し、
六項側の全miniversal基底を三十八個の二次式で記述する。
全Artin基底上の節点型の補完条件と任意に遅い障害を証明し、原稿IXの幾何から非被約性・非純次元性を導く。
\emph{標準六項への補完}という量化を保ち、別の例外列の埋め込みは未決として分ける。
\item \textbf{重み付き射影三次元空間のPoisson成分と非被約構造}：全重み $m\ge3$ に対する
一般点の長さ $m+3$ と交叉の厚み、標準例外列代数の二次miniversal基底、
引戻し族の明示的4-Calabi--Yau量子化を主結果とする。
全重みの接線成分系列、一般点長さ、非線形CGG族のNakayama公式とDu Val中心も扱う。
重み三では十成分、三角形の局所環、留数による混合核、$c=0$ の被約分類をまとめる。
全成分の分類、最大次元、任意の量子化のAS正則性は未証明として分ける。
\item \textbf{重み付き量子射影圏の変形環と対称性}：全ての共鳴における固定非零量子化のhull、
反標準Hochschild--Serre環、族内の完全な同値分類を扱う。
偶数端点の全変形のAS正則実現、引戻し型の非完全交叉hull、量子VK四項圏の導来非同値を含む。
全形式変形の純粋部分への還元、対称性保持部分の完全な二重構造、任意弧の消失次数を決定する。
端点族では関数斜体の完全分類と非有界PI次数の集積も扱う。
全重みの純粋部分への還元と通常hullの分解、導分付き障害複体による標準補完の自然な障害比較を含む。
\item \textbf{重み付きCalabi--Yau量子化の中心と表現}：中心の部分解消と特殊化、
全primitive idealの分類、全域の葉の閉包位相、根基とDixmier--Moeglin同値を扱う。
端点の一の冪根chartでは、中心・Azumaya部分・特異中心ファイバーとMcKay模型を決定する。
さらに一変数混合項の全中心とその無標識な区別、最深点の完備局所Brauer障害、一般境界の巡回被覆と例外留数を扱う。
射影圏の同値分類から独立した、次数を忘れた通常CY座標環の表現論としてまとめる。
\end{enumerate}

先行する定理の既知性については、各稿の序論で一般理論との関係を記した。
特に原稿IVの代数を、新しい種類の generalized down-up algebra と呼ぶべきではない。
公表価値は明示形式族の完全判定と、その研究上の帰結として評価する。

\section{標準部分代数の持上げを他の問題へ使う}

$B_n=(\Sym(V)[w])_{(0,1,3,4)}$、$\dim V=n\ge3$ とし、$A^{(0)}$ を $w$ を含まない部分とする。
四頂点 $E=\kk^4$ を固定する。
\[
 HH^1(B_n)_{w>0}=HH^2(B_n)_{w>0}=0,
 \qquad C^0_E(B_n)_{w>0}=0.
\]
二番目の消滅は積の規格化を与え、一番目と零次 cochain の消滅は比較同型の一意性を与える。
任意の Artin 基底上で、$A^{(0)}$ を持ち上げる自由直和因子の部分代数は一意である。
この組合せが、単なる「座標変換で存在する」という主張より強い。

\subsection{一般化のための最小条件}
同じ帰納法を他の有向代数へ使うには、次を確認すればよい。
\begin{enumerate}
\item 積・微分・Gerstenhaber bracket と両立する重み分解がある。
\item 非正の cochain が bracket で閉じ、正の二次 cohomology が消える。
\item 正の一次 cocycle が零である。
\end{enumerate}
最後は $HH^1_{>0}=0$ だけでは一般には足りない。正の零次 cochain が存在する場合には、
内部導分が正の一次 cocycle を生じ、部分代数は内部共役までしか一意でない可能性がある。
また、持ち上げた頂点冪等元を動かす内部自己同型まで許すと、同じ部分集合としての一意性は失われ得る。
この二点を落とさずに、別の切り出しや exceptional collection へ応用する。

\subsection{比較研究で使える結論}
$n=3$ では二段の合成の像は任意 Artin 基底上で自由階数十の直和因子であり、
十一行十一列 minor は冪零成分まで含めて零である。
Loginov--Orlov の標準有限代数では階数十一なので、直接の平坦特殊化は不可能である。
一方、中間 Ore 代数を通る節点族は存在する。
これを、標準生成対象を忘れた圏の変形の不可能性へ拡張するには、新しい議論が必要である。

\section{点配置モジュライが与える次の対象}

$V,W$ は三次元、$U$ は六次元とする。
左関係 $E\subset V\otimes U$ は七平面で、$6\times7$ 一次式行列 $L$ を与える。
最大小行列式の零点が有限なら、長さは二十一であり、
\[
 0\longrightarrow R(-7)^6\xrightarrow{L^t}R(-6)^7
 \longrightarrow I_Z\longrightarrow0,\qquad R=\Sym V
\]
が最小分解になる。
従って $Z\subset\PP^2$ は $L$ を $\PGL_6$ まで復元する。
右関係は、
\[
 F_E:U\otimes E\longrightarrow V\otimes\bigwedge^2U,
 \quad Q\longmapsto(L_iQ^t-QL_i^t)_i
\]
の三次元核の射影枠から復元される。

第二成分の一般のモジュライは、$\Hilb^{21}(\PP^2)$ の開部分における
$\rk F_E\le39$ の二十四次元成分 $\mathcal H_C$ を $\PGL_3$ で割ったものとして記述できる。
商の次元は十六である。
この表示は $\mathcal H_C$ の有理性を含まず、rank locus 全体の既約性も含まない。

\subsection{一般の剛性と uniform position}
一般の点配置の Galois 群は $S_{21}$ である。
配置の射影自己同型群は $S_{21}$ の正規部分群となるが、射影枠を固定する三-cycle は
非自明な射影変換から来ない。そのため一般の射影安定化群は自明であり、
復元定理から $\Out(A_E)=1$ が従う。
これは $HH^1(A_E)=0$ だけでは排除できない有限外部自己同型も排除する。

また、任意の部分集合の次数 $d$ の評価写像は最大階数である。
特に三点の共線性、六点の二次曲線上への集中、十点の三次曲線上への集中、
十五点の四次曲線上への集中は、一般点で起こらない。
この情報はコンパクト化の境界を識別する際に使える。

\subsection{反対代数の非自明な対合}
有限代数を反対代数へ移し頂点順を反転する操作は、第二成分の十六次元商に
双有理対合 $\iota$ を与える。この対合は非自明であり、一般には
\[
 A_E\not\simeq A_E^{\mathrm{op}}
\]
である。従って商の不変体に対して
$[k(\mathcal M):k(\mathcal M)^{\langle\iota\rangle}]=2$ となる。

証明には三つの入力を使う。一般の $\Out(A_E)=1$ により、もし左右配置が射影同値なら
その射影同型は一意で関係空間の関数体 $K=\mathbb C(Y_C)$ へ降下する。
thin moduleから定まる左右の標準対応との差は $S_{21}$ の中心化群に入り、自明である。
一方、明示点 $A_*$ の十個の被約点で標準対応 $x\mapsto z$ を射影変換 $G$ に要求すると、
$z\wedge Gx=0$ の有理線形系は階数九となる。
九次小行列式が四となる証明書があり、非零の $G$ は存在しない。
これは $A_*$ を含む有限平坦なthin-module族上のベクトル束写像の階数である。
階数九は開条件なので、一般での射影同値性と矛盾する。

この結果は一般の有限代数の同型に関するものである。
反対代数との導来同値や、対合の全固定点の分類は未決定である。

\subsection{自己反対部分族の次元}
$J_0=\operatorname{diag}(-1,1,1,-1,1,-1)$、
$D=\operatorname{diag}(1,-1,1)$、$P=(12)$ とする。
中央矢を $U=\Sym^2V$ と同定し、反対合の三代表を
$(J_0,I)$、$(J_0\Sym^2D,D)$、$(\Sym^2P,P)$ と取る。
境界点 $B_-$ を通る固定成分 $Z_a$ の確定した次元は次の通りである。
\[
\begin{array}{c|c|c|c|c|c}
 a&\text{中央行列の固有次元}&\dim Z_a&
 \dim\overline G^{\tau_a}&\dim\overline{\overline GZ_a}&
 \text{同型類の次元}\\\hline
 1&(3,3)&32&25&58&7\\
 2&(1,5)&44&33&62&11\\
 3&(4,2)&35&27&59&8
\end{array}
\]
$Y_C$ は $B_-$ の近くで、滑らかな $D_0$ 上の射影的核枠束と実際に同型なので滑らかであり、
標数零の対合の固定スキームもその近くで滑らかである。
第三の場合には、全関係スキームの固定接次元三十六から、
$Y_C$ に横断する正固有値の法線一方向を除いて三十五を得る。
各族には独立な二つのトーラス重みを持つ不変接方向があり、
Hochschild微分のSchur行列の二次小行列式の先頭係数は $\pm4$ となる。
従って一般に $HH^1=0$ であり、軌道の固定部分の次元を引いて上表を得る。

三つの飽和の包含関係と全固定部分の分類は未決である。
三族全てと被約二十一点の開部分との交わりは非空である。
以下のHensel構成は各族について同じ基点を通る固定成分上で行う。
飽和の閉包に加わる全境界点まで可逆な反対合が延びるとは限らず、
$HH^1=0$ だけでは有限な全外部自己同型を排除できない。
$Z_2$ の一般点については、後述の全monodromyから別に $\Out=1$ が従う。
$Z_1,Z_3$ の一般の全外部自己同型群は未判定である。

\subsection{十一固定点と二次曲線上の五対}
十一パラメータの同型類を持つ $Z_2$ の非空開部分で、
thin-moduleスキームと左右の点スキームは被約な二十一点である。
外側の反対合は $D=\operatorname{diag}(1,-1,1)$ なので、
右平面の座標を $y=Dz$ により左平面と同一視する。
\textbf{この同一視後}の標準対応が、十一点を固定し、残る十点を五対に交換する対合である。

中央の固有基底で $J_2=\operatorname{diag}(1,-1,-1,-1,-1,-1)$ とし、
左関係行列を $L=\binom{a}{B_5}$ と書く。
可逆な対称七次行列 $A$ が関係のpairingを与え、
\[
 aAB_5^t=0,\qquad q(x)=a(x)A a(x)^t
\]
となる。右frameは元の座標では
$Q_j=\varepsilon_j J_2L_jA$、$\varepsilon=(1,-1,1)$ である。
この二次形式 $q$ は非退化であり、点配置は
\[
 Z=Z_{\mathrm{fix}}\sqcup Z_{\mathrm{pair}},\qquad
 |Z_{\mathrm{fix}}|=11,\quad |Z_{\mathrm{pair}}|=10,
\]
\[
 Z_{\mathrm{fix}}\cap V(q)=\varnothing,\qquad
 Z_{\mathrm{pair}}=Z\cap V(q)
\]
と分かれる。従って一般の二十一点配置のuniform positionは、
この自己反対部分族に制限すると成立しない。

十一という数は、$B_5$ の階数落ちをベクトル束
\[
 \mathcal E=\ker(\OO^{7}\xrightarrow{aA}\OO(1))(1)
 \simeq\Omega^1_{\PP^2}(2)\oplus\OO(1)^4
\]
の五つの切断の退化スキームとして数えることで得られる。
$\rk\mathcal E=6$、$c_2(\mathcal E)=11$ である。
固定点が $q=0$ 上にあると $Aa^t\in\ker L$ が階数落ちの法線微分を消し、
被約性に反する。このため十一固定点は二次曲線の外にあり、
残り十点が二次曲線上の五対となる\cite{PaperII}。

非空性の証明は、別の固定成分上の数値例ではなく、
境界点 $B_-$ を通る同じ $Z_2$ 上での代数的Hensel構成を使う。
$L$ と正規化した対称 $A$ の百五十三変数に対する六十方程式には、
$B_-$ で行列式二のJacobianブロックがある。
残る九十三変数を一つの有理接方向に沿って動かすと、
$\mathbb Z_{(101)}[[T]]$ 上の一意な形式解が得られる。
$T=101$ での代数的特殊化について、二十一点の評価代数の判別式と
$q$ の非退化性が厳密な $101$-進付値で非零と判定される。
被約性と非退化性は開条件なので、これで所要の非空開部分を得る。
再計算の入口は \nolinkurl{verify_self_opposite_geometry_v24.py} であり、
全六十方程式と有限代数の乗法行列を入力から再構成する。

この分解により、指定反対合を持つ枠付き関係族上のmonodromyは
$S_{11}\times(C_2\wr S_5)$ の部分群となる。
次の標準形を用いると等号が成立し、一般の全外部自己同型群の消滅と
標識を忘れた後の復元まで得られる。

関連する補助事実として、奇数個の集合への推移的群作用と可換な対合は恒等である。
実際、奇数個の集合上の対合は固定点を持ち、その固定点集合は推移的群で不変である。
基礎関数体上で定義された二乗恒等の反対合を持ち、thin-moduleスキームが
左右の点スキームと同型な奇数次数の有限エタール被覆となり、
そのGalois作用が推移的である場合、反対合が誘導する点置換は恒等となる。
従って左右の標準対応は射影変換で実現され、対応する線形系の階数は八以下となる。
この適用では反対合を既に指定しているため、$\Out=1$ を別に仮定する必要はない。
$Z_2$ の二十一点には実際に五つの二点軌道があるため、その標識付きmonodromyは推移的ではない。
上の十一点と十点への分解は、この一般的な群論的制約の具体的な実現でもある。

\subsection{重み付き五直線の標準形と単有理性}
滑らかな二次曲線 $Q=V(q)$、五本の異なる直線 $H_i=V(\ell_i)$、
重み $w_i\ne0$、$\sum_iw_i=1$ を取る。
$Q+\sum_iH_i$ は単純正規交叉とし、
\[
 U=\PP^2\setminus(Q\cup\textstyle\bigcup_iH_i),\qquad
 \omega_w=\sum_{i=1}^5w_i\frac{d\ell_i}{\ell_i}-\frac12\frac{dq}{q}
\]
と置く。Eulerベクトル場への評価は零なので、これは射影平面上の有理的一形式である。
$Z_2$ の一般点は、$\omega_w$ の十一零点と五組の $Q\cap H_i$ の和として表される。
十一零点は境界を避け、$Q\cap H_i$ は指定反対合の二点軌道である。

関係行列も平方根を使わず構成できる。
$q=x\cdot x$、$\ell_i=v_i\cdot x$ とし、
\[
 K=\{z\in\kk^5:\textstyle\sum_i z_i=0\},\quad
 \langle z,z'\rangle_K=\sum_i\frac{z_i z'_i}{w_i},\quad d_i=e_i-w
\]
と置く。この形式は非退化で、
\[
 \langle d_i,d_j\rangle_K=\frac{\delta_{ij}}{w_i}-1,
 \qquad \sum_iw_i d_i=0.
\]
$\kk^3\mathbin\perp K$ の直交和形式を $A$ とすれば、六行は
\[
 a(x)=(x,0),\qquad B_i(x)=(v_i\times x,\ell_i(x)d_i)
\]
となる。外積の恒等式と上のGram恒等式により、混合ブロックだけでなく
全自己反対incidenceを満たす。$B$ の階数落ち条件は、$U$ 上で
$\sum_iw_i v_i/\ell_i=x/q$、すなわち $\omega_w=0$ に一致する。

逆構成で再利用する点は、全incidenceから得られる五次の可換な自己随伴行列族である。
その行列式を $Q$ に制限すると、非固定十点の被約和を切り出す。
従って同時固有指標は五つ全て異なり、五本の直線が順序を除き一意に定まる。
固有ベクトルを $v_i^W=(r_i,d'_i)\in\kk\oplus K$、その長さの二乗を $g_i$ とすると
\[
 w_i=\frac{r_i^2}{g_i},\qquad \sum_iw_i=1
\]
となる。この比は固有ベクトルのスカラー倍、補助点、関係枠の変更に依存しない。
有限性により $r_i\ne0$ なので重みも非零である。

従って、指定反対合付きの同型類は一般に
\[
 \{(Q;(H_i,w_i)_{i=1}^5):\sum_iw_i=1\}
 \big/(\PGL_3\times S_5)
\]
と双有理である。次元は $5+10+4-8=11$ であり、
$Q$ を固定した十四次元の有理パラメータ空間から支配されるので単有理である。
この結論は有理性そのものの証明ではない。
同じ成分 $Z_2$ を記述することは、既存の非空開部分への逆構成と、
固定スキームの滑らかさによる成分の一意性で保証される\cite{PaperII}。

小さな有理例には
\[
 (\ell_i)=(x_1,x_2,x_3,x_1+x_2+x_3,x_1+2x_2+3x_3),
 \qquad (w_i)=\frac1{15}(1,2,3,4,5)
\]
を使える。全incidence、Torelliの階数条件、有限被約な二十一点性の再計算入口は
\nolinkurl{verify_master_normal_form_v25.py} である。

\subsection{独立な二因子と全monodromy}
指定反対合を持つ枠付きの関係族で、幾何学的monodromyは
\[
 \boxed{G_{Z_2}=S_{11}\times(C_2\wr S_5)}
\]
である。粗モジュライ上の普遍被覆を先に仮定する必要はない。
この決定を他の族へ応用するときは、次の三段階を分ける。
\begin{enumerate}
\item $Q,H_i$ を固定し、重みだけを動かす。
$\theta_i=d\log\ell_i-\frac12d\log q$ の評価写像
$\OO_U^5\to\Omega_U^1$ は全射で、核の射影束である臨界点対応は既約である。
対数余接束の第二Chern類は十一なので、一般の零点被覆は次数十一で推移的となる。
\item 重み空間に単独のfoldがある。
厳密な有理例では、原点のHessianは $\operatorname{diag}(0,-203/8)$、
核方向の三階微分は $-65$、横断する重み方向の一次係数は一である。
飽和した零スキームは長さ二の原点と九つの被約点に分かれ、境界と無限遠に余分な零点はない。
従ってmonodromyは一つの互換を含む。
素数次数十一の推移的群は原始的なので、固定点上の群は $S_{11}$ となる。
この変形では十個の $Q\cap H_i$ を一つずつ固定するため、部分群 $S_{11}\times1$ を得る。
\item 一本の直線を双対二次曲線の周囲で動かすと、その交点対だけが交換される。
五本で独立に行い、さらに $(H_i,w_i)$ の置換を実現すると、第二因子への全射を得る。
枠の変更は指定反対合を保つ連結な一般線形群内の道で閉じられる。
第一因子全体を既に含むので、二因子の独立性も従う。
\end{enumerate}
単独foldの有理数証明書の入口は \nolinkurl{master_monodromy_v25.py} である。
有理数上の特殊化Galois群から複素monodromyを推測する議論ではなく、
既約な対応と実際の複素ループを使う。

\subsection{標識を忘れた配置からの復元と $\Out=1$}
固定十一点 $T$ 自体にはuniform positionが成立する。任意の部分集合 $T'\subset T$ で
\[
 \operatorname{rank}\bigl(H^0(\PP^2,\OO(d))\to H^0(T',\OO_{T'}(d))\bigr)
 =\min\{\tbinom{d+2}{2},\#T'\}.
\]
退化束の分解が $T$ のHilbert関数を与え、$S_{11}$ が同じ個数の部分集合に推移的であることから従う。
従って六固定点が同じ二次曲線に乗ることはない。
B\'ezoutを用いると、配置の十点以上を含む二次曲線は $Q$ に限る。

一般の配置の射影自己同型群 $H$ は有限で、点分解体のGalois群により正規化される。
$Q$ の一意性から $H$ は固定点と非固定点をそれぞれ保つ。
$h\in H$ と $\sigma\in S_{11}\times1$ の交換子は $Q$ 上の十点を点ごとに固定する。
そのうち四点は射影枠なので交換子は恒等である。
十一文字上の $S_{11}$ の中心化群は自明であり、固定点も射影枠を含むので $h=1$ となる。
二十一点Torelliによる安定化群の同型と合わせ、非空開部分で
\[
 \operatorname{Aut}_{\PP^2}(Z)=1,\qquad \Out(A_Z)=1
\]
を得る\cite{PaperII}。

二つの反対同型の合成は代数自己同型なので、反対同型は内的自己同型を除いて一意となる。
内的自己同型はthin moduleの同型類に自明に作用するため、点配置の対合も一意に定まる。
従って標識を忘れた配置から、$Q$、五つの交換対、それらを結ぶ直線、重みを復元できる。
先の双有理記述と単有理性は、標識を忘れた十一次元自己反対モジュライにも成立する。

逆向きの判定で必要なのは、十点が $Q$ 上にあることに加えて、
残り十一点が対応する $\omega_w$ の零点となり、標準形の関係データがTorelliの開条件を満たすことである。
「十点が二次曲線上」という条件だけを自己反対性の十分条件とはしない。
この復元は三固定成分の全分類や、その飽和の閉包間の包含も決定しない。

\subsection{全三族を統一する署名公式}
右平面との同一視と中央基底の変更の後、
\[
 J=\operatorname{diag}(I_p,-I_q),\quad p+q=6,\quad 1\le p\le3,
 \qquad L=\binom{P}{B_q}
\]
とする。$A$ は対称可逆で、全incidence
$L_iAL_j^tJ=JL_jAL_i^t$ を仮定する。
小ブロック $P$ が平面上で常に階数 $p$ を持ち、点スキームが有限被約なら、
\[
 E=\ker(\OO(1)^7\xrightarrow{PA}\OO(2)^p),\qquad
 c(E)=\frac{(1+H)^7}{(1+2H)^p}
\]
より固定点数は $f(p)=21-12p+2p^2$ である。
残る $2p(6-p)$ 点は $p(6-p)$ 対となり、
$\det(PAP^t)=0$ という次数 $2p$ の曲線上にちょうど乗る。
固定点でGram核があれば $k=AP^tc$ が非零の法線方向を消してしまうので、固定点はその外にある。
$p>1$ では、混合式 $PAB_q^t=0$ に加え、両方の同符号ブロックのincidenceも必要である。

\[
\begin{array}{c|c|c|c|c|c}
 & (p,q)&\text{固定点}&\text{非固定点}&\text{二点軌道}&\text{曲線次数}\\\hline
 Z_1&(3,3)&3&18&9&6\\
 Z_2&(1,5)&11&10&5&2\\
 Z_3&(2,4)&5&16&8&4
\end{array}
\]
この曲線は指定反対合と全階数の固有ブロックに付随する。
$L\mapsto LM$, $A\mapsto M^{-1}AM^{-t}$ でGram行列は不変であり、
ブロック基底変更では行列式が非零定数倍されるだけである。
従って関係枠の任意性を除いて比較できる。

三族全てで、この表の条件を満たす非空開部分が存在する。
基点の小ブロックの最大小行列式は全次数 $p$ の単項式を生成するので、
全階数条件は基点とその近傍で成り立つ。
全方程式数、基点Jacobian小行列式、入力精度、乗法行列の保証精度、
判別式の付値は次のとおりである。
\[
\begin{array}{c|r|r|r|r|r}
 &\text{式数}&\det J_{\rm pivot}&N&N-w&v_{101}(\operatorname{Disc})\\\hline
 Z_1&72&-1&1536&1533&1120\\
 Z_2&60&2&1024&1022&700\\
 Z_3&69&12288&1024&1022&686
\end{array}
\]
各行で判別式の付値が保証精度より小さいため、真のHensel解で非零である。
追加の二署名の入口は \texttt{hensel\_other\_signatures.py 3} と
\texttt{hensel\_other\_signatures.py 2} であり、既定では証明書を変更しない。
境界の $B_-$ に点族の有限平坦性を仮定しない。
関係族の形式曲線で成分への所属を保証し、被約点配置の開部分で有限エタール被覆を取る。

三族とも、$\Gamma_P=V(\det(PAP^t))$ が滑らかで幾何学的に既約となる非空開部分を持つ。
六次と四次の種数はそれぞれ十と三である。証明に使う有限な非消滅判定は次節に記す。
指定反対合を保つ点monodromyについて、現在の確定範囲は
\[
 \begin{aligned}
 G_{Z_1}&\subseteq S_3\times(C_2\wr S_9),\\
 G_{Z_2}&=S_{11}\times(C_2\wr S_5),\\
 G_{Z_3}&\subseteq S_5\times(C_2\wr S_8)
 \end{aligned}
\]
である。$Z_1$ の固定三点への射影は $S_3$ と確定している。
$Z_1$ の非固定側の最大性と固定側からの独立性、
$Z_3$ の両射影の全射性と独立性は未決である。

同じ計数の再利用先も保持する。
$m\times(m+1)$ の一次式行列、署名 $(p,m-p)$、対称可逆pairing、
全incidence、$P$ の全階数、有限被約性という同じ仮定では、
第二Chern類の計算は
\[
 |Z|=\binom{m+1}{2},\quad
 |Z_{\rm fix}|=\binom{m+1}{2}-2p(m-p),\quad
 |Z_{\rm pair}|=2p(m-p)
\]
に置き換わる。これは仮定下の計数公式であり、そのような関係成分や被約点配置の実在を
一般の $m$ について主張するものではない。

\subsection{滑らかな四次・六次曲線とPrymデータ}
小ブロックのGram行列を $Q_P=PAP^t$、曲線を $C=V(\det Q_P)$ とする。
同じ固定成分上の代数的Hensel点で滑らかさを保証する証明書は、
$F=\det Q_P$ の三偏微分が次数 $3d-5$ の全斉次式を生成することを検査する。
ここで $d=2p$、各偏微分に掛ける単項式の次数は $2d-4$ である。
\[
 \begin{array}{c|c|c|c|c}
  &d&\text{Macaulay行列}&\text{全列階数小行列式の付値}&\text{保証精度}\\\hline
 Z_3&4&45\times36&11&256\\
 Z_1&6&135\times105&57&256
 \end{array}
\]
付値は素数 $101$ に関するものであり、いずれも保証精度未満である。
従って真の標数零の点で偏微分の射影的共通零点はなく、滑らかさの非空開部分を得る。
これは被約二十一点性と同じ成分上の同じHensel点を使う。
入口は \nolinkurl{gram_curves_v25.py}、証明書は
\nolinkurl{gram_curves_v25_certificate.json} である。

$i:C\hookrightarrow\PP^2$ とし、対称行列表示
\[
 0\longrightarrow\OO(-2)^p\xrightarrow{Q_P}\OO^p
 \longrightarrow i_*\mathcal L\longrightarrow0,
 \qquad \eta=\mathcal L(1-p)
\]
から線束を定める。滑らかさにより $Q_P$ は $C$ 上で常に階数 $p-1$ なので、
$\mathcal L$ は局所自由階数一である。
この分解を $\OO(-2)$ で双対化すると
\[
 \mathcal L^{\otimes2}\simeq\OO_C(2p-2),\qquad
 \eta^{\otimes2}\simeq\OO_C
\]
を得る。$p=2,3$ では、分解を $1-p$ だけねじったコホモロジーから
$H^0(C,\eta)=0$ となるので $\eta$ は非自明である。
$p=1$ では $\eta=\OO_C$ となる。

非固定点の被約除数を $D_{\rm mov}$、平面の直線切断を $H$ とすると
\[
 \boxed{\quad\OO_C(D_{\rm mov})\simeq\eta(q),\qquad
 \eta\simeq\OO_C(D_{\rm mov}-qH).\quad}
\]
この式の再利用には次の束の解釈が有効である。
$\mathcal S=\operatorname{im}P^t\simeq\OO_C(-1)^p$ の根基は
$\mathcal T=\mathcal S\cap\mathcal S^\perp\simeq\eta(-p)$ である。
固定点は $C$ の外にあるため
$\mathcal B=\operatorname{im}B^t\simeq\OO_C(-1)^q$ も部分束となり、
$\mathcal S^\perp/\mathcal B\simeq\OO_C(q-p)$ である。
自然な射 $\mathcal T\to\mathcal S^\perp/\mathcal B$ の零点は、
$L$ の階数が落ちる非固定点ちょうど全部である。
切断の次数 $2pq$ が既知の非固定点数と一致するので、全零点の重複度は一となる\cite{PaperII}。

従って $Z_3$ と $Z_1$ は、それぞれ種数三と十の曲線と非自明二捩れ線束を与える。
$\operatorname{Spec}_C(\OO_C\oplus\eta)$ は連結エタール二重被覆であり、
その種数は五と十九である。特に六次曲線の場合は $p=q=3$ なので
\[
 K_C=\OO_C(3),\qquad
 D_{\rm mov}\in|K_C\otimes\eta|,\qquad
 h^0(C,K_C\otimes\eta)=9.
\]
十八非固定点はPrym-canonical線形系の被約除数となる。
四次曲線では $K_C=\OO_C(1)$ であり、十六非固定点の線束は $\eta(4)$ である。

\subsection{補助Gram曲線の二重交差と適用範囲}
反対側のGram行列式
\[
 \Gamma_B=V(\det(BAB^t))
\]
は次数 $2q$ の非零曲線で、$C$ を成分に持たず、
$C$ 上のCartier除数として
\[
 \Gamma_B|_C=2D_{\rm mov}
\]
となる。非退化対称束 $\mathcal H=\mathcal S^\perp/\mathcal T$ への
$\mathcal B\to\mathcal H$ の行列式が $D_{\rm mov}$ を切り出し、
Gram行列式がその二乗となるためである。
六次の場合、二つの六次曲線の交差は十八点に交差重複度二ずつ集中する。
また固定点では $B$ の階数が落ち、行列式の微分も消えるので、
全ての固定点は $\Gamma_B$ の特異点である。
一般の補助曲線は幾何学的に既約で、固定点だけが通常二重点である。
六次では三節点、八次では五節点となり、非固定点では滑らかである。
証明と再利用できる法線形式は第\ref{sec:notes-prym-v27}節に記す。

$\eta$ のエタール二重被覆は曲線 $C$ の上の被覆であり、
指定反対合による $D_{\rm mov}$ の点の二個ずつの交換とは異なる対象である。
両者の同一視や、点の交換が $C$ 全体の自己同型へ延長することは主張しない。
$p=3$ では $y^tQ_P(x)y=0$ が $\PP^2\times\PP^2$ の双次数 $(2,2)$ 超曲面を与え、
第一射影のconic bundleの判別式と二捩れ線束は $(C,\eta)$ となる。
本族は特別な部分族なので、一般のconic bundleの有理性に関する定理をそのまま転用しない。

\subsection{復元定理から広がる幾何学的問題}
最初に調べるべきものは、点配置上の rank 条件の意味である。
二十四次元という値だけでなく、特殊な線形系・極性・対応や determinantal な対称性として説明できれば、
十六次元商の有理性や境界を同時に扱える可能性がある。
左右交換の対合については、固定部分、境界への延長、商の有理性が具体的な対象となる。

一方、代数の Torelli 定理が得られても導来圏の Torelli 定理は未解決である。
球面 twist が生成する自己同値は、有限代数の外部自己同型とは別物である。
この違いを維持しながら、例外列の変異による点配置の変換を調べる。

\section{Prym 写像の像と逆構成を使うためのデータ}
\label{sec:notes-prym-v27}

主要な分類の結論は次のように分かれる。ここで $\mathcal R_g$ は
曲線と非自明二捩れ線束の対のモジュライであり、Prym 多様体のモジュライではない。
\begin{center}\small
\begin{tabular}{p{.15\linewidth}p{.20\linewidth}p{.53\linewidth}}\toprule
署名と次元&像&復元と残るファイバー\\\midrule
$(2,4)$、八次元&$\mathcal R_3$ を支配&一般関数体上で $E\times D$ と双有理。$E$ の周期・指数は三、$D$ は非分裂二次曲線、ファイバーの指数・最小多重切断次数は六。\\
$(3,3)$、七次元&$\mathcal R_{10}$ 内の七次元像&抽象 $(C,\eta)$ から双有理復元。$W_6^2(C)=\{H\}$ は被約。指定反対合を保てば固有ブロックの順序も不要。\\\bottomrule
\end{tabular}\end{center}
幾何と六次の復元は公表稿V、net写像の相対有理性と一般関数体上の算術は公表稿VIに全証明を置く。
指定反対合そのものを忘れた代数の粗モジュライへの写像は、この表の結論に含めない。

\subsection{全incidenceから係数Gram行列を読む}
$P$ の行数を $p$ とし、その係数行を $(u,i)$ の順に並べた行列を $\widehat P$ とする。
$Q=PAP^t$ に対して
\[
 G(Q)_{(u,i),(v,j)}=
 \begin{cases}\left[x_i^2\right]Q_{uv},&i=j,\\
 \left[x_ix_j\right]Q_{uv}/2,&i\ne j\end{cases}
\]
を定めると、同符号の全incidenceから
$G(Q)=\widehat P A\widehat P^t$ となる。
混合直交性だけではこの別々の対称性は得られない。
この係数の偏極が、二つの復元法の共通の出発点である。

\subsection{四次の場合：支配性と残余符号}
正規化チャートは百五十三変数、六十九方程式、接次元八十四である。
同じ seed $17029$ の Hensel 分枝上で
\[
 dQ|_{T}=D_t-D_zJ_z^{-1}J_t
\]
の十八次小行列式は $101$ 進付値二、入力精度は $101^{32}$ である。
$J_z$ は単元で、逆行列の更新を最終補正座標で再確認しているため精度損失はない。
従って十八個の Gram 係数は独立に動く。
任意の非超楕円種数三 Prym 対は、基点を持たない $|K_C\eta|$ の二つの切断の積から
対称二次行列表示を持つので、これが $\mathcal R_3$ への支配性を証明する。

一般に $G(Q)$ は可逆で、関係枠を変えて
\[
 P(x)=\begin{pmatrix}x^t&0&0\\0&x^t&0\end{pmatrix},\qquad A=\operatorname{diag}(G,1)
\]
とできる。一次直交 syzygy 空間は $K=K_6\oplus K_3$、$\dim K_6=6$、$\dim K_3=3$ である。
$B$ の四行は $\operatorname{Gr}(4,K)$ の点 $W$ を与え、
三つの交代形式について同時等方的でなければならない。
標準化に残る $\tau(c,\ell)=(c,-\ell)$ を必ず商する。
この符号を忘れると、逆構成に現れる二価性を余分な被覆と取り違える。

pairing を正規化しないチャートの次元は八十五であり、
固定 Gram ファイバーの商は $85-18-49-16=2$ 次元である。
$A_{77}=1$ の八十四次元チャートでは、Gram を固定して動ける関係枠は四十八次元である。
八十四から四十九を差し引くことは、同じスカラーを二度消している。

\subsection{楕円線織ファイバーの逆構成}
$D_i=\operatorname{diag}([e_i]_\times,[e_i]_\times)$ とし、
\[
 \Omega_{ij}=D_i^tG^{-1}D_j-D_j^tG^{-1}D_i,\qquad
 E_Q=V\bigl(\operatorname{Pf}(z_1\Omega_{23}+z_2\Omega_{31}+z_3\Omega_{12})\bigr)
\]
とする。一般に $E_Q$ は滑らかで、その上の核束 $\mathcal K_Q$ は階数二である。
\[
 (X_Q\text{ の有限配置開部分})/\langle\tau\rangle
 \dashsim\mathbb P_{E_Q}(\mathcal K_Q).
\]
必要な条件 $\operatorname{rank}(W\to K_3)=3$ は同じ Hensel 点で証明されている。
有限配置条件は $W\cap K_3=0$ を強制する。非零な交わりがあれば、
$B$ の一行が $\ell(x)e_7$ となり、直線 $\ell=0$ 全体が階数低下集合に入るためである。

逆構成では $k=W\cap K_6$ を取り、
\[
 F_k=\{v\in K_6:k^t\Omega_{ij}v=0\text{ for all }i<j\}
\]
を作る。滑らかな三次曲線が直線を含まないことから、
$k\subset\ker\Omega(z)$ を満たす $[z]\in E_Q$ は一意で、$\dim F_k=4$ である。
$W$ は $F_k\to K_3$ のグラフとなり、三次元商 $F_k/k$ 上で
指定された可逆写像を $\bigwedge^2\phi$ に持ち上げる。
$\bigwedge^2:\operatorname{GL}_3\to\operatorname{GL}_3$ の核は $\{\pm1\}$ なので、
二つの持上げが $\tau$ でちょうど同一視される。

一般の平面四次曲線に自己同型がないことと、対称最小分解の一意性により、
粗 Prym ファイバーにさらに有限商は生じない。
ただし、上の階数二核束は一般関数体上で実際のベクトル束へは降りず、
その射影化は非零Brauer classを持つ二次曲線束へ降りる。
この算術的な違いを第\ref{subsec:notes-brauer-v27}節で記述する。

\subsection{六次の場合：有理式による再構成}
九行七列の $\widehat P$ は階数七、混合方程式 $PAb^t=0$ の十八行二十一列の行列は階数十八を持つ。
いずれも同じ Hensel 点で単元小行列式を持つので、非空開条件である。
$G(Q)$ の可逆な主七次部分行列 $G_I$ を選び、
\[
 A_0=G_I,\qquad \widehat P_0=G(Q)_{*,I}G_I^{-1}
\]
とすると、$G(Q)=\widehat P_0A_0\widehat P_0^t$ である。
他の階数七の分解との差は $\operatorname{GL}_7$ だけである。
次に三次元核 $K_{P,A}$ の基底を取れば、$B$ も行変換を除いて復元される。

この手順は像の判定法にもなる。任意の三行三列の対称二次式行列に対して、
まず $\operatorname{rank}G(Q)=7$、混合核の次元三を調べ、
その核上の三つの交代形式が消えることを検査する。
これに既存の開階数と有限配置条件を加えると全incidenceを再構成できる。
階数七だけから残りの九つの同符号方程式が自動的に成立すると考えてはならない。
どの既約成分が得られるかも、指定した $Z_1$ との照合が必要である。

平面 Prym データは対称最小分解から $Q$ を合同変換まで復元する。
さらに以下の一意性により、$H$ を指定せずに抽象対 $(C,\eta)$ から復元できる。
従って $\mathfrak Z_1^+\dashrightarrow\mathcal R_{10}$ は七次元の像への双有理写像であり、
全二十一点配置と特定の $D_{\rm mov}\in|K_C\eta|$ も戻る。

\paragraph{再利用できる平面直線束の一意性.}
標数零の代数閉体上の滑らかな平面曲線 $C$ の次数を $d\ge4$ とする。
次数 $d$ 以下で三つ以上の切断を持つ線束は $H=\mathcal O_C(1)$ だけであり、
Brill--Noether軌跡 $W_d^2(C)$ は被約な一点である。
これはNoether理論の特別な場合である\cite{HartshorneNoether,CoppensKatoPlaneSystems}が、
本研究で必要な証明は次のように短く閉じる。

固定因子を除いた線束 $M'$ の次数を $n\le d$ とし、一般の被約除数を $D$ とする。
随伴公式とRiemann--Rochにより、次数 $d-3$ の平面形式の評価階数は
\[
 \operatorname{rank}(\mathrm{ev}_D)=n+1-h^0(M')\le n-2.
\]
$n\le d-2$ なら、各点を、他の $n-1$ 点を一つずつ通る直線の積で分離でき、階数は $n$ となる。
$n=d-1$ で非共線なら、一点 $p$ を避けて残りの二点を結ぶ直線と、
残る $d-4$ 点を通り $p$ を避ける直線の積で全点を分離できる。
$d=4$ でも三つの非共線点を一次式で分離する証明そのものである。
共線なら階数は $d-2$ なので、いずれも上界に反する。
従って $n=d$ である。非共線な $D$ は非共線な $d-1$ 点を含むので再び矛盾し、
$D$ は直線上にある。B\'ezoutにより $D=C\cap\ell$、従って $M'=H$ となり、
次数から元の固定因子も零である。

$[H]$ における $W_d^2(C)$ の接空間は、Serre双対性によって乗法
\[
 H^0(H)\otimes H^0(K_C-H)\longrightarrow H^0(K_C),
 \qquad S_1\otimes S_{d-4}\longrightarrow S_{d-3}
\]
の像の零化空間である。この乗法は全射なので接空間は零となり、Nakayamaの補題が被約性を与える。
また、平面曲線の抽象同型は $H$ を保つため、射影変換へ一意に延長する。
六次で $H^3=K_C$ の根だけを数えると失われる情報は、切断数 $h^0(H)=3$ である。
この復元は曲線を保持しており、Prym多様体だけからのTorelliを意味しない。

\paragraph{固有ブロックの順序を戻す規則.}
署名 $(3,3)$ の有限配置開部分では、$P$ は平面全体で階数三、
$B$ はちょうど三固定点で階数低下する。
従って、二つのブロックのうち全点で最大階数を持つものを $P$ として選べる。
pairingのスカラー正規化を戻した枠付き開部分を $\mathcal U$、
順序を保つ枠・座標・pairing変更群を $\mathcal G_{\rm ord}$ とする。
$s(P,B,A)=(B,P,A)$ は $J$ を $-J$ に移し、全incidenceを保つ。
$\mathcal U\cap s\mathcal U=\varnothing$ であり、
\[
 (\mathcal U\sqcup s\mathcal U)/(\mathcal G_{\rm ord}\rtimes\langle s\rangle)
 \dashsim \mathcal U/\mathcal G_{\rm ord}
 \dashsim\mathfrak Z_1^+.
\]
交換は別の枠付き成分を動かし得るので、順序を忘れる商はこの合併の商として定義する。
一つの成分に交換が作用すると無条件に仮定しない。
ここで不要になったのは固有ブロックの順序であり、指定反対合そのものの一意性と
$Z_1$ の一般の全外部自己同型群は未判定である。

\subsection*{六次Prym像の再利用可能な標準形}

階数七の係数Gramの核は二次元pencilである。
行列式が相異なる三根を持つ開部分では、これを対角の跡零pencilに正規化できる。
正規化後の $P$ と三次元の一次直交syzygy基底 $R$ は
論文Vの「九本の余因子方程式」で与えた固定行列となり、
全ての変形情報は十九変数の対称形式 $A$ に移る。
負ブロックは $B=RA^{-1}$ であり、残る条件は
\[
 \begin{aligned}
 (A^{-1})_{14}&=(A^{-1})_{15}=(A^{-1})_{16}=0,\\
 (A^{-1})_{23}&=(A^{-1})_{24}=(A^{-1})_{26}=0,\\
 (A^{-1})_{35}&=(A^{-1})_{36}=(A^{-1})_{45}=0.
 \end{aligned}
\]
である。添字は $0,\ldots,6$。
従って、後続の計算を百五十三変数・七十二方程式のチャートから、
十九変数・九余因子方程式のチャートへ縮小できる。
この正規化に必要な三根の順序付けは、最後に $\mathfrak S_3$ で除く。

$A_0$ は $(0)$ の成分が $1$ で、$(1,3),(2,5),(4,6)$ を交換する行列である。
$\det A_0=-1$。
上記の順序の九余因子の $(p,m,i,r,o,l,b,g,f)$ に関するJacobi行列は $I_9$ である。
従って局所計算には $(s-1,a,c,d,e,h,j,k,n,q)$ の十変数を自由変数として用いられる。
これは九個の変数を形式的・エタール局所的に消去できるという主張であり、
大域的な有理パラメータ表示を与えるものではない。
また $A_0$ ではトーラスの安定化群が一般点より大きいため、
十次元チャートの滑らかさだけから粗商の基点の滑らかさを推論しない。

正則核pencilを持つ階数七Gram軌跡は、十四次元のpencil軌道上の
十九次元線形空間族の開部分であり、滑らかな三十三次元空間である。
九方程式は既存成分に沿う一般点で余次元九を持つ。
選択した十次元成分の有効安定化群は
$\mathbb G_m^3\rtimes\mathfrak S_3$ であり、その商は七次元である。
全余因子軌跡の既約性や他の成分は未判定。
成分の選択は既存の整数基点 $B_-$ と同定した $A_0$ により固定する。

\paragraph{標準的な直線とPrym-canonical切断.}
核pencilの三つの特異元から、両三次元空間の固有直線分解を取り戻せる。
対応する三つの階数一テンソルが張る三次元空間 $E$ は根の順序に依存せず、
$\ell=E/K\subset (U\otimes V)/K$ は標準的な直線となる。
これは標準形の $\langle e_0\rangle$ である。
$v\in\ell\setminus\{0\}$ に対する $t=PAv$ の余核像は
$H^0(C,K_C\otimes\eta)$ の標準的な一次元部分空間を定める。
十八点除数は直接
\[
 I_{D_{\rm mov}}=(\det Q,\operatorname{adj}(Q)t)
\]
で復元される。
その最小分解の次数は、生成元が $(5,5,5,6)$、syzygyが $(7,7,7)$。
三つの五次生成元は固定三角形の三頂点で可逆な対角評価行列を与える。
従って $h^0(I_D(5))=3$、$h^0(I_Z(5))=0$。
固定三点は核pencilだけから、移動十八点はこの標準切断から得られる。

補助六次曲線は、標準形で $v=e_0$ と正規化し、$s=v^tAv$ と置くと
\[
 \det A\det(RA^{-1}R^t)
 =s\det Q-t^t\operatorname{adj}(Q)t
\]
から得られる。この零曲線は $v$ の非零定数倍に依存せず、$Q$ と標準直線から計算できる。

\paragraph{次に調べる価値のある問題.}
十次元の余因子成分と三次元トーラスの表示により、有理性・不変量・分岐の計算が具体化した。
当面は九余因子方程式の大域的な消去や低次の不変量を調べることができるが、
形式局所解を有理解と取り違えないことが必要である。
全monodromyと反対合の標識忘却には、標準形が与える固定三角形以外の情報が要る。
また、九条件を単なる階数七条件から自動的に導いてはならない。
階数七・正則pencilの三十三次元軌跡と、本研究の二十四次元枠付き像は異なる。

\subsection{Hermite--Pfaffian恒等式と二次同種写像}
\label{subsec:notes-hermite-v27}

四次の対称表示を
\[
 Q=\begin{pmatrix}a&b\\b&c\end{pmatrix},\qquad
 a=x^tA_0x,\quad b=x^tA_1x,\quad c=x^tA_2x,
 \qquad G=\begin{pmatrix}A_0&A_1\\A_1&A_2\end{pmatrix}
\]
と書き、二次曲線のnetを $N=\langle a,b,c\rangle$ とする。
双対平面上の普遍階数二部分束 $\mathcal S$ に対する制限
\[
 r_N:N\otimes\mathcal O\longrightarrow\operatorname{Sym}^2\mathcal S^*
\]
の行列式を $H_N$ と書く。これは三次式であり、$H_N(z)=0$ は
netの非零な二次曲線が直線 $\ell_z$ を含むことと同値である。
具体的には $x=(z_3u,z_3v,-z_1u-z_2v)$ を代入し、
$a,b,c$ の係数を $u^2,uv,v^2$ の順に並べた行列 $R_N$ を使って
$H_N=\det(R_N)/z_3^3$ と計算できる。
分母は消え、全双対平面上の式となる。

既存の外積・Pfaffianの規約では、正確な恒等式は
\begin{equation}\label{eq:notes-hermite-pfaffian-v27}
 2\det(G)\operatorname{Pf}\bigl(z_1\Omega_{23}+z_2\Omega_{31}+z_3\Omega_{12}\bigr)
 =H_N(z)
\end{equation}
である。証明は一つの直線での相補小行列式に帰着する。
$M=G^{-1}$ とすると
\[
 \operatorname{Pf}\Omega_{12}
 =\det M_{(1,3,6),(5,3,6)}-\det M_{(2,3,6),(4,3,6)}.
\]
Jacobiの公式により、$\det(G)$ を掛けた右辺は
\[
 \det\begin{pmatrix}
 (A_0)_{11}&(A_1)_{11}&(A_2)_{11}\\
 (A_0)_{12}&(A_1)_{12}&(A_2)_{12}\\
 (A_0)_{22}&(A_1)_{22}&(A_2)_{22}
 \end{pmatrix}=\tfrac12 H_N(0,0,1)
\]
となる。中央の係数が二倍になることが定数 $2$ の由来である。
座標変更 $x=gx'$ では両辺とも
$F(z)\mapsto\det(g)^3F(g^{-t}z)$ と変わるので、全座標で恒等式が成立する。
従って補助曲線 $E=E_Q$ は、古典的なnetのHermite曲線そのものである
\cite[著者公開原稿 \S3.2.2, Exercise 3.12]{DolgachevCAG}。
その同型類は、$Q$ の三成分の基底よりもnet自体に依存する。

判別式三次曲線を
\[
 \Delta_N=V\bigl(\det(\lambda_0A_0+\lambda_1A_1+\lambda_2A_2)\bigr)
 \subset\mathbb P(N)
\]
とする。$E$ と $\Delta_N$ が滑らかなら、直線成分をそれを含む退化二次曲線へ送る写像
\[
 \pi_N:E\longrightarrow\Delta_N
\]
は連結有限エタール二重被覆である。
実際、判別式の滑らかさにより退化二次曲線の階数は常に二で、二直線は相異なる。
一つの直線が独立な二つのnetの元に含まれると、そのpencilは全て退化し、
$\Delta_N$ に直線成分を生じるため、像は一意に定まる。
局所的には階数二の非退化二次形式の等方直線を選ぶエタール二重被覆となる。

被覆のtrace零線束は対称表示
\[
 0\longrightarrow\mathcal O_{\mathbb P(N)}(-2)^3
 \xrightarrow{\lambda_0A_0+\lambda_1A_1+\lambda_2A_2}
 \mathcal O_{\mathbb P(N)}(-1)^3\longrightarrow i_*\theta_N\longrightarrow0
\]
の余核である。対称性から $\theta_N^2=\mathcal O_{\Delta_N}$、
コホモロジーから $h^0(\theta_N)=0$ となる。
同定を直接確認するには、$\mathscr K=\ker(\mathcal O^3\to\mathcal O(1)^3)$、
$\mathscr V=\mathcal O^3/\mathscr K$ と置く。
等方直線被覆の判別式線束は
$\det\mathscr V\otimes\mathcal O(-1)=\mathscr K^*\otimes\mathcal O(-1)$ であり、
上の余核 $\theta_N$ と一致する。従って
\[
 (\pi_N)_*\mathcal O_E=\mathcal O_{\Delta_N}\oplus\theta_N,
 \qquad
 \operatorname{Pic}^0(E)
 \simeq\operatorname{Pic}^0(\Delta_N)/\langle\theta_N\rangle.
\]

対称最小分解の一意性により、$(\Delta_N,\theta_N)$ は一般Prym対の不変量である。
これが支配的有理写像
\[
 \Psi:\mathcal R_3\dashrightarrow Y_0(2),\qquad
 (C,\eta)\longmapsto
 (\operatorname{Pic}^0(\Delta_N),\langle\theta_N\rangle)
\]
を与える。支配性は後述の試験族の非定数な $j$ で確認できる。
Jacobianを用いるので原点の選択を要しない。
この二次同種写像から、種数三Prym対の二次元Prym多様体に楕円因子があるとは結論しない。
また、核束までnetだけに依存するという主張は含めない。

\subsection{net写像の相対有理性と降下で残るもの}
\label{subsec:notes-relative-rational-v28}

$\Psi:\mathcal R_3\dashrightarrow Y_0(2)$ は
Dolgachev\cite[Remark~4.1]{DolgachevRationality2008}の古典的なnet写像である。
本研究で得た相対的な結論は
\[
 \mathbb C(\mathcal R_3)=\mathbb C(Y_0(2))(u_1,\ldots,u_5)
\]
という、指定された底を保つ純超越表示である。公表稿VIに全証明を置く。
再利用すべき点は、一般netの非自明な安定化群を消さずに扱う方法である。

一般のnetの射影安定化群は
$G=(\mathbb Z/3)^2\rtimes C_2$ である。
$\operatorname{SL}_3$ 内の逆像 $\widetilde G$ を、その一意のSylow三部分群による
符号指標 $\chi$ でねじると、線形持上げは
三乗根の対角群と $S_3$ の半直積 $G(3,3,3)$ となる。
その不変式環は
\[
 k[p,q,r],\qquad \deg(p,q,r)=(3,3,6)
\]
であり、分裂形の生成元は
$p=x^3+y^3+z^3$、$q=xyz$、$r=x^3y^3+y^3z^3+z^3x^3$ である。
非定数群の形に対しても、次数三の二次元空間と次数六の一次元補空間を
基礎体上で選ぶことにより同じ多項式環が降りる。
従って $\mathbb P^2/G\simeq\mathbb P(1,1,2)$ は分裂している。

さらに $V'=V\otimes\chi$ とすると
\[
 \operatorname{Sym}^2V^*\simeq V'\otimes M,\qquad
 M=\operatorname{Hom}_{\widetilde G}(V',\operatorname{Sym}^2V^*),\quad\dim M=2.
\]
二つの部分表現を個別に降ろす代わりに、このHom空間を使う。
$\mathbb P(V'\oplus V')\dashrightarrow\mathbb P(V')$ の開部分は
階数三のHom束であり、中心 $\mu_3$ はこの束上で相殺される。
自由作用開部分で商を取れば真のベクトル束となるため、二次曲線空間の商は五次元有理である。
この議論は定数群の標準形を基礎体へ無条件に移す議論ではない。

モジュラー関数体上のnet代表は
\[
 \langle y^2+2xz,\ 2yz,\ -x^2+2xz-2y^2+uz^2\rangle
\]
と取れる。判別式のWeierstrass表示と $j$ は
\[
 Y^2=X(X^2+6X+9+u),\qquad
 j=\frac{1728(u-3)^3}{u(u+9)^2}.
\]
唯一の非零有理二捩れ $(0,0)$ と次数三の $j$ 写像により、$u$ は $Y_0(2)$ の座標となる。
入力の二次曲線は $\mathbb P(N)$ 内にあり、その双対を評価写像で引き戻すと四次Prym曲線になる。

底の関数体の有理性は、一般Hermite曲線やPrym-canonical二次曲線の分裂とは独立である。
有限安定化群による降下が残るため、周期三・指数二・積の指数六と両立する。
また $\Psi$ 全体が正則な射影空間束であるという結論ではない。


\subsection{一般Prymファイバーの非零Brauer障害}
\label{subsec:notes-brauer-v27}

$K=\mathbb C(\mathcal R_3)$ とし、曲線の自己同型が自明な開部分上の一般対を
\[
 C/K,\qquad 0\ne\eta\in\operatorname{Pic}^0_{C/K}(K)[2]
\]
とする。$\eta$ はまずPicardスキームの点であり、実際の $K$ 上の線束とは限らない。
その降下障害を
\[
 \alpha_{\mathrm{Prym}}=\delta_C(\eta)\in\operatorname{Br}(K)[2]
\]
と書く。これは元の十一点族のpotential障害 $\beta$ とは別の対象である。
Brauer classと射影線形群の降下cocycleの対応は
\cite[Tags 03R1, 0A2J]{StacksBrauerV27}を用いる。

$V_{\rm can}=H^0(C,\omega_C)$ は真の三次元 $K$ ベクトル空間だが、
$W=H^0(C_{\overline K},\omega_C\eta)$ は $\alpha_{\rm Prym}$ による階数二のtwisted空間である。
その射影化はBrauer--Severi二次曲線 $D/K$ へ降りる。
$\eta^2\simeq\mathcal O_C$ を選ぶ $\mu_2$-gerbe上では、
非自明慣性は $W$ に $-1$ として作用する。
乗法 $\operatorname{Sym}^2W\to\operatorname{Sym}^2V_{\rm can}$ は $Q$ を与え、
係数Gram形式を $W\otimes V_{\rm can}^*$ 上の形式と解釈できる。
この規約で余核は正確には
$(\omega_C\eta)\otimes\det(W)^{-1}$ であるが、$\det(W)$ は慣性が自明な $K$ 上の線であり、
その自明化による調整は障害を変えない。

一次直交syzygyの空間は、Gram形式で値を双対化すると
\[
 T\simeq W^*\otimes\bigwedge^2V_{\rm can}
\]
となるので、慣性は $T$ にも $-1$ として作用する。
三つの交代形式は自然な写像
$\bigwedge^2T\to\bigwedge^2V_{\rm can}$ を与える。
そのパラメータ平面は
$\mathbb P((\bigwedge^2V_{\rm can})^*)\simeq\mathbb P(V_{\rm can})$ であり、$K$ 上分裂している。
従って $E$ は $K$ 上の平面三次曲線へ降りる。
その上の階数二Pfaffian核にも慣性は $-1$ として作用するので、
射影化はBrauer class $\alpha_{\rm Prym}|_E$ を持つ二次曲線束 $S\to E$ へ降りる。
残余符号商を取ったファイバーの逆構成はこの降下と可換である。
同じ次数二のBrauer classを持つ二次曲線は同型なので、一般点上で
\begin{equation}\label{eq:notes-prym-product-v27}
 S\dashsim_K E\times_KD
\end{equation}
を得る。これは一般ファイバーの双有理モデルであり、束全体の積表示ではない。

\paragraph{非消滅を検出する支配族.}
$F=\mathbb C(s,t)$ 上の二次曲線
$D_0:Z^2-sX^2-tY^2=0$ は非分裂で、独立変数を添加しても非分裂である。
$s$ 進付値を使うと、剰余体上の $Z_0^2=tY_0^2$ から $Z_0=Y_0=0$ となり、
残る $sX^2$ の付値一を打ち消せないことが直接の証明になる。
十八個の独立係数からなる一般の三二次式で $f:\mathbb P^2\to\mathbb P^2$ を作り、
$C=f^{-1}(D_0)$ とする。一般には基点がなく、$C$ は滑らかな四次曲線である。
$D_0$ の幾何学的次数一のPicard classを $\ell$ とし、
$M=f_C^*\ell$、$\eta=M\omega_C^{-1}$ と置く。
$\ell^2=\mathcal O_{D_0}(1)$ から $\eta^2=\mathcal O_C$、
対称表示の分解から $h^0(M)=2\ne3=h^0(\omega_C)$ なので $\eta\ne\mathcal O_C$ である。
自然性により
\[
 \delta_C(\eta)=\delta_{D_0}(\ell)=[D_0]=(s,t)\ne0.
\]
代数閉包上ではこの族は全ての一般二行二列対称二次式表示を動かし、
$\mathcal R_3$ を支配する。従って一般の $\alpha_{\rm Prym}$ は非零で、位数・指数ともに二である。
ここで非分裂性を保つ純超越拡大は $F\subset F(\text{十八係数})$ であり、
$K$ からこの体への拡大が純超越であるとは仮定していない。

\paragraph{曲線へ制限しても分裂しない理由.}
$E$ は分裂平面内の三次曲線なので、次数三の零cycleを持つ。
もし $\alpha_{\rm Prym}|_{K(E)}=0$ なら有理写像 $E\dashrightarrow D$ が存在し、
滑らかな曲線から固有曲線への写像なので全体に延びる。
次数三の零cycleを押し出すと $D$ 上に次数三の因子ができ、
その線束と $\omega_D$ の積は次数一になる。
種数零のRiemann--Rochから $D(K)\ne\varnothing$ となり矛盾する。
従って一般の核束の射影化は、$E$ 上の真の階数二ベクトル束の射影化にはならない。
一方、$S$ の一般点はその二次曲線上の点を与えるので、$\alpha_{\rm Prym}|_{K(S)}=0$ となる。

\subsection{補助曲線の指数三と最小多重切断次数六}
\label{subsec:notes-index-v27}

曲線や固有多様体の指数は、全閉点の剰余体次数の最大公約数とする。
種数一曲線の周期は、そのJacobianに対するtorsor class
$[E]\in H^1(K,\operatorname{Pic}^0(E))$ の位数とする。
一般の $E/K$ は原点を持たず、周期と指数はともに三である。
上界 $\operatorname{ind}(E)\mid3$ は平面直線切断から従う。
下界には、単なる一つの特殊化でなく、次の\emph{支配的な}付値族を使う。
超越元 $r$ と離散付値パラメータ $s$ に対して
\[
 M_s(u,v,w)=\begin{pmatrix}
 ru&sw&sv\\sw&rsv&su\\sv&su&rs^2w
 \end{pmatrix}
\]
という二次曲線netを取ると、判別式とHermite曲線の式は
\begin{align*}
 \det M_s&=-rs^2\left(u^3+sv^3+s^2w^3-\frac{r^3+2}{r}suvw\right),\\
 E_s&:\quad r^2(s^2z_1^3+sz_2^3+z_3^3)+(r^3-4)sz_1z_2z_3=0.
\end{align*}
$t^3=s$ を入れると、$M_s$ は座標変更 $(u,v,w)=(U,V/t,W/t^2)$ と
$\operatorname{diag}(1,t,t^2)$ による合同変換で $M_1$ から得られる。
二つの三次曲線の滑らかさを保証する因子はそれぞれ
$(r^3-1)^2(r^3+8)$ と $(r^3-1)(r^3+8)^2$ であり、$r$ が超越元なら非零である。

支配性の再現可能な判定は次である。$M_1$ の三係数二次式を $(q_0,q_1,q_2)$ とし、
平面座標 $g\in\operatorname{GL}_3$ とnet基底 $h\in\operatorname{GL}_3$ を加えた
\[
 (r,g,h)\longmapsto h\,(q_0(gx),q_1(gx),q_2(gx))^t
\]
を考える。十九変数から十八係数への接行列は $r=2,g=h=I$ で階数十八を持ち、
一つの十八次小行列式は $-462422016$ である。
これにより全順序付きnetの空間、従って一般対称表示を通じて $\mathcal R_3$ を支配する。
$g,h$ の独立な成分を $\mathbf a$ とすれば
$K\hookrightarrow F=\mathbb C(r,\mathbf a)((s))$ であり、$E_F\simeq E_s$ となる。
この支配性は、指数の計算をモジュライの一般点へ戻すために不可欠である。

有限拡大 $L/F$ の次数が三で割れなければ、分岐指数 $e=v_L(s)$ も三で割れない。
三座標の付値を $a,b,c$ とすると、三つの立方項の付値
\[
 2e+3a,\qquad e+3b,\qquad3c
\]
は三を法として相異なる。混合項の付値 $e+a+b+c$ はこの三つの平均なので、
唯一の最小付値より真に大きく、方程式は成立しない。
零の座標があっても同じ付値比較で矛盾する。
従って全閉点の次数は三の倍数である。
一方、$t^3=s$ 上には $[t^{-2}:-t^{-1}:0]$ があるので $\operatorname{ind}(E_F)=3$ である。
基礎体拡大で指数は元の指数を割るため、$\operatorname{ind}(E)=3$ を得る。
種数一曲線の周期は指数を割り、周期一ならJacobianの自明torsorとして有理点を持つので、周期も三である。

式\eqref{eq:notes-prym-product-v27}の両射影から
$\operatorname{lcm}(\operatorname{ind}E,\operatorname{ind}D)=6$ は積の指数を割り、
零cycleの外積から積の指数は $3\cdot2=6$ を割る。従って
\[
 \operatorname{ind}(S)=6.
\]
滑らかな射影曲面の点吹上げは指数を変えないので、同じ値は全ての滑らかな射影双有理モデルに成立する。

さらに、指数六の定義だけに依存せず、任意の非空開部分内に次数六の閉点を作れる。
$U\subset E\times D$ の像の補集合を避ける一般 $K$ 直線で $E$ を切ると、
指数三により一つの次数三の閉点 $x$ を得る。
$L=\kappa(x)$ 上でも $D_L$ は非分裂である。分裂すれば $D$ に次数三の零cycleができるからである。
$D_L$ を $U_x$ の補集合を避ける一般 $L$ 直線で切ると、一つの次数二の閉点 $y$ を得る。
その合成次数は $[\kappa(y):K]=2\cdot3=6$ である。
双有理写像の定義される開部分を通じて、同じ構成を一般ファイバーへ戻せる。
正則な一般ファイバーの閉点は固有な滑らかモデルへのLang--Nishimuraにより指数六に従うので、
全ての多重切断次数は六の倍数であり、最小次数は六である。
この結論は正則な商開部分または滑らかなモデルに関する。
任意に収縮した特異境界に入る切断まで、同じ双有理議論で禁止してはならない。

\subsection{Pl\"ucker二次曲線と内在的二次拡大}
\label{subsec:notes-quadratic-splitting-v28}

一般関数体 $K=\mathbb C(\mathcal R_3)$ 上のweight一のsyzygy空間 $T$ と
$\Omega:\bigwedge^2T\to B$、$\dim T=6$、$\dim B=3$ から
\[
 c_\Omega:\bigwedge^3T\to T\otimes B,\qquad
 uvw\mapsto w\otimes\Omega(u,v)-v\otimes\Omega(u,w)+u\otimes\Omega(v,w)
\]
を作る。一般に階数十八、核 $\Pi$ はweight一の階数二空間である。
従って $D_\Pi=\mathbb P(\Pi)$ はPrym Brauer class
$\alpha_{\rm Prym}$ を持つ二次曲線であり、$D_\alpha$ と $K$ 同型となる。
同型の選択まで標準的であるとは主張しない。

Pl\"ucker関係式の制限は一つの二次式を生成し、
\[
 Y=\operatorname{Gr}(3,T)\cap D_\Pi=\operatorname{Spec}L
\]
は次数二のエタール因子となる。
$D_\Pi$ は非分裂なので $L/K$ は真の二次体であり、$\alpha_{{\rm Prym},L}=0$ である。
これは二つの共通等方三平面を区別する内在的な分裂体である。
一方、$\operatorname{per}(E_L)=\operatorname{ind}(E_L)=3$ は残り、
$S_L\dashsim E_L\times\mathbb P^1$、$\operatorname{ind}(S_L)=3$ となる。

基礎体上で二次代数を計算するには、次数六の中心単純代数 $\operatorname{End}(T)$ の中で
\[
 \mathfrak a_\Omega=\{J:\Omega(Ju,v)+\Omega(u,Jv)=0\}
\]
という一次元空間を解けばよい。
その非零元は $J^2=\delta I$、$\operatorname{Trd}(J)=0$ を満たし、
$L=K[J]=K[t]/(t^2-\delta)$ である。
$rI+sJ$ は $\Omega$ を $r^2-\delta s^2$ 倍するため、
$\operatorname{Res}_{L/K}\mathbb G_m/\mathbb G_m$ は $S\to E$ に忠実に作用する。
その固定部分は $E_L\to E$ という次数二の有限エタール多重切断である。
二次拡大上で核束は次数 $-3$ の二つの線束に分かれる。

試験族 $G(t)$ では
\[
 q_t(a,b)=(t-4)a^2+(t-4)ab-b^2,\qquad \delta=t(t-4).
\]
従って $\sqrt{\det G}=\sqrt{t-4}$ では正しい二次代数を得ない。
$D_\Pi$ を射影直線として座標化できるのはBrauer classの分裂体上である。
その体 $k$ でも二次因子が分裂し得るので、一般式は
$k[t]/(t^2-\operatorname{disc}q)$ と書き、常に体だとは仮定しない。
計算入口は \nolinkurl{isotropic_plucker_v28.py} である。


\paragraph{次の判定に残るもの.}
一般Brauer classの非消滅、$E$ の原点不存在、指数三・六は確定している。
内在的な二次分裂体も確定した。残るのは境界での分岐と相対コンパクト化の降下、全八次元空間の有理性、
および指定反対合自体を忘れたモジュライの幾何である。
元の十一点族の相対potential障害 $\beta$ の消滅は、この議論からは判定されない。


\subsection{補助曲線の節点と使い回せる長さ判定}
固定点 $x_0$ で $M=L(x_0)$、左核 $u=(0,b)$、$w_i=uL_i$ と置く。
全incidenceは $w_iAM^t=0$ を与え、被約な階数低下点の法線微分は
\[
 T_{x_0}\mathbb P^2\xrightarrow{\sim}(\ker M)^*
\]
である。従って $(w_iAw_j^t)$ は非退化であり、補助 Gram 行列式の二次初項を与える。
固定点は常に通常二重点である。

補助曲線の三偏微分に対する次数 $D=3d-5$ の Macaulay 行列で、
$p=3$ では $102$ 次小行列式の付値 $48$、$p=2$ では $205$ 次小行列式の付値 $63$ を得る。
入力・保証精度はともに $101^{256}$ である。
これにより指定次数の Jacobian 商の次元はそれぞれ三、五以下となる。
再利用できる一般補題は
\[
 \dim_k(S/I)_D\le f\le D
 \quad\Longrightarrow\quad
 \operatorname{length}\operatorname{Proj}(S/I)\le f
\]
である。飽和化して一般一次非零因子で割れば、標準次数付き環の Hilbert 関数の差が、
零になるまで正であることから従う。
一点の次数の商次元を、根拠なく射影スキームの長さと同一視しない。

既知の節点がこの長さを尽くすので、他の特異点はない。
節点数が $d-1$ より少ないため、B\'ezout によって幾何学的既約性を得る。
六次では三節点・正規化種数七、八次では五節点・正規化種数十六である。
非固定点はすべて滑らかな通常二重接触点になる。
補助 Gram 行列は節点を含めて corank 一なので、節点曲線上にも
非自明二捩れ線束 $\eta_B$ があり、連結な有限エタール二重被覆を定める。
正規化への引き戻しの非自明性は未決定である。

\subsection{固定三点のmonodromyと残る独立性}
$Z_1$ では固定三点への幾何学的 monodromy の射影が $S_3$ である。
基点の退化イデアルの飽和は $(xy,xz,yz)$ であり、
連結な枠変更群に属する境界安定化元が三巡回置換と互換を与える。
三点の有限エタール被覆が境界まで延びるため、全二十一点が被約な稠密開部分にもこの作用が残る。
ここから $S_3\times\{1\}$ の包含は導けず、非固定側との独立性は別の課題である。

\subsection{補助楕円曲線の変動とGrassmann境界}
公表稿Vの一変数行列 $G(t)$ は
\[
 \det G(t)=t-4,\qquad
 E_t:\ 2tz_1^3+z_2^3+tz_3^3+(t-8)z_1z_2z_3=0
\]
を与える。これは三次曲線の滑らかさの明示証明であるとともに、
次の補助不変量の変動を追うための小さな試験族になる。
\paragraph{試験族の二つの $j$ 不変量.}
Hermite--Pfaffian恒等式により $E_Q$ の基底不変性はnetの幾何から分かる。
上の試験族に対応する判別式曲線は
\[
 \Delta_t:\quad\lambda_0^3+4\lambda_1^3+t\lambda_2^3
 -(t+4)\lambda_0\lambda_1\lambda_2=0.
\]
$t=1$ では $E_t,\Delta_t$ ともに滑らかで、$\det G=-3\ne0$ である。
通常の $j(i)=1728$ の正規化では
\begin{align*}
 j(E_t)&=-\frac{(t-8)^3\bigl((t-8)^3-432t^2\bigr)^3}
 {2t^2(t-2)^3(t+16)^6},\\
 j(\Delta_t)&=\frac{(t+4)^3\bigl((t+4)^3+864t\bigr)^3}
 {4t(t-2)^6(t+16)^3}.
\end{align*}
Hesse形 $X^3+Y^3+Z^3+3hXYZ=0$ の公式
$27h^3(8-h^3)^3/(1+h^3)^3$ から得られる
\cite[\S2]{ArtebaniDolgachevHesse}。
いずれも非定数であり、$\mathcal R_3\dashrightarrow Y_0(2)$ の支配性を確認できる。
十八次元の行列空間上の有理関数の非定数性なので、元の四次Gram曲線が滑らかな
稠密開部分へ制限しても変わらない。
二つの $j$ が満たす次数二のモジュラー方程式は恒等式として再検査できるが、
被覆の存在とtrace零線束の同定は第\ref{subsec:notes-hermite-v27}節の幾何学的証明に依拠する。

\paragraph{三成分閉包と局所モデル.}
一般の $Q$ について、$X_Q$ は滑らかな楕円線織面 $R$ と二つの平面
$\Pi_\pm\simeq\mathbb P^2$ からなる単純正規交差曲面である。
二つの共通等方三平面を $U_\pm$ とすると
$\Pi_\pm=\{U_\pm\oplus\ell:\ell\in\mathbb P(K_3)\}$ である。
射影 $W\to K_3$ の階数二は、Pfaffian三次曲線に直線を強制するため存在しない。
階数一の点はちょうどこの二平面、階数三の点は滑らかな一つの既約開部分となる。

$K_6=U_+\oplus U_-$ に対する交差ブロックを $M(z)$ とすれば、
境界近傍で全方程式は正確に
\[
 M(\ell)d=0
\]
へ消去できる。
$\det M=0$ は滑らかな三次曲線なので、局所式は $fv=0$ となる。
これは成分や交差を点集合だけでなくスキームとして決める小さなモデルである。
$R$ は $E$ 上の分解可能な射影直線束で、二重曲線 $D_\pm$ は互いに交わらない二つの切断となる。
境界平面の $\ell$ と線織面の底の $z$ を同じ座標と仮定しない。

\[
 \deg N_{D_\pm/R}=0,\qquad
 \deg N_{D_\pm/\Pi_\pm}=9,\qquad
 \deg\mathcal T_X^1|_{D_\pm}=9.
\]
一般一次摂動は種数二十K3曲面を与え、その全空間には二つの二重曲線上に九個ずつ普通二重点が生じる。
この九点は一次平滑化切断の単純零点である。
二つの平面成分のイデアルを吹き上げる局所式は $ab=tc$ の $(a,t)$ のblow-upであり、
射影的な小さなcrepant解消を与える。
中心ファイバーは、中央の $R$ を保ち、両側の平面をそれぞれ九点で吹き上げた三成分のType II Kulikovモデルとなる。
法束は平滑化零因子を引くことで相殺する。
九点を吹き上げた各平面を、追加条件なしに有理楕円曲面とは呼ばない。

\paragraph{スクロールの分極と再利用可能な数値.}
$\det\mathcal K=\mathcal O_E(-2)$ なので、
直線部分空間をパラメータ化する $\mathbb P_E(\mathcal K)$ の超平面類の平方は六である。
GrassmannのPl\"ucker類については $\deg X_Q=38$、両平面の次数は一ずつ、
主成分 $R$ の次数は三十六である。
これら二つの分極を同一視しない。
一般関数体上で二平面を区別するには、前述の二次拡大 $L/K$ が必要であり、
Kulikov修正の粗モジュライ上への降下とBrauer classの境界での分岐は残る問題である。


\section{行列式線束が結ぶ線織面・K3端成分・被覆}
\label{sec:notes-new-structures-v29}
以下は公表稿V・VIの線束と被覆の証明から、再利用する構成を取り出したものである\cite{PaperV,PaperVI}。

\subsection{二つの線織面を区別する線束}
順序付き共通等方三平面に関して
$\Omega=\left(\begin{smallmatrix}0&M\\-M^t&0\end{smallmatrix}\right)$ と書き、
$E=\{\det M=0\}$、$H=\mathcal O_E(1)$ とする。
\[
 0\to\mathcal O_{\mathbb P^2}(-2)^3\xrightarrow{M}
 \mathcal O_{\mathbb P^2}(-1)^3\to i_*\theta\to0
\]
で定める $\theta$ は次数零の非自明線束である。
netの判別三次曲線上の二位数線束 $\theta_N$ とは異なる。
転置は $\theta^{-1}$ を与え、核束と二つの線織面は
\[
 \mathcal K=H^{-1}(\theta^{-1}\oplus\theta),\qquad
 R=\mathbb P_E(\mathcal O\oplus\theta),\qquad
 S=\mathbb P_E(\mathcal O\oplus\theta^2)
\]
である。$R\to S$ は残余符号 $\tau$ による大域商であり、各ファイバーでは二乗写像である。
二つの二重曲線の法束は $\theta$ と $\theta^{-1}$ となる。

境界側の座標 $x=\ell$ とスクロールの底の座標 $y=z$ の差は平行移動である。
$d\in\ker M(x)$、$u^tM_i d=0$、$u\in\ker M(y)^t$ に対して
\[
 0\to\mathcal O_E(x)\to\operatorname{im}(M:\mathcal O_E^3\to H^{\oplus3})
 \to H(-y)\to0
\]
を作り、行列式を取ると $\theta=\mathcal O_E(y-x)$ となる。
従って $p|_{D_+}=t_\theta$、$p|_{D_-}=t_{\theta^{-1}}$ である。
この差を無視して二つの座標を同一視しない。
平行移動はPicard零次部分に恒等的に作用するので、法束の式自体はどちらの座標でも一致する。

二共通平面が被約な二点である開部分では $\theta^2\ne\mathcal O_E$ である。
二位数なら $M$ を対称化でき、グラフ $(u,au)$ が共通等方三平面の一次元族となるためである。
ここには、共通平面の二次被覆の分岐と行列式線束の二位数条件を比較する入口がある。
その境界での重複構造は未計算である。

\subsection{第二の楕円moduliを保証する有限入力}
十八次元のブロック対称Gram空間 $\mathcal H$ に対し、
\[
 \mathcal H^\times\times\operatorname{GL}_6
 \longrightarrow(\bigwedge^2 k^{6*})^{\oplus3}
\]
は支配的である。既存Hesse族の $t=1$ の点で、四十五次微分小行列式は
\[
 \frac{962948}{1162261467}
 =\frac{2^2\,7^2\,17^3}{3^{19}}\ne0
\]
となる。出力三形式の基底変更を追加する必要はない。
\nolinkurl{ruled_moduli_dominance_v29.py} は、逆行列の微分と合同作用の微分から行列を作り直し、
九倍した整数行列のBareiss消去でこの非消滅を検証する。

一次行列式表示と線束の古典的対応を合わせると、
逆Gram族は $(E,\{\theta,\theta^{-1}\})$ の二次元モジュライを支配する。
$E$ の同型類は一次元のnetモジュライだけで決まるため、
一般netを固定しても $\theta$ のKummer像は非定数となる。
従って $\theta$ は幾何学的一般点で無限位数、複素数上では非常に一般の点で無限位数である。
各 $n$ に対する $\theta^n=\mathcal O_E$ の真閉部分を可算個除くという意味であり、
単一の非空Zariski開集合の全点に対して非捩れ性を述べるものではない。

\subsection{反標準束の半豊富性とHalphen指数}
小解消後の端成分 $Y=\operatorname{Bl}_9\mathbb P^2$ では
$D=-K_Y$、$D^2=0$、$\mathcal O_D(D)=\theta^{\pm1}$ である。
$D$ はnefであり、完全列
\[
 0\to\mathcal O_Y((n-1)D)\to\mathcal O_Y(nD)
 \to\theta^{\pm n}\to0
\]
が全ての反標準線形系を制御する。
無限位数なら $h^0(nD)=1$ が全 $n\ge1$ で成立し、$\kappa(D)=0$、$\nu(D)=1$ である。
位数が $m$ なら、$j<m$ に対するRiemann--Rochと $H^2(jD)=0$ から $H^1(jD)=0$ となり、
$h^0(mD)=2$ を得る。新しい切断は $D$ 上で非零なので、$|mD|$ は基点を持たない。
$D$ の原始性と $m$ の最小性から一般ファイバーは連結で、種数一となる。
従って半豊富性は有限位数性と同値で、Halphen指数は $m$ である。
良い二平面開部分の有限位数は $m\ge3$ である。第30版では十分大きい全 $m$ の実現を得た。
低位数を含む全 $m\ge3$ の実現と有効な下限は未決である。

\subsection{二つの二次指標を混同しないための被覆図式}
$F=\mathbb C(Y_0(2))$ 上のnet代表 $N_u$ を固定し、まず $k=\overline F$ 上で考える。
$X=\mathbb P^5$ を二次曲線の空間、$G=H\rtimes C_2$、$H\simeq(\mathbb Z/3)^2$、
$M=X/G$ とする。$J=\operatorname{Pic}^0(E_N)$ は固定した楕円曲線である。
$Y\to X$ は共通等方三平面の一方を選ぶ被覆、$\iota$ はその交換とする。
行列式線束は支配的写像 $\theta:Y\dashrightarrow J$ を定め、
\[
 \theta(\iota y)=-\theta(y),\qquad
 \theta(gy)=\chi(g)\theta(y)
\]
を満たす。$H$ は $E_N$ に平行移動として作用し、Picard零次部分への作用は自明である。
$G$ の残りの対合は反転を与える。

$Y$ が分裂すると各成分は $\mathbb P^5$ と双有理となり、
楕円曲線への写像が定数になるため、固定net上の非定数性に反する。
従って $Y$ は幾何学的に既約である。
$W=Y/H$ とし、$G/H$ の対合を $\rho$ とすると、$W\to M$ の群は $C_2^2$ である。
\begin{center}\small
\begin{tabular}{c|c|c}
中間被覆&関数体&固定する対合\\\hline
$M_\chi=X/H$&$K_0(\sqrt\epsilon)$&$\iota$\\
$M_L=Y/G$&$K_0(\sqrt\delta)$&$\rho$\\
$Z$&$K_0(\sqrt{\delta\epsilon})$&$\rho\iota$
\end{tabular}
\end{center}
ここで $K_0=k(M)$、$\delta$ は元の共通平面被覆、$\epsilon$ はnet安定化群の符号である。
二指標は独立であり、合成体の次数は四である。
二つの反転を相殺した $Z$ だけに、この構成による固定 $J$ への写像が降りる。
\[
 Z\dashsim\operatorname{Norm}(M\times_{J/\{\pm1\}}J),
 \qquad q(\overline Z)=1.
\]
従って幾何学的一般netファイバーの $Z$ は非単有理、特に非安定有理である。
元の $M_L$ の有理性や、全六次元空間の不正則数・有理性へこの結論を移してはいけない。
一般理論によるAbel多様体への写像の延長を用いる\cite[I, \S3]{MilneAbelianVarieties2008}。

この構成は固定した $N_u/F$ に対して $F$ 上へ降りる。
gerbe上の余核線束と $\det U_+$ はともに重み一なので、後者の逆をテンソルして重みを相殺できる。
相対Picard classは基底からの線束のテンソル積を忘れる。
これにより、任意基底上でgerbe全体が中立であると仮定せずに写像を降ろせる。
一方、$T\oplus B$ 上の慣性作用は $(-1,+1)$ であり、元のGrassmann曲面には自明でない。
共役な二平面の和が不変であることだけで、Kulikov模型の粗基底上への降下を結論しない。

分裂Hesse座標では符号被覆は
\[
 w^2=\Delta(a,b),\qquad
 \Delta=a^2b^2-4b^3-4a^3-27+18ab
\]
である。$p=x^3+y^3+z^3$、$q=xyz$、$r=x^3y^3+y^3z^3+z^3x^3$ として
$a=p/q$、$b=r/q^2$、$w=(x^3-y^3)(y^3-z^3)(z^3-x^3)/q^3$ を用いる。
第三の被覆の平方類は $\delta\Delta$ となる。
$F$ 上で特定Hesse座標を仮定せず、次数九の一次元半不変式空間を降ろすことが必要である。

\section{種数二十K3の偏極付き境界格子}
\label{sec:notes-boundary-lattice-v29}
端成分の基底を $H_\pm,E_i^\pm$、反標準類を $d_\pm=3H_\pm-\sum E_i^\pm$ とする。
中央成分の数値基底は $S,F$、$S^2=F^2=0$、$SF=1$ である。
正規化の格子と垂直部分格子は
\[
 \Lambda=I_{1,9}\oplus U\oplus I_{1,9},\qquad
 J=\langle(d_+,-S,0),(0,S,-d_-)\rangle.
\]
$J$ は原始等方二平面であり、
$L=J^\perp/J\simeq U\oplus E_8(-1)^{\oplus2}$ である。
明示的な双曲基底は
\[
 e=(E_9^+,S+F,E_9^-),\qquad f=(d_+,0,0)
\]
で、残りは各端の $E_i-E_{i+1}$ $(i\le7)$ と $H-E_1-E_2-E_3$ である。
Mayer--Vietorisと特殊化から整数上の等長写像を作り、
この明示格子のユニモジュラー性を用いると、極限の $I^\perp/I$ との同定に追加指数は生じない。

Pl\"ucker類は
\[
 h=(H_+,6S+3F,H_-)\pmod J,\qquad h^2=38,
 \qquad h\alpha_8^+=1.
\]
$L_h=h^\perp$ は階数十七の負定値偶格子、判別群は $\mathbb Z/38$、
判別形式は $-1/38$ である。
全ての平方 $-2$ の根は各端の $E_i-E_j$ で尽くされ、根系は $A_8\oplus A_8$、根数は百四十四である。
根格子 $A$ の飽和指数は三で、原始な直交ベクトルは
\[
 v=(3H_+,-20S+9F,3H_-),\qquad v^2=-342.
\]
$[L_h:A\oplus\mathbb Zv]=27$、$[L_h:\overline A\oplus\mathbb Zv]=9$ である。
根系のみから全境界格子を同定せず、これらの貼合せと整数Gram行列を保持する。

全根を排除する短い証明も再利用できる。
$x=(aH-\sum bE,cS+dF,a'H-\sum b'E)\in h^\perp\cap J^\perp$ に対して
$\sum b=3a-d$、$\sum b'=3a'-d$、$a+a'+3c+6d=0$ となり、
Cauchyの不等式は $-x^2\ge38d^2/9$ を与える。
根なら $d=0$。$a,a'$ を三を法として正規化すると、
$(1,2),(2,1)$ の場合は負ノルムが少なくとも四となるので、両者は零である。
残る平方和二の整数解が各端の差だけを与える。

二重曲線は二本なので、モノドロミーは原始等方基底 $\delta,\mu$ に対して
\[
 N(x)=2\bigl((x\mu)\delta-(x\delta)\mu\bigr),\qquad
 N^2=0,\quad \operatorname{rank}N=2,\quad \operatorname{Im}N=2I
\]
となる。最大整数因子は二、群の指数 $[I:\operatorname{Im}N]$ は四である。
一般K3のPicard階数や、端成分上の有効根の本数をこれから推定しない。
\nolinkurl{kulikov_lattice_v29.py} は全格子と独立な有限根列挙を検算する。

\section{標準七次被覆を次の計算へ使う}
\label{sec:notes-septic-v29}
本節の主張と完全なファイバーの検査は公表稿VIIに証明した\cite{PaperVII}。
六次側の標準十八点 $D$ に対し、$t$ の係数行列が可逆な開部分では
$Qg=(\det Q)t$、$g=\operatorname{adj}(Q)t$ から
$\widetilde{(g_0,g_1,g_2)}=\mathcal I_D$ を得る。
従って $\operatorname{Bl}_D\mathbb P^2$ 上の完全線形系
\[
 L=5H-\sum_{1}^{18}E_i,\qquad L^2=7
\]
は一般に有限平坦な次数七の射を与える。$L|_C=\eta(2H)$、その次数は十二、切断数は三である。
源側のramification因子は $R_f\sim2\widetilde C$、
押し出しcycleの次数は二十四である。
線形同値 $R_f\sim2\widetilde C$ はスキームとしての等式ではない。
第30版ではさらに、一般の $R_f$ が滑らかな種数三十七の整曲線で、
被約分岐曲線が整な二十四次曲線、$R_f$ がその正規化であることを証明した。

同じ成分に属する厳密な単純foldの点は、十九変数標準形で
\[
 s=1,\qquad (a,c,h,j,n,q)=(3,-1,4,-2,5/4,3/4),
 \qquad\text{他の変数は零}
\]
である。$A_0$ を通る六次元の対角変化族の中にあり、成分所属を保っている。
$t=x$、$Q=\operatorname{diag}(F_0,F_1,F_2)+2xx^t$、
\[
 (F_0,F_1,F_2)^t=
 \begin{pmatrix}-1&3&-1\\4&-1&-2\\5/4&3/4&-1\end{pmatrix}x^{[2]},
 \qquad g=(x_0F_1F_2,x_1F_0F_2,x_2F_0F_1)
\]
となる。$Q$ 自体を対角行列と取り違えない。
基点は $F_i=F_j=0$ の十二点と $x_i=F_i=0$ の六点である。
六次曲線は三チャート全てで $\mathbb F_{101}$ 上滑らかであり、標数零でも滑らかである。

標的 $[1:1:1]$ のファイバーは、代表元を正規化した $X=MX^{[2]}$ の
原点以外の解と同一である。
消去多項式は $z(z-1)^2q_5(z)$、$q_5$ は次数五でsquare-freeであり、
原点の単純解を除いた長さは七である。
各座標のgcd条件が座標零の非零解を排除し、同時に六本の例外曲線の像が標的を通らないことを保証する。
残る十二本の例外曲線の像は標的座標直線である。
従って二重点一個と単純点五個を漏れなく数えている。
二重点の左右Jacobi核による二次微分は $-679/1521$ で非零となる。

相対blow-upと有限平坦階数七の局所分解により、この単純foldは
パラメータの非空開集合にも残る。
幾何学的monodromyは推移的な素数次数七の群なので原始的であり、
単純foldの互換と、その共役を辺とする連結グラフから $S_7$ を得る。
この $S_7$ は標準七次写像の群であって、元の十八非固定点の群ではない。
\nolinkurl{canonical_septic_cover_v29.py} が全ての有限入力を読取り専用で再計算する。

\paragraph{分岐計算に使える六変数族の式.}
一般の $M=\left(\begin{smallmatrix}-1&a&c\\h&-1&j\\n&q&-1\end{smallmatrix}\right)$ でも
$F=Mx^{[2]}$、$g_i=x_iF_jF_k$ と書ける。
全成分が非零の標的 $[u_0:u_1:u_2]$ に対して、
正規化されたファイバーの方程式は
\[
 X=\operatorname{diag}(u_0,u_1,u_2)M X^{[2]}
\]
となる。$\det M\ne0$ なら斉次化した三つの二次式には無限遠解がなく、
完全交叉の長さは八である。原点が単純なので残余長は七となる。
ただし座標零の非零解は例外曲線上の点に対応し得るため、
トーラス内の消去で捨ててはならない。

同じ族にはJacobian恒等式
\[
 \det\left(\frac{\partial g_i}{\partial x_j}\right)
 =-5F_0F_1F_2
 \det\left(\operatorname{diag}(F)
       -2\operatorname{diag}(x)M\operatorname{diag}(x)\right)
\]
がある。$h_i=x_i/F_i$、$P=F_0F_1F_2$ とすると $g=Ph$、
Eulerの式は $(dh)x=-h$、$(dP)x=6P$ を与える。
行列式補題で $\det(dg)=-5P^3\det(dh)$ となり、上の式が従う。
三つの二次曲線 $F_i=0$ は標的の三座標点へ縮約される。
残る六次因子は、分岐源の整性と滑らかさを証明する特殊化の中心成分となる。
特殊ファイバーの分岐像の正規化種数は十であり、一般ファイバーの種数三十七と混同しない。

\paragraph{九余因子式の大域的消去に残る具体的構造.}
$A$ はLaurent単項式 $1,x/y,x/z,y/x,y/z,z/x,z/y$ の積の線形汎関数を並べたmoment行列と見られる。
$A^{-1}$ の九零条件は、添字 $1,2,5,6,4,3$ を巡回順に置くと、
六頂点上の巡回三重対角形と中心頂点 $0$ を持つ疎行列条件となる。
$A=\left(\begin{smallmatrix}s&v^t\\v&C\end{smallmatrix}\right)$ と分けると、
逆行列の下右ブロックは $(C-vv^t/s)^{-1}$ である。
従って $s\ne0$ の開部分では、巡回三重対角行列 $T$ を用いて
$C=T^{-1}+vv^t/s$ と書き、十九変数の線形条件を課し直す方法がある。
これは方程式の別表示であり、有理性や全成分の既約性を証明するものではない。
六次像の大域有理性を進めるなら、まずこの線形条件の消去が有限多価になる箇所を特定する。

\section{境界periodとHalphen層の継続入口}

\paragraph{到達点と使うべき座標。}
両平面の零点 $q_i^\pm$ を使えば、底曲線への平行移動を全て吸収した
periodは $\psi(x)=-\sum b_iq_i^++\sum b_i'q_i^-+(2d-c)a$ となる。
$a=\sum q_i^+=-\sum q_i^-$ である。
固定 $E$ の配置空間は $\mathcal P_E\simeq E^{17}$、
偏極periodへの写像は核 $\Delta E[3]$ の次数九のisogenyである。
整数計算をするときは $L_h$ ではなく $L/\mathbb Zh$ を最初に使う。
$L_h\hookrightarrow L/\mathbb Zh$ は指数三十八なので、
そこへ制限した係数行列の指数は百十四になる。
根系の直交方向は $v=3h-38f$ であり、別の独立なねじれ方向ではない。

\paragraph{次の優先問題。}
境界periodの支配性は証明済みとなったので、次の有効な目標は、
実際の逆Gram族の像が各 $m$-torsion層と交わるかの有効判定である。
現状では、ある固定した一般の $E_0$ 上で有限個の $\theta$ を除き実現するため、
十分大きい全 $m$ のHalphen指数は実現する。
全 $m\ge3$ を得るには、像の有限補集合と低位数torsionとの位置関係を
明示的に制御する必要がある。単なる支配性だけで全位数の実現を主張しない。
固定したnetの第三の二次被覆と $\theta$-Kummer写像は第30版で明示された。
従って $R$ の有限個の悪い値、三次pencilの基底、および低位数torsionの逆像を
同時に調べることが、有限個の候補を処理する具体的な手段になる。

\paragraph{一般性を減らす補題。}
滑らかな平面三次曲線を行列式に持つ任意の $3\times3$ 一次行列 $M$ を取り、
その線束を $\theta$ とする。$\theta^2\ne\mathcal O_E$ なら、
cross-block交代形式の共通等方三平面のスキームは
$U_+,U_-$ の二つの被約点だけである。

証明は以下の通り。任意の非零列 $v$ に対し
$\operatorname{span}_i M_i v$ の次元は少なくとも二である。
さもなければ基底を変えた行列の一列が $(\ell,0,0)^t$ となり、
$\det M$ が一次式と二次式に分解する。同様に任意の非零行 $u^t$ に対して
$\operatorname{span}_i u^tM_i$ の次元も少なくとも二である。
共通等方三平面 $W$ が $U_+$ と非零交差すれば、その一つの非零元 $u$ から
$\operatorname{pr}_{U_-}W$ は高々一次元となる。
その次元が一なら $W\cap U_+$ は二次元であり、
射影の非零元 $v$ に対する $\operatorname{span}_iM_i v$ の消滅空間に
含まれることと矛盾する。従って $W=U_+$ である。
残る場合は $U_-$ 上のグラフ $W=\{(Av,v)\}$ であり、等方条件は
$A^tM=M^tA$ である。この可換図式は
$\theta\to\theta^{-1}$ を誘導するが、$\theta^2\ne\mathcal O_E$ なので零写像である。
従って $A^t:\mathcal O(-1)^3\to\mathcal O(-1)^3$ は
$M^t$ の像を通る。しかし $\operatorname{Hom}(\mathcal O(-1),\mathcal O(-2))=0$
なので $A=0$ となる。接空間も同じ一次方程式であるため両点は被約である。

この補題は、任意の $(E,\theta)$ から作る拡大Grassmann族を調べる際に使える。
ただし、これだけでは制約付き逆Gram像への属することまでは保証しない。

\paragraph{将来の跳躍層。}
固定 $E$ のmarked period空間で、任意の非零
$\lambda\in L/\mathbb Zh$ の貼合せ可能性は $\psi(\lambda)=0$ という
余次元一の条件である。$\lambda$ が原始的なら核は連結であり、
$\lambda=n\lambda_0$ の場合は $n^2$ 個の平行移動された成分を持つ。
従って点の衝突以外のCartier類の跳躍も、十八点への整数一次条件として
具体的に書ける。ただし、中央ファイバーでの貼合せ可能性だけから
任意の一変数平滑化の一般ファイバーでのNoether--Lefschetz条件を導かない。
平滑化方向を同時に制約する議論が次に必要になる。

\paragraph{一次資料。}
Kond\=o, \emph{Type II degenerations of K3 surfaces}, Nagoya Math. J.
\textbf{99} (1985), 11--30, \S2.8--2.9 は、制限束の比をperiodと定義し、
Carlson拡大と同一視している。Beauville,
\emph{Determinantal hypersurfaces}, Michigan Math. J. \textbf{48} (2000),
Propositions~3.1, 3.3 は行列式と非自明次数零線束の対応である。
一般のType~II period理論自体を本研究の新規性と数えない。

\section{三次pencil・二次同種写像・Albaneseの継続入口}
\label{sec:notes-modular-v30}

\paragraph{今回確定した構造.}
幾何学的一般netを固定した二次曲線係数空間を、双対二次曲線の係数による
$X=\mathbb P^5$ として扱う。$C=a_*c_*^t+c_*a_*^t-2b_*b_*^t$ について
\[
 \det G=I_m(C),\quad\det C=2(\det L)^2,\quad
 R=2I_m(C)/\det C
\]
である。$I_m$ は六係数の三次式。従って $R$ のファイバーは四次元三次超曲面となる。
元の $D\subset\mathbb P(N)$ の係数と $D^\vee$ の係数を混同しない。
三次になるのは $D^\vee$ の係数であり、元の係数への双対化は二次Cremona変換である。

$q=(m-2)^3$、$a_m=m^2-40m-32$ と置くと
\[
 B_m:\ v^2=-R(4R^2-a_mR-8q),\qquad
 \kappa(R)=-4R+a_m+8q/R.
\]
$B_m$ は幾何学的に $\operatorname{Jac}(\Delta_N)$ であり、
$\kappa$ は $B_m\to\operatorname{Jac}(E_N)$ の二次同種写像のKummer商である。
従って、二次曲線の一般変動は、行列式線束のKummer座標へ直接ではなく、
もう一つの楕円曲線のKummer座標 $R$ を経由している。

$Y,W,Z$ は全て $B_m$ へ降りる。各写像の幾何学的一般ファイバーは整かつ単有理で、
$\operatorname{Alb}(Y)=\operatorname{Alb}(W)=\operatorname{Alb}(Z)=B_m$、不正則数はちょうど一である。
$\theta_Z:Z\dashrightarrow\operatorname{Jac}(E_N)$ のStein有限部分は次数二なので、
その幾何学的一般ファイバーはちょうど二成分になる。
元の共通平面被覆 $M_L=Y/G$ の不正則数は零である。
この差は $G/H$ が $B_m$ 上で反転し、第三被覆では $\rho\iota$ が打ち消し合うことによる。

\paragraph{有限二重被覆の明示モデル.}
$D=\det C$、$I=I_m(C)$、$w=vD^2$ と置けば
\[
 w^2=-4ID(8I^2-a_mID-4qD^2)
\]
となる。これは四つの三次pencil memberを分岐因子とする二重被覆である。
分岐因子の総次数は十二だが、四memberは共通の $(3,3)$ 基底部分集合を持つ。
従って滑らかな十二次分岐被覆との比較では、この共通基底による特異性を明示的に扱う必要がある。
このモデルの特異性と解消は、Albanese構造を方程式から説明する次の入口となる。

\paragraph{計算を再開するときの最短入口.}
$A=\Omega_{23}^{-1}\Omega_{31}$、$B=\Omega_{23}^{-1}\Omega_{12}$、$K=[A,B]$ と置く。
次の二式が主要な入力である。
\[
 \operatorname{tr}K^2=\frac{-3R+(m-20)(m-2)+6q/R}{m^2},\qquad
 \mathscr J=K^3-\tfrac14\operatorname{tr}(K^2)K.
\]
$\mathscr J$ は共通等方平面上で $\pm\ell$ として作用し、
\[
 \ell^2=-R(4R^2-a_mR-8q)\frac{(R^2+2q)^2}{64m^6R^4}.
\]
この普遍生成元は、毎回 $45\times36$ 線形系の核を求める手続きを置き換える。
ただし $\ell=0$ の特殊点では生成元が消えるため、共通平面の特殊化をこの式だけで判定しない。

trace恒等式のexact gateは、五次元conicの三角sliceとnet変数 $m$ の六変数上で
多項式全係数を照合する。平方式はMoore標準形の三つの固有値の差と
Cayley--Hamilton恒等式から導く。
一般六変数行列で $\mathscr J^2$ をそのまま展開する必要はない。

\paragraph{ファイバーについて既に十分な証明.}
三次pencilには $\det C=0$ が含まれるので、一般memberの幾何学的整性と
非錐性は直ちに従う。Koll\'arの三次超曲面の単有理性定理を適用すれば、
一般memberの滑らかさや特異軌跡の完全分類を先に求める必要はない。
有限群商のファイバーにも単有理性が残るため、AlbaneseとStein分解はここまでで確定する。

\paragraph{次に価値のある問題.}
明示的な $m=1,R=1$ のmemberでは、Jacobianイデアルの有理数体上の
Gr\"obner計算が射影的特異スキームのHilbert多項式 $6n$ を与える。
従ってこのmemberは滑らかではなく、一次元の特異部分を持つ。
この入力は \nolinkurl{notes/cubic_jacobian_QQ_certificate_v30.json} に保持する。
一つのmemberから一般の特異軌跡の次元・整性・滑らかさまでは結論しない。
次に調べるなら、特異軌跡を一般の $(m,R)$ で同定し、群作用と整合する射影構成を探す。

四次元の三次超曲面の特異構造を求めれば、その幾何学的有理性を強められる可能性がある。
しかし本体 $Z\to B_m$ のファイバーはさらに有限商なので、三次超曲面の有理性だけでは
$Z\dashsim B_m\times\mathbb P^4$ は従わない。
本体の双有理型を進めるには、群作用に適合する有理パラメータ化または
$k(B_m)$ 上の有理性を扱う必要がある。
元のnetモジュライ関数体 $F$ 上で $B_m$ の捻りを同定する問題も残る。
本節の $m$ は有限なHesse型被覆上の座標であり、これを $F$ の座標と同一視しない。

\section{有限七次被覆と分岐曲線の継続入口}
\label{sec:notes-septic-new-v30}

\paragraph{一般点で確定した幾何.}
選択した十次元余因子成分 $Y_0$ の非空開部分で
\[
 f:X=\operatorname{Bl}_{18}\mathbb P^2\longrightarrow\mathbb P^2,
 \qquad L=5H-\sum E_i
\]
は有限平坦な次数七の射であり、$L$ は豊富である。
一般の幾何学的monodromyは $S_7$ である。
分岐源 $R$ は滑らかな幾何学的整曲線で、
\[
 R\sim12H-2\sum E_i,\quad R^2=72,\quad K_XR=0,\quad g(R)=37.
\]
被約分岐曲線 $B$ は整な二十四次曲線で、$\nu:R\to B$ は正規化となる。
従って $p_a(B)=253$、$\sum\delta_p(B)=216$ である。
微分 $d\nu$ の零因子の長さは七十二、双対曲線 $B^*$ の次数は四十八である。
公表稿VIIではこれらを一般の同じ成分上で証明する\cite{PaperVII}。

\paragraph{有限性の証明に再利用できる判定.}
$h^0(\eta(H))=0$、$h^0(\eta(2H))=3$ と後者の基点自由性から
$h^0(\eta(H+p))=0$ を得る。従って十八点のうち任意の十七点を通る四次曲線はない。
$L$ が縮約する既約曲線を $dH-\sum m_iE_i$ とすると
\[
 \sum m_i=5d,\qquad \sum m_i^2\le d^2+2d+2.
\]
Cauchyから $d\le6$、整数の最小平方和から $d=5,6$ を除き、
上の十七点条件から $d=4$ も除く。
残るのは五点直線、十点二次曲線、十五点三次曲線、
または一点二重と十三点単純の三次曲線である。
同じ成分上の六変数特殊族では、縮約曲線は三本の $F_i=0$ だけである。
十八点を標識した有限エタール被覆上で三つの効果的条件が真閉集合となることを
一次変形で示せば、一般点の有限性が従う。
標識被覆の全成分を一つと見なさず、証人を通る既約成分を選んでから開集合を降ろす。

\paragraph{分岐源の整性を証明する特殊化.}
単独foldの証人 $A_*=(3,-1,4,-2,5/4,3/4)$ では、残余critical六次式は二次式と四次式の積に分解する。
この点を滑らかな残余曲線の証人として再利用してはならない。
六変数族の点 $(a,c,h,j,n,q)=(1,2,3,4,5,6)$ では
\[
 R_*=R_0+T_0+T_1+T_2.
\]
$R_0$ は滑らかな種数十の六次曲線のstrict transform、
$T_i$ は互いに交わらない収縮二次曲線のstrict transformである。
全成分の係数は一、$LR_0=24$、$LT_i=0$ となる。
一般点では $L$ が豊富なので、一般分岐源の各成分の特殊化には
唯一の正次数成分 $R_0$ が必要になる。その係数が一だから一般分岐源は幾何学的に整である。
単独foldの持続と合わせ、分岐源から被約分岐曲線への次数が一になる。
この特殊族自体の分岐曲線の正規化種数十を、一般点の種数と同一視しない。

\paragraph{一次で見えない節点を平滑化する二次補正.}
$R_*$ は三十の通常節点を持ち、そのうち二十四点はblow-up基点の外にあり、
六点は例外曲線上にある。前者では一次平滑化が検出できる。
後者では動く基点を二次まで追い、blow-up座標で実際の分岐式を取る必要がある。
局所分岐式の二次部分を
\[
 r(z,\varepsilon)=\tfrac12z^tHz+\varepsilon c^tz+\varepsilon^2b+
 \text{三次以上},\qquad \det H\ne0
\]
と書くと、真の平滑化係数は $b$ 単独ではなく
\[
 b-\tfrac12c^tH^{-1}c
\]
である。これは臨界点 $z=-\varepsilon H^{-1}c+\cdots$ を代入した値である。
六点でこの値は $\mathbb F_{163}$ 上
$118,118,14,14,31,31$ となり、全て非零である。
動く基点に対し $h=g/u$ として $\det[h;h_u;h_w]$ を計算する方法と、
斉次Jacobianを $5u^2$ で割る方法が独立に一致する。
係数の $\varepsilon^2$ 項を二階微分そのものと取り違えない。

標数零へ戻すには、孤立した有限体の飽和計算だけでは不十分である。
三十節点の交差族が固有かつ有限エタールとなる局所環を取り、
例外上六点を閉開部分として区別する。
単元Hessianと非零一次・二次係数を持ち上げ、相対特異点スキームの固有性を使えば、
有限個の真の判別条件を避けた一般点で $R$ が滑らかとなる。

\paragraph{七十二と四十八を特異点型なしで使う.}
$R$ が滑らかなので $df$ は $R$ 上で階数一であり、
$\operatorname{coker}(df)=i_*\mathcal M$ で $\mathcal M$ は線束となる。
曲面Riemann--Rochから
\[
 \chi(T_X)=-28,\quad\chi(f^*T_{\mathbb P^2})=8,
 \quad\deg\mathcal M=72.
\]
$\mathcal I=\operatorname{im}(df|_R)$ の次数は $3LR-72=0$ なので、
$T_R\to\mathcal I$ の零因子 $Z_\nu$ は次数七十二である。
Gauss写像の線束は
$K_R\otimes\nu^*\mathcal O(2)(-Z_\nu)$ で、次数は四十八となる。
標数零の双対性によりGauss写像は双有理である。

尖点七十二個・節点百四十四個という個数は、分岐曲線が通常の節点と尖点だけを持つ場合の帰結である。
次の有効な検証対象は、$Z_\nu$ の被約性、分岐源の二点を同じ像へ送る対応の横断性、
三重像と非通常接触の排除である。これらを七次元の制約付き族の中で証明する必要があり、
無制約な一般の曲面射の安定性をそのまま引用して済ませない。


\section{相対 critical 問題の最小限のデータ}

$B=\kk[[t_1,\ldots,t_{16}]]$、$T=B[[x,y]]/(F_1,F_2)$ は実際の残留十一点族である。
$T$ は正則局所整域、$B$ 上有限自由階数十一で、一般点ではエタールである。
\[
 P=B[[x,y]]/(F_1,F_2)^2,\quad
 D:P\longrightarrow T^2,\quad
 Q=\Omega^1_{T/B}/dT
\]
と置く。$P$ は階数三十三で、$D$ は $22\times33$ 行列である。
中心での階数十九の単元ブロックを消すと
\[
 D\sim I_{19}\oplus[\,0\ b\ C\,],\qquad C\in\operatorname{Mat}_{3\times12}(B)
\]
となる。$M=\coker C$、$\beta=[b]\in M$ なら $Q=M/B\beta$ である。

\subsection{存在判定として保持すべき同値}
正規化した全相対 potential の存在は $\beta=0$ と同値である。
中心 class は四次必要条件によって一つの $[W_9]$ に固定される。
形式解が存在すれば、代数的冪級数環の忠実平坦性により、代数的で収束する解を選べる。
従って収束性を別個の未解決問題として数えない。

分岐加群
\[
 \mathscr R=\ker(d:T\to\Omega^1_{T/B})
\]
は $\ker D$ と同型であり、階数十一の反射的 $B$-加群である。
$Q$ は $B$-加群であるが、$dT$ は一般に $T$-部分加群ではないため、
$Q$ を $T$-加群と見なしてはならない。
また $Q$ 全体は一次 de Rham cohomology ではない。
後述の $Q\to\Omega^2$ の核が一次 cohomology である。

\subsection{三方向ファイバーと一次近傍}
\[
 A_0=\kk[[t_i:i\notin\{5,15,16\}]],\quad
 \mathfrak a=\mathfrak m_{A_0}B,\quad
 s=t_{15},\quad u=t_5-11t_{15},\quad v=t_{16}.
\]
このとき
\[
 M/\mathfrak aM\simeq Q/\mathfrak aQ
 \simeq\Hom_\kk\left(
 \frac{\kk[[s,u,v]]}{(s^3,su,sv,u^2,uv,v^2)},\kk\right).
\]
長さ五、Hilbert 関数 $(3,1,1)$、最小 $\kk[[s,u,v]]$-Betti 数 $(3,8,6,1)$ である。
全 $B$-関係数は $M$ で十二、$Q$ で十二または十三なので、$A_0$ 上の非平坦性が従う。

$A_1=A_0/\mathfrak m_{A_0}^2$ に制限すると
\[
 M/\mathfrak a^2M\simeq Q/\mathfrak a^2Q
 \simeq A_1^4\oplus A_1/(\ell_1,\ell_2,\ell_3),
\]
\[
 \ell_1=t_8+9t_9-4t_{10}+5t_{14},\quad
 \ell_2=t_{11}-t_{14},\quad\ell_3=t_{13}-t_{14}.
\]
長さは六十七であり、平坦な階数五の場合の七十と異なる。
ここでの直和は $A_1$-加群としてのものに限る。
全次数でも自由階数四と循環torsionへの分解が成立する。次節でその根拠と残る障害を記す。

\section{Serre双対性が確定するcorank二の幾何}

有限可換代数の幾何的ファイバーについて、$\Omega^1/dA=0$ は
各局所Artin因子の埋込次元が高々一であることと同値である。
したがって $Q$ の台は、有限被覆のcorankが二以上となる場所の像と一致する。
二変数の場合には外微分から完全列
\[
 0\longrightarrow H^1(\Omega^\bullet_{T/B})
 \longrightarrow Q\longrightarrow\Omega^2_{T/B}\longrightarrow0
\]
を得る。全射性は二形式の係数を一変数で形式積分すれば分かる。

\subsection{全次数での消滅 $\psi=0$}
実際の五式から
\[
 \Lambda=\frac{B[[x,y]]}{(F_1,F_2,F_{1x},F_{1y},F_{2x})},
 \qquad\psi=[F_{2y}]\in\Lambda
\]
を作る。$\Lambda$ は十三次元正則局所環、$A_0$ 上有限自由階数四で、
閉ファイバーとパラメータの作用は
\[
 \Lambda/\mathfrak m_{A_0}\Lambda
 \simeq\kk[[x,y]]/(xy,y^2-2x^2),\qquad
 t_5=11x^2,\quad t_{15}=x^2,\quad t_{16}=0
\]
である。\textbf{実際の全形式族では $\psi=0$ である。} 従って
\[
 \Omega^2_{T/B}\simeq\Lambda,\qquad
 \dim\Supp_BQ=13,\qquad\rk_{A_0}Q=4.
\]

全次数の証明で再利用すべき構造は次の通りである。
普遍thin moduleを $\mathcal M$ とする。
自己Ext複体の最上微分は $D_2=[Q(z),-L(x)]$ であり、
関係式から $u^tD_2=0$ となる。$u$ はprimitiveで、五次単位小行列もあるので
$\Ext^3(\mathcal M,\mathcal M)\simeq T$ が、Serre固定性に先立って成立する。
相対Serre像の各頂点は閉点の計算とperfect complexの最小化により自由階数一である。
相対双対性によるHomの生成元は閉点で同型になり、全 $T$ 上で
$S_T(\mathcal M)\simeq\mathcal M[3]$ を得る。

単位ブロックを消した自己Ext複体は
\[
 T\xrightarrow{0}T^2\xrightarrow{J_F}T^2\xrightarrow{0}T
\]
である。鎖自己双対性から可逆行列 $U,V$ について $UJ_F=J_F^tV$ が得られる。
中心の二次部分
\[
 J_{F,2}=\begin{pmatrix}2xy&x^2\\y^2&2xy\end{pmatrix}
\]
の係数比較は、$U(0)=V(0)=c\left(\begin{smallmatrix}0&1\\1&0\end{smallmatrix}\right)$、
$c\ne0$ を強制する。
$\Lambda$ 上では $J_F=\operatorname{diag}(0,\psi)$ なので、
$(1,2)$ 成分から単元倍 $U_{12}\psi=0$ を得る。
この証明には、循環的な全potentialや鎖同型の対称性を仮定しない。

\subsection{全族の相対almost closed一形式表示}
同じ行列恒等式を転置して平均すると
\[
 P_0=(U+V^t)/2\in\operatorname{GL}_2(T),\qquad
 P_0J_F=J_F^tP_0^t
\]
を得る。$S=B[[x,y]]$ への可逆な持上げ $\widetilde P_0$ を選び、
$(a_1,a_2)^t=\widetilde P_0(F_1,F_2)^t$、
$\alpha=a_1dx+a_2dy$ と置けば
\[
 (a_1,a_2)=(F_1,F_2)=I,\qquad
 d_{S/B}\alpha\in I\Omega^2_{S/B}
\]
となる。従って全十六変数族は、全次数の相対almost closed一形式の零点である。
ここで使う微分は基底 $B$ を固定した相対微分である。
さらに、$P_0$ の原点値を固定した対称化条件は有限の非斉次線形方程式である。
代数的冪級数環からその完備化への忠実平坦性により解を降下できるため、
代数的かつ収束するalmost closed表示を選べる。
これは $d_{S/B}\alpha=0$ を意味せず、potentialの存在もまだ結論しない。
一般にalmost closedな一形式の零点が関数の臨界スキームにならない例は
Pandharipande--Thomas\cite{PT}にある。

\subsection{四枚被覆の全形式標準形}
長さ四の閉ファイバーの可換代数変形の接空間は四次元である。
実際の $\Lambda/A_0$ のKodaira--Spencer写像は全射であり、
$t_8,t_{11},t_{13},t_{14}$ の四方向の小行列式は $-8/5$ である。
切断を要求しない有限平坦代数変形の最小versal性から、ある形式座標で
\begin{align*}
 A_0&\simeq\kk[[u,v,w,z,s_1,\ldots,s_9]],\\
 \Lambda&\simeq A_0[[X,Y]]/(XY-u,\ Y^2-2X^2-v-wX-zY)
\end{align*}
となる。座標および分類写像の一意性は主張しない。

この四枚被覆の一般Galois群は $S_4$ である。
証明は四次方程式
\[
 2X^4+wX^3+vX^2+zuX-u^2=0
\]
の順序付き分解代数が、原点での完備化後にも整域で自由階数二十四であることによる。
有理関数体上のGalois群だけを計算して、完備化後の群を推測しない。

被覆の全空間内のramification divisorは
\[
 4X^2+2Y^2+wX-zY=0
\]
であり、線形座標で $\alpha\beta+\gamma\delta=0$ と書ける。
したがって三次元普通二重点と九次元形式円板の積である。
基底内のbranch divisorについても、次節のように正規化と自由因子構造が分かる。
$S_4$ の対象は、残留十一点の $S_{11}$、全点配置の $S_{21}$、
自己反対部分族の非推移的monodromyと区別する。
特に $Z_2$ の固定十一点に作用する $S_{11}$ も、残留十一点族のものとは別の被覆の群である。

\section{四枚被覆の分岐因子とcritical slice}

この節だけは九つの自明な方向を省き、
$R=\kk[[u,v,w,z]]$、
$\mathscr D=R[[X,Y]]/(XY-u,Y^2-2X^2-v-wX-zY)$ とする。
ここでの $u,v$ は最小versal標準形の座標であり、先の三方向ファイバーの
$t_5-11t_{15},t_{16}$ という座標と同一視しない。

\subsection{trace判別式・正規化・接コホモロジー}
次の行列は、相対接コホモロジーの八次表示から単位ブロックを消したものである。
\[
 C_\Delta=\begin{pmatrix}
 uw/4&-v/2&2u&(-3uz+vw)/4\\
 3uz&4u-wz&2v&(-uw-vz+w^2z)/2\\
 (8v-w^2)/4&-3z&w&(16u+5wz)/4\\
 -4u&3w/2&z&(4v-3w^2)/4
 \end{pmatrix},\qquad\delta=\det C_\Delta.
\]
$\delta$ は基底 $(1,X,Y,X^2)$ のtrace判別式であり、
重み $(2,2,1,1)$、次数八の被約既約多項式である。
branch divisor $\Delta=V(\delta)$ は自由因子で、$-C_\Delta$ の列が
対数的ベクトル場の自由基底を与える。具体的な恒等式は
\[
 (\partial_u\delta,\partial_v\delta,\partial_w\delta,\partial_z\delta)(-C_\Delta)
 =\delta\,(w/2,0,-8,3z/2)
\]
である。自由性の一般理論自体は既知であり、ここで再利用する成果は
この四枚被覆のtrace、different、接コホモロジーを同じ行列で表す明示式である\cite{PaperIII}。

$q=4X^2+2Y^2+wX-zY$、$p(X)=2X^4+wX^3+vX^2+zuX-u^2$ とすると
\[
 \operatorname{Nm}_{\mathscr D/R}(q)=-8\delta,\qquad
 \operatorname{disc}_X p=16u^2\delta.
\]
余分な $u^2$ は分岐そのものから来ない。
基底 $(1,X,X^2,X^3)$ を $(1,X,Y,X^2)$ で表示する行列式は $-u/2$ なので、
単一生成元 $X$ による表示の指数がこの因子を生じる。

有限射 $V(q)\to\Delta$ は正規化である。実際、$V(q)$ は正規な普通二重点であり、
$R_{(\delta)}$ 上で $q$ 倍写像の余核の長さが
$\operatorname{ord}_\delta\operatorname{Nm}(q)=1$ となる。
この長さが残余関数体の拡大次数なので、有限射は双有理である。
この論証は形式局所環の上で成立する。

さらに
\[
 J=\begin{pmatrix}Y&X\\-4X-w&2Y-z\end{pmatrix},\qquad
 N=T^1_{\mathscr D/R}=\coker J
\]
について
\[
 \End_R(N)\simeq\mathscr D/(q)
 \simeq\kk[[\alpha,\beta,\gamma,\eta]]/(\alpha\beta+\gamma\eta)
\]
である。$\operatorname{adj}(J)$ による $N\hookrightarrow(\mathscr D/(q))^2$ が
階数一のtorsion-free性を与え、共通の関数体内の有限自己準同型環を
整閉包と比較すれば等号を得る。この表示は、今後conductorやrank一加群を調べる入口になる。
$N$ は元の十一点族のpotential障害加群 $M$ とは異なる。

\subsection{四変数族全体は相対対称obstruction theoryを持たない}
四枚被覆の相対余接複体は
$[\mathscr D^2\xrightarrow{J^t}\mathscr D^2]$ である。
最小二項複体の自己双対性があれば、可逆行列 $U,V$ について
$UJ=J^tV$ が必要となる。
しかし独立な一次変数 $X,Y,w,z$ の係数比較は
$U(0)=V(0)=0$ を強制する。従って
\[
 \mathbb L_{\mathscr D/R}\not\simeq\mathbb L_{\mathscr D/R}^{\vee}[1].
\]

ここで使える一般的な補助事実は次である。
整な全空間を持つ、generically etaleな有限平坦局所完全交叉について、
相対対称perfect obstruction theoryがあれば、それは相対余接複体と擬同型になる。
対称性から一般点で複体は非輪状となり、負一次コホモロジーは
自由加群のtorsion部分加群なので零になるためである。
したがって $\mathscr D/R$ は相対対称perfect obstruction theoryを持たない。
また任意の滑らかな相対ambientのalmost closed零点表示は、
局所完全交叉の場合に余接複体を自己双対にするため、この被覆には存在しない。
九つの自明な変数を加えた実際の $\Lambda/A_0$ にも同じ結論が成り立つ。

\subsection{pencilの二つの例外と明示的な $D_4$ potential}
$(p,q)\ne(0,0)$ に対し、$(w,z)=(pt,qt)$ に制限した三変数族を考える。
\[
 \mathscr D_{p,q}
 =\frac{\kk[[u,v,t,X,Y]]}
 {(XY-u,\ Y^2-2X^2-v-ptX-qtY)}.
\]
次の条件は同値である。
\[
 p^2+2q^2\ne0;\quad
 \text{相対critical表示の存在};\quad
 \text{相対almost closed表示の存在};
\]
\[
 \text{相対対称perfect obstruction theoryの存在};\quad
 \mathbb L_{\mathscr D_{p,q}/\kk[[u,v,t]]}
 \simeq\mathbb L_{\mathscr D_{p,q}/\kk[[u,v,t]]}^{\vee}[1].
\]
存在する場合、potentialは多項式
\begin{align*}
 W_{p,q}={}&p\left(XY^2-\frac23X^3\right)
 +q\left(\frac13Y^3-2X^2Y\right)\\
 &-\frac t2(pX+qY)^2
 +(4qu-pv)X-(2pu+qv)Y
\end{align*}
と取れる。実際、
\[
 \binom{\partial_XW_{p,q}}{\partial_YW_{p,q}}
 =\begin{pmatrix}-4q&p\\2p&q\end{pmatrix}
 \binom{XY-u}{Y^2-2X^2-v-ptX-qtY},
\]
変換行列の行列式は $-2(p^2+2q^2)$ である。
中心の二変数三次形式の判別式は
$\frac83(p^2+2q^2)^2$ なので、二つの例外方向では
中心三次形式そのものが退化する。

良い方向では、Jacobi代数の非定数基底 $(X,Y,X^2)$ に対する
$u,v,t$ の接写像の行列式は $(p^2+2q^2)^2$ である。
したがってこの三変数族は、相対座標変更と基底だけの関数の加算を同一視した
$D_4$ 関数変形の最小versal族となる。
一方、長さ四の\emph{代数}の最小versal変形には四変数が必要である。
この差を単なる次元比較にとどめず、非自己双対性と上の明示的なslice分類で記述できる。

中心potentialを一つ固定してから許される接方向だけを数えると、
このpencilを正しく捉えられない。
中心のpotential自体が二次元の三次形式のpencilをなし、
slice方向 $[p:q]$ がその中心potentialの方向を選ぶからである。
また、$\Lambda/A_0$ の非critical性は、元の有限射 $T/B$ の
almost closed表示や残るpotential障害の値を決定しない。

\subsection{同じ非critical被覆を生む長さ十一のcritical族}
corank二の被覆の非critical性から元の点族の非critical性を導くことはできない。
この推論を排除する、正則全空間を持つ明示的な反例を保持する。
\[
 \widetilde B=\mathbb C[[a,b,c,u,v,w,z]],\qquad R=\mathbb C[[u,v,w,z]]
\]
上で
\begin{align*}
 \mathcal W={}&\frac12X^2Y^2+\frac15(X^5+Y^5)
 -\frac w6X^3+\frac z{12}Y^3-2uXY\\
 &+\frac c2X^2+\frac c4Y^2+\frac v4Y^2+aX+bY
\end{align*}
を取ると、
\[
 \widetilde T=\widetilde B[[X,Y]]/(\mathcal W_X,\mathcal W_Y)
 \simeq\mathbb C[[X,Y,c,u,v,w,z]]
\]
は $\widetilde B$ 上有限自由階数十一であり、全次数でcritical族である。
中心のJacobianイデアルは
$(X^4+XY^2,X^2Y+Y^4)$ で、局所的に $XY^3,Y^6$ を加えた
Gr\"obner基底を持ち、標準単項式の個数は十一である。
中心potentialの像も $-Y^5/10\ne0$ である。

そのcorank二の環は、$a,b,c$ を消去すると
\[
 \widetilde\Lambda=
 \frac{R[[X,Y]]}{(XY-u,\ Y^2-2X^2-v-wX-zY+4X^3-8Y^3)}
 \simeq\mathbb C[[X,Y,w,z]]
\]
となる。閉ファイバーでは $1-2X,1-8Y$ が単元であることを使い、
\[
 (XY,Y^2-2X^2+4X^3-8Y^3)=(XY,Y^2-2X^2)
\]
を得る。従って $R$ 上有限自由階数四である。
$u,v,w,z$ のKodaira--Spencer写像は同型であり、
Jacobianの独立な四列とパラメータの四列を合わせた小行列式は八に取れる。
最小versal性により、基底の形式座標変更も許せば
$\widetilde\Lambda/R$ は非criticalな四枚被覆 $\mathscr D/R$ と形式同型である。

全環に同じ九個の自明な形式変数を加えると、基底次元は十一点族で十六、
四枚被覆で十三となり、後者は実際の $\Lambda/A_0$ と同じ形式型を持つ。
これは元の $T/B$ とこのcritical族を同一視する主張ではない。
元の $\beta$ を判定するには、corank二の被覆への縮約で失われる情報が必要となる\cite{PaperIII}。
微分、中心の長さ、Kodaira--Spencer写像の独立検算入口は
\nolinkurl{critical_source_noncritical_corank_v25.py} である。

\section{残る相対potential障害は循環torsionの理想比較}

$M\twoheadrightarrow Q\twoheadrightarrow\Lambda\simeq A_0^4$ を $A_0$-線形に分裂させる。
既知の最小生成数五と非平坦性から $M,Q$ はともに一般階数四であり、
全torsion部分は循環加群である。真の非零イデアル $J_M\subseteq J_Q\subsetneq A_0$ に対し
\[
 M\simeq A_0^4\oplus A_0/J_M,\qquad
 Q\simeq A_0^4\oplus A_0/J_Q,\qquad
 H^1(\Omega^\bullet_{T/B})\simeq A_0/J_Q
\]
となる。これは $A_0$-加群の分解であり、$B$-加群分解とは限らない。

$\beta$ は $M$ のtorsion部分に属し、その生成元を選んで $b_0\bmod J_M$ と表すと
\[
 J_Q=J_M+(b_0),\qquad b_0\in\mathfrak m_{A_0}^2,
\]
\[
 \text{正規化した全相対potentialの存在}
 \ \Longleftrightarrow\ \beta=0
 \ \Longleftrightarrow\ b_0\in J_M
 \ \Longleftrightarrow\ J_M=J_Q.
\]
循環torsionは $B$ に安定で、$B$ の作用は $A_0/J_M$ の乗法を通るため、
$B\beta=A_0\beta$ である。イデアルは内在的に
\[
 J_M=\operatorname{Fitt}^{A_0}_4(M),\qquad
 J_Q=\operatorname{Fitt}^{A_0}_4(Q)
\]
と書ける。$b_0$ は生成元の変更で単元倍されるが、イデアルの一致は変わらない。

\subsection{三つの一次関係の後に真の高次関係が残る}
全 $B$ 上の最小関係数は $\beta_1^B(M)=12$ で、
$Q$ では $\beta=0$ なら十二、$\beta\ne0$ なら十三である。
四次までの持上げ可能中心が $\kk[W_9]$ だけであるため、
$C$ の全syzygyの中心値が零となることが、この最小性を保証する。
三方向ファイバーの第一Betti数は八なので、Koszul二重複体から
\[
 \beta_1^B(N)\le8+\beta_1^{A_0}(N),\qquad N=M,Q
\]
が従う。循環分解と合わせて
\[
 \mu(J_M)\ge4,\qquad\mu(J_Q)\ge4,
 \qquad\beta\ne0\Longrightarrow\mu(J_Q)\ge5
\]
を得る。

両イデアルの一次部分は独立な $\ell_1,\ell_2,\ell_3$ である。
$J_M$ 内の三つの持上げ $g_1,g_2,g_3$ を選ぶと、
\[
 A'=A_0/(g_1,g_2,g_3)\simeq\kk[[r_1,\ldots,r_{10}]],\qquad
 J'_M=J_M/(g_1,g_2,g_3)
\]
であり、\textbf{$0\ne J'_M\subset\mathfrak m_{A'}^2$} が成立する。
従って $J_M$ は三元生成ではない。
また $A'$ は整域なので
\[
 \operatorname{ht}J_M,\operatorname{ht}J_Q\ge4,
 \qquad \dim A_0/J_M,\dim A_0/J_Q\le9.
\]
これにより
\[
 \pd_BM\ge7,\qquad\pd_BQ\ge7,\qquad\pd_B\mathscr R\ge5
\]
となる。最後の不等式は
$0\to\mathscr R\to B^{33}\to B^{22}\to Q\to0$ の正確な二次元シフトによる。

\subsection{判定を前進させるために必要なもの}
\textbf{$J_M=J_Q$ は未判定である。}
高次関係の最初の次数、$J'_M$ の最小生成系、$b_0$ の像を同じ座標で求めることが
次の具体的な目標となる。単に三つの\emph{線形項}を零にしたsliceは、
非線形持上げ $g_1,g_2,g_3$ を消去した商と同じではない。
この違いを保った上で、非所属を検出する厳密な有限商、
または所属を全次数で強制する循環構造を選ぶ。

\subsection{小商の有限化と二次関係を取り出す順序}
$\mathfrak n=\mathfrak m_B$ とする。三方向ファイバーの長さ五の加群は
縦方向の極大イデアルの三乗で消えるため
\[
 \mathfrak n^3M\subset\mathfrak aM,\qquad
 \mathfrak n^{3k}M\subset\mathfrak a^kM\quad(k\ge1)
\]
であり、商加群 $Q$ にも同じ包含が成り立つ。特に
\[
 M/\mathfrak a^3M
 =M/(\mathfrak a^3+\mathfrak n^9)M
\]
なので、$A_0$ の第二近傍は、基底方向の次数二以下・全パラメータ次数八以下の
有限なSchur計算で厳密に求められる。
三つの既知の一次関係方向と二つの残余方向だけを基底に残し、
縦三方向を保つ場合、必要な係数単項式は
\[
 \binom{11}{3}+5\binom{10}{3}+15\binom{9}{3}=2025
\]
個である。これは必要な有限化の上界であり、計算済みの証明書を意味しない。
この全次数八の第二近傍証明書はまだ完成しておらず、
より低い全次数で切った商の消滅を $M/\mathfrak a^3M$ の消滅に読み替えてはならない。

三つの一次関係を消去した後の二次関係には、次の順序が必要である。
自由な四生成元を分離して関係を同じ座標へ運んだ後、まず一次関係写像の核を取る。
その核から生じる二次項を抽出し、$\ell_1=\ell_2=\ell_3=0$ の共通核上で比較する。
全ての関係の二次項を先にこの線形部分空間へ制限すると、
既知の一次関係の非線形持上げを新しい高次関係と誤認し得る。
全次数での消去には引き続き $g_1,g_2,g_3$ 自身を使う必要がある。

自己Ext複体の鎖自己双対性、almost closed表示、四枚被覆の正規化は、
いずれもこの一元所属を自動的には決定しない。
四枚被覆の非critical性を元の十一点族のpotential非存在へ移すこともできない。
新たな計算は、$b_0\bmod J_M$ を変化させる残余方向を特定できるときに行う。

\section{非可換側で有限 jet から結論しないための比較族}

$\rho^2-\rho+1=0$ とし、$F(s)=s+\sum_{j\ge2}f_js^j$ を取る。
原稿IVは
\[
 A_F=\kk\langle\!\langle X,Y,s\rangle\!\rangle/
 (Xs-\rho sX,\ Ys-\rho^{-1}sY,\ YX-\rho XY-F(s))
\]
を扱う。全ての $F$ で連続 $3$-CY であり、接錐は次の $B_0$ である。
\[
 B_0=\kk\langle X,Y\rangle/
 (X^2Y-XYX+YX^2,\ XY^2-YXY+Y^2X)
\]
\[
 \mathcal R_F(s)=\sum_{j\equiv5\ (6)}f_js^j
\]
と置くと
\[
 A_F\text{ が PI}\quad\Longleftrightarrow\quad\mathcal R_F=0.
\]
$\mathcal R_F\ne0$ のときは全中心が $\kk[[s^6]]$ になる。
$H=F^{-1}$、$K'=H$ とすれば二生成元 potential も
\[
 \Phi_F=-\rho^{-1}K(YX-\rho XY)
\]
と明示できる。

$F=s+c s^{6n-1}$、$c\ne0$ を取ると、中心が小さい非 PI の連続 $3$-CY 代数が得られる。
二生成元関係の最初の補正は次数 $12n-1$ なので、$n$ を増やせば
$B_0$ の完備化と任意に高い jet を一致させられる。
この一致は $B_0$ との比較であり、固定した別の代数 $\mathcal A$ の任意の jet を保存する主張ではない。

\subsection{固定した残留代数に関する保持事項}
固定代数 $\mathcal A$ については、三次 down-up 接錐、厳密な完備 PBW 表示、
連続 twisted $3$-CY 分解、最小 dg model が利用できる。
同次中心の候補 $q,U,D$ の任意の二元の全中心持上げは PI 性と同値であり、
その成立から全連続 CY 性と非可換 potential が従う。
この逆、すなわち固定 $\mathcal A$ の CY 性から PI 性が従うかは未決定である。
比較族によって、同じ接錐を持つ全代数に対する逆向きの一般定理は成立しないと分かった。

保存してある中心計算の有効範囲は次の通りである。
offset 十二・十四・十六の先頭 Poisson 障害は消え、奇数段階も合わせて共通 offset 十八まで進んでいる。
保存した有限 Hochschild 基本 cycle が保証する閉性は offset 十八未満に限られる。
次の段階で閉性を用いるには、内的単元の精度を延長して再検証する必要がある。
これを全形式 CY 性や全中心持上げと呼んではならない。

\section{公表稿から独立に保持すべき確定結果}

次の結果は継続研究で使う可能性が高く、計算入力と証明を保持する。
公表稿の主要な論旨から独立して再利用できるため、証明と入力の所在をまとめる。
\begin{enumerate}
\item 四パラメータ正則領域の全 miniversal hull は、滑らかな型と
\[
 \kk[[u_1,\ldots,u_{1+r},s,x_1,\ldots,x_{15}]]/(sx_1,\ldots,sx_{15})
\]
の型に分かれる。関係スキームには少なくとも次元五十三の $X_0$ と
次元六十七の $Y_E,Y_C$ がある。この局所分類を粗モジュライや全関係成分の分類と混同しない。
\item $Y_E,Y_C$ の一般点の Hochschild 次元はともに $(1,0,16,21)$。
この一致だけから二成分の導来非同値性は分からない。
\item 残留十一点の可換中心は
$x^2y^2/2-(x^5+y^5)/5$ の Jacobian 環と形式同型。
正規化 critical class の曖昧さと Jacobian 環の元の曖昧さは異なる。
標識付き class の延長可能空間の次元はパラメータ次数零から五で
$13,5,5,3,1,1$ となる。
\item 点対象の直交する球面 twist が忠実な $\mathbb Z^{21}$ を生成し、
$S_{21}$-monodromy とその降下が計算されている。
これは一般の有限代数の $\Out=1$ と両立する。
\item 分岐加群の正規化・conductor・判別式格子の公式。
特に相対 closed function の原始元の判別式には被覆判別式の三乗が現れる。
これらは $\beta=0$ の証明ではなく、存在した場合の臨界値の幾何と整環の構造に使う。
\end{enumerate}

原証明の所在は順に、計算付録の
\nolinkurl{global_full_hull_classification_section_ja.tex},
\nolinkurl{cusp_normal_form_section_ja.tex},
\nolinkurl{relative_critical_filtration_section_ja.tex},
\nolinkurl{relative_critical_quartic_section_ja.tex},
\nolinkurl{relative_critical_quintic_section_ja.tex},
\nolinkurl{point_orthogonal_spheres_section_ja.tex},
\nolinkurl{twist_descent_section_ja.tex},
\nolinkurl{relative_ramification_section_ja.tex},
\nolinkurl{ramification_reflexive_section_ja.tex}である。
計算付録内の版番号や未決という記載は、そのファイルの記録時点を表す。
継続判断には本ノートと入口文書を用いる。

\section{再現と独立監査の境界}

検証を再利用するときは、有限計算が確認する恒等式と全次数の構造定理を分ける。
\begin{itemize}
\item Serre自己双対性では単位消去と中心の二次係数比較、
四枚被覆ではKodaira--Spencer写像と四次乗法表を有理数上で検査する。
\item 分岐因子ではtrace判別式、norm、Saito行列、単一生成元の指数を多項式恒等式として検査する。
独立実装は保存済み乗法行列を使わずGroebner正規形から再構成する。
\item 非critical性は八未知数の一次線形系、slice分類は記号的な $p,q$ の恒等式で検査する。
有限個の方向への数値代入でpencil全体を代用しない。
critical族内の非critical被覆の例は、中心イデアルの長さ、Hessianの消去、
Kodaira--Spencer写像を有理数上で確認し、形式型の一致には最小versal性を用いる。
\item torsionの高さ下界は、五次元ファイバーのKoszul微分と三つの一次関係を独立に照合し、
その上でスペクトル系列と正則パラメータの議論を用いる。
\item 自己反対三族の被約性は、各々の同じ固定成分上のHensel解の存在と、
有限評価代数の判別式の非零性を組み合わせる。
有限精度で残差が小さいことだけでは、全方程式の解の存在を意味しない。
四次・六次の滑らかさは、同じHensel点での偏微分のMacaulay小行列式の非消滅で確認する。
\item 重み付き五直線の標準形は全incidenceを有理数上で検査し、
逆構成は可換な自己随伴行列の同時対角化による。
全monodromyは既約な臨界点対応、単独foldの厳密な零スキーム、
直線と二次曲線の交点被覆を組み合わせた幾何学的証明である。
\item Prym線束と非固定点除数の同定は、対称行列表示の双対性と部分束の行列式による。
滑らかさの証明書だけから、この除数同定や二重交差を推測しない。
\item Gram微分の階数は単元Jacobianを用いた接空間制限で検証する。
楕円線織面の同定は残余符号商と外積二乗の逆構成による。
六次の復元は階数七のGram分解と三次元核を用い、有限個の例から一意性を推測しない。
\item 平面直線束の一意性はRiemann--Rochと点の分離、被約性は乗法写像の全射性による。
$H^3=K_C$ の根を数えるだけでは双有理性を証明できない。
\item Hermite--Pfaffian恒等式は相補小行列式で証明し、全座標の記号計算で再現する。
二重被覆とtrace零線束は幾何学的に同定し、二次モジュラー方程式の検査と区別する。
\item Brauer障害はPrym線束の降下cocycleから求める。非分裂二次曲線を引き戻す支配族が非消滅を示す。
補助曲線の指数三には支配的net族と離散付値を使い、一般ファイバーの指数六には零cycleの積を使う。
最小次数六の閉点の実在は、指数の定義だけに任せず、次数三と二の直線切断から構成する。
\item 相対有理性では群の全五十四要素と表現の分解を厳密に照合し、不変式環の全次数の生成と降下を証明する。
Pl\"ucker縮約核とLie代数は独立な線形実装でも階数を確認する。
九余因子式は形式係数の恒等式、SNC境界は全局所式の消去で確認し、有限例だけから一般定理を推測しない。
\item 節点以外の特異点の排除は、Jacobian商の一つの次数の次元から長さを抑える補題を使う。
固定三点の全対称射影は連結群軌道の閉路から導き、算術Galois群を代用しない。
\item 逆Gramの四十五次小行列式を厳密に再計算し、行列式表示の古典的対応から二次元モジュライの支配性を導く。
第三被覆の指標表と判別式は記号的に照合し、非単有理性は非定数な楕円写像の幾何学的証明による。
\item K3境界格子は整数基底と貼合せを検算し、全根は不等式による証明と独立な有限列挙を照合する。
七次被覆では完全なファイバーと例外曲線の像を検査し、単純foldの一般点への持続を有限平坦族で証明する。
\item 境界periodは $L/\mathbb Zh$ の指標格子のSmith標準形を照合する。
次数九と $L_h$ へ制限した写像の次数を区別し、点配置から全 $T^1$ への持上げは束の完全列で証明する。
\item Kummerのtrace式は六変数の全係数、三次式 $I_m$ は十変数の全係数で検査する。
Moore標準形の三倍写像と原点同定を用い、有限同種写像の不定性を残さずJacobian座標を得る。
\item 分岐源の平滑化は三十節点を全て検査し、例外上六点では臨界点移動を含む二次係数を使う。
標数零への持上げと相対特異点スキームの固有性は、有限体の数値に加えて必要な幾何学的証明である。
\end{itemize}
各検証の入力、実行範囲、結果は
\nolinkurl{verification/verification_v30.json} および対応する証明書にまとめる。
v30では四十三の主検証と三つの補助検証を実行し、code・入力・証明書の
全snapshotについて実行前後の同一性を検査した。今回その既存記録と入力は保持する。結果と正確なファイル数は上記報告に記録する。
v31の新規検証は \nolinkurl{verification/verification_v31.json} にまとめ、
通常実行で入力・証明書が変わらないことも検査する。
v32の六ゲートの記録は \nolinkurl{verification/verification_v32.json} に保持する。
v33の検証記録は保持する。その統合検証器 \nolinkurl{code/verify_v33.py} は一般重みの混合消去、
微分Ore族、$c=0$ の正規形、接線型成分、可解Lie作用の星積と、
再利用する基底切片を検査する。
全出力JSONを保存済み証明書と比較し、実行前後のコード・入力・証明書のhashを確認する。
正確な実行範囲とファイル数は \nolinkurl{verification/verification_v33.json} および整合性報告に記録する。
旧版全体のゲートを再実行したという意味ではなく、旧版の入力と記録を保持する。
v34では \nolinkurl{code/verify_v34.py} が、接線系列の全重み恒等式、
重み三の全Schur式、Nakayama公式の規約、不変環と中心を検査する。
実行範囲と完全JSON比較は \nolinkurl{verification/verification_v34.json}、
コード・証明書の不変性は同じ版のintegrity報告に記録する。
全整数の生成性と成分性には、本文の構造的証明が必要である。
v35の通常ゲート \nolinkurl{code/verify_v35.py} は、固定量子化の導分、反標準環、
中心特殊化、四項Hochschild複体の四つを検査する。
全JSON比較と入力ハッシュ不変性を \nolinkurl{verification/verification_v35.json} に記録する。
中心化環と形式弧移送の独立数学監査も同梱する。
この検証は、以下の構造的証明や全重みの証明の代用とはしない。
v36の新規ゲートは、Gerstenhaber作用、偶数端点の全混合写像、primitive境界を検査する。
作用の全冪階数には、商行列と全コサイクル基底の二つの厳密な計算を保存する。
通常実行は既存証明書との全JSON比較を行い、入力と証明書を変更しない。
旧版のゲートは、それぞれの保存済み記録と入力を保持する。

v37のゲートは、引戻し型、端点の明示Ore族、VKの二次障害と独立な同時可解性、
原点の冪零指数、および正規スライスの恒等式を検査する。
全ての通常実行は凍結証明書との完全JSON比較と入力不変性を確認する。

v38の五ゲートと監査記録も保持する。
v39の \nolinkurl{code/verify_v39.py} は、純粋部分へのフィルトレーション還元、
五変数と平方零変数の有理式族、特異部分の降下模型、混合項を検出する中心境界を検査する。
有理式族は、分母を払った全 $2556$ 多項式方程式を元の積から直接再構成して検証する。
有限jetの外挿は用いない。境界と降下の全位数の主張は本文の一般証明による。
第五のゲートは、標準補完接線と純粋曲線の接線の正確な一致を、独立な二つの証明書から確認する。
実行前後のコード・入力・証明書の一致と完全JSON比較を、同版の検証報告に記録する。
独立監査では商複体から全形式還元への帰納法、非零平方関係からの全環判定、
巡回降下の逆行列表現、Brauer周期の付値証明を照合する。

全次数の結論の根拠も保持する。
$\psi=0$ には相対Serre双対性と全環上の単元恒等式、
四枚被覆の標準形には最小versal性と接写像の全射性、
$S_4$ には順序付き分解代数の完備化後の整域性を使う。
分岐因子の正規化にはnormの付値と正規な普通二重点、
三つの固定成分の開条件の非空性には代数的Hensel解と有限な付値証明を使う。
$\beta=0$ と固定 $\mathcal A$ の全PI性・全CY性は未判定である。


v40の \nolinkurl{code/verify_v40.py} は、全重みのフィルトレーション還元、
導分付き障害の推移、全中心と最深ファイバー、一般境界の形式降下・留数の四ゲートを実行する。
全整数の結論は本文の一般証明による。
全ての通常実行で完全JSON比較とコード・入力の不変性を確認する。
純粋部分の $m=3,4,5$ の表は任意実行の厳密プローブであり、一般定理と区別して保持する。

\section{次回の研究方向を選ぶ基準}

方向の優先順位は、未解決である期間や既に使った計算量では決めない。
得られる定理の範囲、一般化できる仕組み、厳密な判定への距離を比較し、
新しい構造が見つかれば順序を変更する。
現時点では、短期の具体性と成果の広がりについて次の判断が妥当である。

\subsection{全重みの純粋部分と高次の接合障害}

全重みの正の $w$ 重みの消滅と、純粋部分の一意な持上げは証明した。
同定を固定した相対変形の次元は $(m^2+9m+28)/2$ であり、同じ点の行基本変形を増やすことは次の主目的としない。
純粋部分の $HH^1$ は全ての重みで一次元と分かり、その生成元が四項へ持ち上がるため接続写像は零となった。
従って $R_{B_{j,r}}\simeq R_{N_m}[[k_m]]$ は全ての共鳴で成立する。
次の具体的な問題は、疎な四対象代数 $N_m$ の $HH^2$ と全変形環を決めることである。
$m=3,4,5$ の厳密計算は $HH^2=3$ を示すが、これだけから全重みの滑らかさを結論しない。
非可換射影平面との比較が有効かを検討する際も、
通常のtilting列ではないため、逆向きの高次Extを落とした比較を無条件に使わない。

標準六項の結合性障害線と四項の対称性障害との自然な比較は、導分付き障害複体で解決した。
一般重みでは、小拡大ごとの増強された障害を追い、
純粋部分の対称性曲線の切詰め次数と外側の $\gamma^{j+1}$ を比較する。
これは全ての $HH^3$ の間に未構成の写像が存在すると仮定する方法を避けられる。
別の例外列による半直交埋込みは、この標準列の補完問題とは別に残る。

\subsection{混合項の射影的な区別と捻れた解消}

一変数混合項 $H=f(Q,u)+xG$ の全中心と、$f=0$ の全閉ファイバーが決まった。
$f\in Q\mathbb C[Q,u]$ なら最深点のファイバーと完備局所指数はそのまま保たれる。
最深点では完成後にも指数 $\ell$ が残り、一般境界の閉点でのMcKay型NCCRとは異なる。
一般境界の巡回被覆は、横断 $A_{g-1}$ 解消の留数 $i\chi$ を消す最小係数被覆でもある。
この既知の計算を繰り返すより、一変数形に帰着しない混合項で中心をどのように貼るか、
また非零Brauer類を保持した代数層や捻れた解消をどこまで構成できるかを問う価値がある。

重み四・位数五では、$f=0$ と $f=Qu^3$ の全中心の特異部分の次元が二と一に分かれ、
無標識な全中心とアフィンdg導来Morita型を区別できた。
射影圏を区別するには、全中心の特異構造が反標準環またはSerre不変量へ残るかを調べる。
アフィン代数の区別から射影導来圏の非同値へは直接進まない。
正共鳴の非PI部分では、Du Val中心の部分解消上の代数層または相対圏が別の候補となる。
中心上有限なNCCRという語を非PI環へ適用せず、根基の各層と境界加群の貼合せから調べる。

\subsection{成分分類と別成分の量子化を進める条件}
重み三では $c=0$ と非零接線型の四共鳴が完了した。
残る非被約三次形式とその他の特異三次形式について、基底の形式切片と実際の移動核を同時に扱う。
十成分で全分類が尽きるかを問うが、混合核次元の増加だけから数を増やさない。
全重みでの成分数の下界と、一般点のArtin環の形をさらに増やす場合にも、
今回の重量・特異点評価の仕組みが別の三次層へ移るかを判断基準とする。

明示量子化の別の入口は有理型 $R(1,5),R(2,4)$ である。
それらの一般点に実際の座標環を与えられれば、補完可能なVK変形の幾何的範囲が広がる。
PBW性、AS正則性、通常のCY性と、既存の形式量子化の存在は区別する。
構造的な正規形が見えなければ、Prym・七次被覆・potentialの課題と再比較する。

\subsection{並行する候補：七次被覆の安定分岐と復元}

有限平坦性、分岐源の滑らかさ、正規化種数、微分零因子の次数は確定した。
残る特異点型を決定すれば、二十四次曲線の具体的な等特異族に研究を接続できる。
直接の入口は微分零因子の被約性と二点対応の横断性である。
通常尖点七十二個・節点百四十四個を目標とするが、三重像や高次接触を独立に排除する。
さらに分岐曲線または七次被覆から元のPrym対を復元できれば、
既存の七次元像に新たな内在的モジュライ解釈を与えられる。
標準七次被覆の $S_7$ だけではこの復元は従わない。

\subsection{独立に強い候補：三次pencilの商と有理性}
Kummer写像、Albanese、二次同種写像、$q=1$ は確定した。
次は一般三次四次元超曲面とその有限商の双有理型を分けて扱う。
三次memberは滑らかとは仮定しておらず、特異軌跡が有理パラメータ化の入口になる可能性がある。
ただし商の有理性と $k(B_m)$ 上の有理性が得られなければ、Albanese写像の双有理自明性には届かない。
四つの三次memberに沿う二重被覆モデルの共通基底と特異性も、群作用を保持した解消に使える。
元の共通平面被覆 $M_L$ の $q=0$ は有理性を意味しない。
$M_L\to\mathbb P^1_R$ の一般ファイバーの単有理性から有理連結性を強める方針は、
基底・一般ファイバーを扱う定理の仮定を確認した上で検討する。
九余因子成分のwheel型疎逆行列とSchur補による大域消去も、七次元Prym像の有理性に届くなら優先する。

\subsection{K3側：低位数Halphen層とperiodの大域化}
固定した境界型のCarlson periodと十八次元支配性、十分大きい全Halphen指数の実現は確定した。
明示Kummer写像を使い、逆Gramの良い開部分が除く有限個の $R$ 値と低位数torsionを照合することが次の具体的課題である。
全 $m\ge3$ の実現、最小の有効な閾値、および有限標識を忘れたperiod写像の次数は未決である。
中央ファイバーのCartier類の跳躍を、一般K3のNoether--Lefschetz条件へ移すには平滑化方向も制約する必要がある。
Kulikov模型の粗モジュライへの降下では、$T\oplus B$ の非自明なgerbe慣性を処理する。
全境界型の支配や、極限混合Hodge構造の全ての平滑化依存情報を得たとはしない。

\subsection{並行する課題：完全なmonodromyと標識の一意性}
$Z_1$ の固定三点への射影 $S_3$、$Z_2$ の全群、補助曲線の節点は確定している。
$Z_1$ の残る課題は、非固定九対の置換と内部交換、固定因子からの独立性である。
$Z_3$ では固定五点の推移性と単独互換、非固定八対の最大群を分けて検討する。
Prym線束が符号条件を強制する可能性を排除してから、wreath product全体を目標にする。
いずれも点配置の自己同型群、反対合の標識の一意性、Prym写像の内在化へ結び付く場合に優先する。

\subsection{条件付きで優先する課題：$J'_M$ と $b_0$ の比較}
元の十一点族の全相対potentialを決めることは、引き続き最も強い未決定理の候補である。
ただし高さ四の下界だけでは $J'_M$ の最初の次数も定まらない。
高次関係と $b_0$ を同じ商で計算するための縮約原理、
または非零性を検出する小さな厳密商が見えた時点で、この方向を第一候補へ上げる。
$\mathfrak n^9M\subset\mathfrak a^3M$ の有限化を使う場合も、
残余方向は $b_0$ と真の高次関係を区別できる構造に基づいて選ぶ。

採用する計算には、どの剰余classを判定し、その結果が所属・非所属にどう結び付くかを要求する。
一般的な低次元sliceでpotentialが存在しても、全基底での所属は従わない。
また、新たなjetで障害が消えただけでは全次数の存在証明にならない。
逆に一つの正確な商で $b_0$ の非零性を証明できれば、全族の非存在判定となる。
明示critical族でも同じ非criticalなcorank被覆を持てるため、
四枚被覆の非critical性だけに戻る方針では元の判定は進まない。

\subsection{構造が得られた場合に進める課題：固定非可換代数}
固定 $\mathcal A$ について、Nakayama自己同型の内的性、
potential複体、最小dg modelのcyclic構造を調べる。
PI性とCY性は別の目標とし、比較族で否定された一般的な逆含意を再び仮定しない。
全次数を制御する正規形や再帰が得られた場合には優先度を上げる。
単に既存中心jetの次の一段を埋めるだけなら、上の幾何と所属判定より優先しない。

\subsection{公表単位の判断}
I・II・III・IVは確定した範囲を保つ。
Vは行列式線束とK3境界period、VIは有限群の不変式・降下・Albanese、
VIIは六次Prym曲線の内在的復元と七次被覆の分岐、
VIIIは半直交成分の標準補完、IXは全重みの成分系列と厚み、Calabi--Yau量子化とDu Val中心、重み三の境界を担う。
Xは固定非零量子圏の分類、全AS正則実現、量子VKの全形式還元と対称性保持部分の二重構造を担う。
XIはDu Val中心、全正共鳴のスペクトル、根基の厚み、端点の中心境界と巡回表現族、Brauer障害を担う。
共通のGram記法は各稿で定義し、主定理と証明の流れを独立に読める構成とする。
今後の分割や統合も、単独で読める主定理と証明の依存関係を基準に判断する。

対称行列表示、楕円スクロール、Mukaiの一般K3構成、net写像自体は既知である。
本研究では関係族との同定と特殊閉包、具体的な像の方程式と標準除数、
安定化群を処理した相対有理性、一般関数体上の障害と指数、
明示period isogeny、Kummerの三次pencil、制約付き七次被覆の分岐を区別して位置付ける。
後続文献全体との新規性の最終照合は、公表前の別の作業として残す。

\begin{thebibliography}{99}
\interlinepenalty=10000
\bibitem{PaperI} 本bundle、公表原稿I、
\emph{Veronese--Kronecker 代数の標準部分代数と変形障害}, v22.
\nolinkurl{vk_canonical_deformations_ja_v22.pdf}.
\bibitem{PaperII} 本bundle、公表原稿II、
\emph{Fano72型関係代数の二十一点Torelliと反対代数}, v26.
\nolinkurl{fano72_point_configurations_ja_v26.pdf}.
\bibitem{PaperIII} 本bundle、公表原稿III、
\emph{有限点族の相対 Serre 双対性と corank 二の幾何}, v25.
\nolinkurl{relative_critical_corank2_ja_v25.pdf}.
\bibitem{PaperIV} 本bundle、公表原稿IV、
\emph{三次 down-up 接錐を持つ完備 Calabi--Yau 代数の共鳴と PI 性}, v22.
\nolinkurl{nc_cy_resonance_ja_v22.pdf}.
\bibitem{PaperV} 本bundle、公表原稿V、
\emph{四次Prymデータに付随する線織面と種数二十K3の境界周期}, v30.
\nolinkurl{self_opposite_prym_fibrations_ja_v30.pdf}.
\bibitem{ArtebaniDolgachevHesse}
M.~Artebani and I.~Dolgachev, \emph{The Hesse pencil of plane cubic curves},
L'Enseignement Math.\ \textbf{55} (2009), 235--273.
\url{https://doi.org/10.4171/LEM/55-3-3}.
\bibitem{LO} K.~Loginov and D.~Orlov,
\emph{Three-dimensional noncommutative Fano minifolds}, arXiv:2610.31250 (2026).
\url{https://arxiv.org/abs/2610.31250}.
\bibitem{PT} R.~Pandharipande and R.~P.~Thomas,
\emph{Almost closed 1-forms}, Glasgow Math. J. \textbf{56} (2014), 169--182.
\url{https://arxiv.org/abs/1204.3958}.
\bibitem{HartshorneNoether}
R.~Hartshorne, \emph{Generalized divisors on Gorenstein curves and a theorem of Noether},
J. Math. Kyoto Univ.\ \textbf{26} (1986), no.~3, 375--386.
\url{https://doi.org/10.1215/kjm/1250520873}.
\bibitem{CoppensKatoPlaneSystems}
M.~Coppens and T.~Kato, \emph{Non-trivial linear systems on smooth plane curves},
Math. Nachr.\ \textbf{166} (1994), 71--82.
\url{https://doi.org/10.1002/mana.19941660106}.
\url{https://arxiv.org/abs/alg-geom/9301003}.
\bibitem{DolgachevCAG}
I.~V.~Dolgachev, \emph{Classical Algebraic Geometry: A Modern View},
Cambridge University Press, 2012.
著者公開原稿、\S3.2.2 and Exercise~3.12.
\url{https://sites.lsa.umich.edu/idolga/wp-content/uploads/sites/1334/2024/08/CAG.21.pdf}.
\url{https://doi.org/10.1017/CBO9781139084437}.
\bibitem{StacksBrauerV27}
The Stacks Project Authors, \emph{The Stacks Project},
Tags 03R1 and 0A2J, Brauer groups and the Brauer group of a scheme.
\url{https://stacks.math.columbia.edu/tag/03R1}.
\url{https://stacks.math.columbia.edu/tag/0A2J}.
\bibitem{PaperVI} 本bundle、公表原稿VI、
\emph{種数三Prymモジュライの相対有理性と二次被覆}, v30.
\nolinkurl{prym_net_arithmetic_ja_v30.pdf}.
\bibitem{PaperVII} 本bundle、公表原稿VII、
\emph{六次Prym曲線の標準七次被覆と分岐曲線}, v30.
\nolinkurl{sextic_prym_septic_covers_ja_v30.pdf}.
\bibitem{MilneAbelianVarieties2008}
J.~S.~Milne, \emph{Abelian Varieties}, version 2.00, 2008,
I, \S3, Theorem~3.2 and Proposition~3.10.
\url{https://www.jmilne.org/math/CourseNotes/AV.pdf}.
\bibitem{DolgachevRationality2008}
I.~V.~Dolgachev, \emph{Rationality of $\mathcal R_2$ and $\mathcal R_3$},
Pure Appl. Math. Q. \textbf{4} (2008), no.~2, 501--508.
\url{https://doi.org/10.4310/PAMQ.2008.v4.n2.a8}.

\bibitem{PaperVIII} 本bundle、公表原稿VIII、
\emph{半直交成分の変形と重み付き射影空間への補完}, v32.
\nolinkurl{vk_ambient_completions_ja_v32.pdf}.
\bibitem{PaperIX} 本bundle、公表原稿IX、
\emph{重み付きPoisson変形の非被約成分とCalabi--Yau量子化}, v34.
\nolinkurl{weighted_poisson_components_ja_v34.pdf}.
\bibitem{Pridham2025}
J.~P.~Pridham, \emph{Quantisation of derived Poisson structures},
Geometry \& Topology \textbf{29} (2025), 3717--3771.
\url{https://arxiv.org/abs/1708.00496}.
\bibitem{PaperX} 本bundle、公表原稿X、
\emph{重み付き量子射影圏の変形環と対称性}, v40.
\nolinkurl{quantum_categories_ja_v40.pdf}.
\bibitem{PaperXI} 本bundle、公表原稿XI、
\emph{重み付きCalabi--Yau量子化の中心と表現}, v40.
\nolinkurl{quantum_primitive_spectra_ja_v40.pdf}.
\bibitem{Schlessinger1968}
M.~Schlessinger, \emph{Functors of Artin rings},
Trans. Amer. Math. Soc.\ \textbf{130} (1968), 208--222.
\url{https://doi.org/10.1090/S0002-9947-1968-0217093-3}.
\url{https://math.uchicago.edu/~amathew/schlessingerdef.pdf}.

\bibitem{KellerDerived2003}
B.~Keller, \emph{Derived invariance of higher structures on the Hochschild complex},
preprint, August 2003, revised May 16, 2018, Theorem~3.2 and its corollary.
\url{https://webusers.imj-prg.fr/~bernhard.keller/publ/dih.pdf}.

\bibitem{Manetti2005}
M.~Manetti, \emph{Deformation theory via differential graded Lie algebras},
Seminari di Geometria Algebrica 1998--1999, Scuola Normale Superiore (1999);
corrected version arXiv:math/0507284 (2005), Corollary~3.2.
\url{https://arxiv.org/abs/math/0507284}.

\bibitem{LinZhangSerre2025}
X.~Lin and S.~Zhang,
\emph{Serre algebra, matrix factorization and categorical Torelli theorem for hypersurfaces},
Math. Ann. \textbf{391} (2025), 163--177.
\url{https://doi.org/10.1007/s00208-024-02915-8}.
\url{https://arxiv.org/abs/2310.09927}.

\bibitem{MinamotoMori2011}
H.~Minamoto and I.~Mori,
\emph{The structure of AS-Gorenstein algebras},
Adv. Math. \textbf{226} (2011), no.~5, 4061--4095.
\url{https://doi.org/10.1016/j.aim.2010.11.004}.
Preprint title: \emph{Structures of AS-regular algebras}.
\url{https://www.math.kyoto-u.ac.jp/preprint/2010/06Minamoto.pdf}.

\bibitem{RickardMorita1989}
J.~Rickard, \emph{Morita theory for derived categories},
J. London Math. Soc. (2) \textbf{39} (1989), no.~3, 436--456.
\url{https://doi.org/10.1112/jlms/s2-39.3.436}.

\bibitem{RickardDerived1991}
J.~Rickard, \emph{Derived equivalences as derived functors},
J. London Math. Soc. (2) \textbf{43} (1991), no.~1, 37--48.
\url{https://doi.org/10.1112/jlms/s2-43.1.37}.

\bibitem{ORS2013}
J.~\"Oinert, J.~Richter and S.~Silvestrov,
\emph{Maximal commutative subrings and simplicity of Ore extensions},
Journal of Algebra and Its Applications \textbf{12} (2013), no.~4, 1250192, 16 pp.
\url{https://doi.org/10.1142/S0219498812501927}.
\url{https://arxiv.org/abs/1111.1292}.

\bibitem{JordanOre2014}
D.~A.~Jordan, \emph{Ore extensions and Poisson algebras},
Glasgow Mathematical Journal \textbf{56} (2014), no.~2, 355--368.
\url{https://doi.org/10.1017/S0017089513000293}.
\url{https://arxiv.org/abs/1212.4063}.

\bibitem{BellWangYee2019}
J.~P.~Bell, X.~Wang, and D.~Yee,
\emph{The Dixmier--Moeglin equivalence, Morita equivalence, and homeomorphism of spectra},
arXiv:1807.11813 (2018), proof of Theorem~1.1.
\url{https://arxiv.org/abs/1807.11813}.

\bibitem{LWWEndpoint37}
L.~Liu, S.~Wang, and Q.~Wu,
\emph{Twisted Calabi--Yau property of Ore extensions},
Journal of Noncommutative Geometry \textbf{8} (2014), 587--609,
doi:\href{https://doi.org/10.4171/JNCG/165}{10.4171/JNCG/165},
arXiv:\href{https://arxiv.org/abs/1205.0893}{1205.0893}.

\bibitem{ReyesRogalski2022}
M.~L.~Reyes and D.~Rogalski,
\emph{Graded twisted Calabi--Yau algebras are generalized Artin--Schelter regular},
Nagoya Mathematical Journal \textbf{245} (2022), 100--153,
doi:\href{https://doi.org/10.1017/nmj.2020.32}{10.1017/nmj.2020.32}.

\bibitem{StacksFitting37}
The Stacks Project Authors,
\emph{The Stacks Project}, Tags 07Z8 and 07ZA,
Fitting ideals and base change.
\url{https://stacks.math.columbia.edu/tag/07Z8}.
\url{https://stacks.math.columbia.edu/tag/07ZA}.

\bibitem{AlevDumas1994Bir38}
J.~Alev and F.~Dumas,
\emph{Sur le corps des fractions de certaines alg\`ebres quantiques},
Journal of Algebra \textbf{170} (1994), 229--265,
doi:\href{https://doi.org/10.1006/jabr.1994.1336}{10.1006/jabr.1994.1336}.

\bibitem{RichardSolotar2006Bir38}
L.~Richard and A.~Solotar,
\emph{Isomorphisms between quantum generalized Weyl algebras},
Journal of Algebra and Its Applications \textbf{5} (2006), 271--285,
arXiv:\href{https://arxiv.org/abs/math/0402413}{math/0402413}.

\bibitem{Bresar2010Bir38}
M.~Bre\v{s}ar,
\emph{An alternative approach to the structure theory of PI-rings},
arXiv:\href{https://arxiv.org/abs/1003.5275}{1003.5275}, 2010.

\bibitem{VanDenBerghMcKay38}
M.~Van den Bergh,
\emph{Non-commutative crepant resolutions},
in \emph{The Legacy of Niels Henrik Abel}, Springer, 2004, 749--770;
corrected version,
arXiv:\href{https://arxiv.org/abs/math/0211064}{math/0211064}, 2009.

\bibitem{MinacSchultzSwallowDescent39}
J.~Min\'{a}\v{c}, A.~Schultz and J.~Swallow,
\emph{Cyclic algebras, Schur indices, norms, and Galois modules},
Annales des sciences math\'{e}matiques du Qu\'{e}bec \textbf{35} (2011), no.~1, 123--136.
\href{https://arxiv.org/abs/math/0410536}{arXiv:math/0410536};
the splitting--norm criterion is recalled in \S1, equation~(2).

\bibitem{DGOBoundary39}
Yu.~A.~Drozd, B.~L.~Guzner and S.~A.~Ovsienko,
\emph{Weight modules over generalized Weyl algebras},
Journal of Algebra \textbf{184} (1996), 491--504.
\url{https://www.imath.kiev.ua/~drozd/Guzner.pdf}.
\bibitem{CGJBoundary39}
N.~Cooney, I.~Ganev and D.~Jordan,
\emph{Quantum Weyl algebras and reflection equation algebras at a root of unity},
Journal of Pure and Applied Algebra \textbf{224} (2020), no.~12, 106440.
\href{https://arxiv.org/abs/1907.11005}{arXiv:1907.11005};
doi:10.1016/j.jpaa.2020.106440.

\bibitem{DasMandal2026}
A.~Das and A.~Mandal,
\emph{Extension, deformation and categorification of AssDer Pairs},
Theory Appl. Categ. \textbf{45} (2026), no.~1, 1--32.
\url{https://tac.mta.ca/tac/volumes/45/1/45-01abs.html}.
\url{https://arxiv.org/abs/2002.11415}.

\bibitem{HarariBrauer40}
D.~Harari, \emph{The Brauer group of fields and schemes}, M2 lecture notes, 2023--2024.
Propositions~4.23 and~5.8.
\url{https://www.imo.universite-paris-saclay.fr/~david.harari/enseignement/m2brauer24/brauer24.pdf}.

\bibitem{GilleSzamuelyBrauer40}
P.~Gille and T.~Szamuely, \emph{Central Simple Algebras and Galois Cohomology},
Cambridge Studies in Advanced Mathematics 101, Cambridge University Press, 2006.
Propositions~2.5.2 and~4.7.3.
\url{https://www.math.ens.psl.eu/~benoist/refs/Gille-Szamuely.pdf}.

\end{thebibliography}
\end{document}
